% Mouvement des planètes et des satellites — Physique Tle D — Cours signé OnBuch+
% Programme officiel MINESEC. Compiler avec : tectonic P07-mouvement-planetes-satellites.tex
\newcommand{\DOCMATIERE}{Physique}
\newcommand{\DOCNIVEAU}{Tle D}
\newcommand{\DOCMODULE}{Module 2 : Mouvements et interactions — évolutions temporelles des systèmes mécaniques}
\newcommand{\DOCLECON}{P07}
\newcommand{\DOCTITRE}{Mouvement des planètes et des satellites}
% OnBuch+ — préambule « cours signés » ENRICHI (compilé avec tectonic / XeLaTeX)
% Système visuel « Pop » d'OnBuch+ : fond crème (jamais blanc), bordures encre
% épaisses, ombres « dures » décalées, titres Archivo Black en capitales.
% Version MATHÉMATIQUES : graphes (pgfplots/tikz), boîtes pédagogiques
% (définition, propriété, méthode, attention, le savais-tu, exercice résolu,
% exercice, TP, bilan), page de garde et signature OnBuch+ ancrée en bas.
%
% Macros attendues dans main.tex AVANT \input{preamble} :
%   \DOCMATIERE  \DOCNIVEAU  \DOCMODULE  \DOCLECON (numéro)  \DOCTITRE
\documentclass[11pt]{article}

\usepackage{fontspec}
\usepackage[french]{babel}
\usepackage[a4paper,margin=2cm,top=3.4cm,bottom=2.8cm,headheight=1.4cm,footskip=1.3cm]{geometry}
\usepackage{amsmath,amssymb,amsfonts}
\usepackage{mathtools} % \xmapsto, \coloneqq... souvent utilisés par les modèles
\usepackage{cancel} % simplifications barrées (bilans de forces)
% \abs{z} (module, valeur absolue) et \norm{u} (norme), utilisés par les modèles
\providecommand{\abs}{}\let\abs\relax\DeclarePairedDelimiter{\abs}{\lvert}{\rvert}
\providecommand{\norm}{}\let\norm\relax\DeclarePairedDelimiter{\norm}{\lVert}{\rVert}
\usepackage[table]{xcolor} % [table] : fournit aussi \rowcolors (alternance de lignes)
\usepackage{pagecolor}
\usepackage{tikz}
\usetikzlibrary{arrows.meta,positioning,calc,shapes.geometric,shapes.misc,shapes.arrows,decorations.pathmorphing,decorations.markings,patterns,shadows,mindmap,trees,fit,backgrounds,matrix,intersections}
\usepackage{pgfplots}
\usepackage[version=4]{mhchem} % équations nucléaires \ce{^{235}_{92}U}
\usepackage[european,siunitx]{circuitikz} % schémas électriques
\pgfplotsset{compat=1.18}
\usepgfplotslibrary{fillbetween} % aires entre courbes (calcul intégral)
\usepackage{enumitem}
\usepackage{fancyhdr}
% [most] charge toutes les bibliothèques tcolorbox (listings, minted, raster,
% documentation...) et gonfle inutilement la mémoire TeX sur un document long
% (~300 boîtes) : on ne charge que ce qui est réellement utilisé.
\usepackage[skins,breakable]{tcolorbox}
\usepackage{graphicx}
\usepackage{colortbl}
\usepackage{booktabs}
\usepackage{tabularx}
\usepackage{diagbox} % case d'en-tête diagonale des tableaux à double entrée
% Colonnes extensibles centrée / gauche / droite (C, L, R), souvent utilisées par les modèles
\newcolumntype{C}{>{\centering\arraybackslash}X}
\newcolumntype{L}{>{\raggedright\arraybackslash}X}
\newcolumntype{R}{>{\raggedleft\arraybackslash}X}
\usepackage{array}
\usepackage{multicol}
\usepackage{siunitx}
% parse-numbers=false : accepte aussi \SI{2,5 \times 10^{-3}}{...}
\sisetup{locale=FR,output-decimal-marker={,},group-minimum-digits=4,parse-numbers=false}
\DeclareSIUnit\division{div} % réglages d'oscilloscope (ms/div, V/div)
\usepackage{qrcode}
% \degree : commande fréquemment utilisée par les modèles pour un angle ou une
% température en dehors de \SI{}{} (non fournie par les paquets chargés).
\providecommand{\degree}{^{\circ}}
% \qed (fin de démonstration), utilisé par les modèles sans amsthm
\providecommand{\qed}{\hfill$\square$}
% \barre : double barre des valeurs interdites dans un tableau de variations (tkz-tab)
\providecommand{\barre}{\ensuremath{\|}}
% environnement proof (amsthm non chargé) utilisé par les modèles
\ifdefined\proof\else\newenvironment{proof}[1][Démonstration]{\par\noindent\textit{#1.}\ }{\qed\par}\fi
\usepackage{lastpage}
\usepackage{hyperref}
\hypersetup{hidelinks}

% ============================================================
% Palette « Pop » OnBuch+ — mêmes hex que la classe P (plus_ui.dart)
% ============================================================
\definecolor{poporange}{HTML}{F59321}
\definecolor{popdark}{HTML}{E07A0C}
\definecolor{popcream}{HTML}{FFF6E8}
\definecolor{popink}{HTML}{1C1714}
\definecolor{poppurple}{HTML}{7A5AE0}
\definecolor{popgreen}{HTML}{2E9E6A}
\definecolor{poppink}{HTML}{E0457B}
\definecolor{popblue}{HTML}{2F7DE1}
\definecolor{popgold}{HTML}{C98A12}
\definecolor{popmuted}{HTML}{8A7F76}
\definecolor{popline}{HTML}{EADFD2}
\definecolor{popgray}{HTML}{8A7F76}
\colorlet{popgrey}{popgray}
% Alias tolérants (le modèle nomme parfois les couleurs différemment)
\colorlet{popwhite}{white}
\colorlet{popblack}{popink}
\colorlet{popred}{poppink}
\colorlet{popyellow}{popgold}
% Teintes claires pour fonds de tableaux / zones
\colorlet{popblueL}{popblue!10!white}
\colorlet{poporangeL}{poporange!14!white}
\colorlet{popgreenL}{popgreen!12!white}
\colorlet{poppurpleL}{poppurple!12!white}
\colorlet{poppinkL}{poppink!12!white}
\colorlet{popgoldL}{popgold!14!white}

\pagecolor{popcream}

% ============================================================
% Typographie
% ============================================================
\defaultfontfeatures{Ligatures=TeX}
\setmainfont{PlusJakartaSans-Medium.ttf}[Path=../../../fonts/,BoldFont=PlusJakartaSans-Bold.ttf]
% Maths en sans-serif (Fira Math) avec chiffres et lettres droites de Plus
% Jakarta, pour que les formules se fondent dans le texte.
\usepackage[math-style=ISO,bold-style=ISO]{unicode-math}
\setmathfont{FiraMath-Regular.otf}[Path=../../../fonts/]
\setmathfont{PlusJakartaSans-Medium.ttf}[Path=../../../fonts/,range={up/num,up/{Latin,latin},\mathup}]
\usepackage{icomma} % virgule décimale sans espace en mode math
% Caractères mathématiques Unicode absents des polices de texte (Plus Jakarta,
% Archivo Black) : rendus par leur équivalent mathématique, en texte comme en
% formule — sinon ils s'affichent en carré vide (ex. « Arithmétique dans ℤ »).
\usepackage{newunicodechar}
\newunicodechar{ℝ}{\ensuremath{\mathbb{R}}}
\newunicodechar{ℤ}{\ensuremath{\mathbb{Z}}}
\newunicodechar{ℂ}{\ensuremath{\mathbb{C}}}
\newunicodechar{ℕ}{\ensuremath{\mathbb{N}}}
\newunicodechar{ℚ}{\ensuremath{\mathbb{Q}}}
\newunicodechar{𝔻}{\ensuremath{\mathbb{D}}}
\newunicodechar{α}{\ensuremath{\alpha}}
\newunicodechar{β}{\ensuremath{\beta}}
\newunicodechar{γ}{\ensuremath{\gamma}}
\newunicodechar{δ}{\ensuremath{\delta}}
\newunicodechar{ε}{\ensuremath{\varepsilon}}
\newunicodechar{θ}{\ensuremath{\theta}}
\newunicodechar{λ}{\ensuremath{\lambda}}
\newunicodechar{μ}{\ensuremath{\mu}}
\newunicodechar{σ}{\ensuremath{\sigma}}
\newunicodechar{τ}{\ensuremath{\tau}}
\newunicodechar{φ}{\ensuremath{\varphi}}
\newunicodechar{ψ}{\ensuremath{\psi}}
\newunicodechar{ω}{\ensuremath{\omega}}
\newunicodechar{Σ}{\ensuremath{\Sigma}}
% majuscules produites par \MakeUppercase dans les titres (α-aminés -> Α)
\newunicodechar{Α}{\ensuremath{\alpha}}
\newunicodechar{Β}{\ensuremath{\beta}}
\newunicodechar{Γ}{\ensuremath{\Gamma}}
\newunicodechar{Δ}{\ensuremath{\Delta}}
\newunicodechar{Ε}{\ensuremath{\varepsilon}}
\newunicodechar{Θ}{\ensuremath{\Theta}}
\newunicodechar{Λ}{\ensuremath{\Lambda}}
\newunicodechar{Μ}{\ensuremath{\mu}}
\newunicodechar{Τ}{\ensuremath{\tau}}
\newunicodechar{Φ}{\ensuremath{\Phi}}
\newunicodechar{Ψ}{\ensuremath{\Psi}}
\newunicodechar{Ω}{\ensuremath{\Omega}}
\newunicodechar{↦}{\ensuremath{\mapsto}}
\newunicodechar{∈}{\ensuremath{\in}}
\newunicodechar{∉}{\ensuremath{\notin}}
\newunicodechar{∧}{\ensuremath{\wedge}}
\newunicodechar{⇔}{\ensuremath{\Leftrightarrow}}
\newunicodechar{⟺}{\ensuremath{\Longleftrightarrow}}
\newunicodechar{⇒}{\ensuremath{\Rightarrow}}
\newunicodechar{⟹}{\ensuremath{\Longrightarrow}}
\newunicodechar{∘}{\ensuremath{\circ}}
\newunicodechar{∪}{\ensuremath{\cup}}
\newunicodechar{∩}{\ensuremath{\cap}}
\newunicodechar{≡}{\ensuremath{\equiv}}
\newunicodechar{⊂}{\ensuremath{\subset}}
\newunicodechar{⊥}{\ensuremath{\perp}}
\newunicodechar{∃}{\ensuremath{\exists}}
\newunicodechar{∀}{\ensuremath{\forall}}
\newunicodechar{∛}{\ensuremath{\sqrt[3]{\,}}}
\newunicodechar{∜}{\ensuremath{\sqrt[4]{\,}}}
\newunicodechar{⃗}{\ensuremath{{}^{\to}}}
\newunicodechar{ⁿ}{\ifmmode^{n}\else\textsuperscript{\ensuremath{n}}\fi}
\newunicodechar{ˣ}{\ifmmode^{x}\else\textsuperscript{\ensuremath{x}}\fi}
\newunicodechar{ᵇ}{\ifmmode^{b}\else\textsuperscript{\ensuremath{b}}\fi}
\newunicodechar{ʳ}{\ifmmode^{r}\else\textsuperscript{\ensuremath{r}}\fi}
\newunicodechar{ᵘ}{\ifmmode^{u}\else\textsuperscript{\ensuremath{u}}\fi}
\newunicodechar{ʸ}{\ifmmode^{y}\else\textsuperscript{\ensuremath{y}}\fi}
\newunicodechar{ᵃ}{\ifmmode^{a}\else\textsuperscript{\ensuremath{a}}\fi}
\newunicodechar{ᵉ}{\ifmmode^{e}\else\textsuperscript{\ensuremath{e}}\fi}
\newunicodechar{ᵏ}{\ifmmode^{k}\else\textsuperscript{\ensuremath{k}}\fi}
\newunicodechar{ᵖ}{\ifmmode^{p}\else\textsuperscript{\ensuremath{p}}\fi}
\newunicodechar{ᴺ}{\ifmmode^{N}\else\textsuperscript{\ensuremath{N}}\fi}
\newunicodechar{ᶜ}{\ifmmode^{c}\else\textsuperscript{\ensuremath{c}}\fi}
\newunicodechar{ʰ}{\ifmmode^{h}\else\textsuperscript{\ensuremath{h}}\fi}
\newunicodechar{ᵠ}{\ifmmode^{q}\else\textsuperscript{\ensuremath{q}}\fi}
\newunicodechar{ₙ}{\ifmmode_{n}\else\textsubscript{\ensuremath{n}}\fi}
\newunicodechar{ₐ}{\ifmmode_{a}\else\textsubscript{\ensuremath{a}}\fi}
\newunicodechar{ᵢ}{\ifmmode_{i}\else\textsubscript{\ensuremath{i}}\fi}
\newunicodechar{ᵣ}{\ifmmode_{r}\else\textsubscript{\ensuremath{r}}\fi}
\newunicodechar{ₖ}{\ifmmode_{k}\else\textsubscript{\ensuremath{k}}\fi}
\newunicodechar{ᵦ}{\ifmmode_{\beta}\else\textsubscript{\ensuremath{\beta}}\fi}
\newunicodechar{ₓ}{\ifmmode_{x}\else\textsubscript{\ensuremath{x}}\fi}
\newfontfamily\popdisplay{ArchivoBlack-Regular.ttf}[Path=../../../fonts/]
\newfontfamily\poplogo{SpaceGrotesk-Bold.ttf}[Path=../../../fonts/]
\newfontfamily\popsemi{PlusJakartaSans-SemiBold.ttf}[Path=../../../fonts/]
\newfontfamily\popxbold{PlusJakartaSans-ExtraBold.ttf}[Path=../../../fonts/]
\newfontfamily\popxboldls{PlusJakartaSans-ExtraBold.ttf}[Path=../../../fonts/,LetterSpace=3.0]

\setlist[enumerate]{leftmargin=1.7em,itemsep=5pt,topsep=4pt}
\setlist[itemize]{leftmargin=1.7em,itemsep=4pt,topsep=4pt,label={\color{poporange}\textbullet}}
\setlength{\parindent}{0pt}
\setlength{\parskip}{5pt}
\renewcommand{\arraystretch}{1.25}

% pgfplots : style Pop par défaut
\pgfplotsset{
  every axis/.append style={axis line style={popink,line width=1pt},tick style={popink},
    grid style={popline,line width=0.5pt},label style={font=\small},tick label style={font=\footnotesize}},
  popcurve/.style={poporange,line width=1.8pt,no markers,smooth},
}
% Même style disponible dans un \draw TikZ simple (hors pgfplots)
\tikzset{popcurve/.style={poporange,line width=1.8pt}}
\tikzset{popgrid/.style={popline,very thin}}

% ============================================================
% Pill / squiggle
% ============================================================
\newcommand{\poppill}[3]{%
  \tikz[baseline=-0.55ex]\node[draw=popink,line width=1.2pt,rounded corners=2.6pt,%
    fill=#2,inner xsep=6.5pt,inner ysep=3pt]{%
    \popxboldls\fontsize{7}{7}\selectfont\color{#3}\MakeUppercase{#1}};%
}
\newcommand{\popsquiggle}[1]{%
  \tikz[baseline=-0.4ex]{
    \draw[#1,line width=2.6pt,line cap=round]
      plot [smooth] coordinates {(0,0.09) (0.34,0.01) (0.68,0.21) (1.02,0.01) (1.36,0.10)};
  }%
}
% Mot-clé surligné (marqueur orange sous le texte)
\newcommand{\cle}[1]{{\popxbold\color{popdark}#1}}

% ============================================================
% En-tête / pied
% ============================================================
\pagestyle{fancy}
\fancyhf{}
\renewcommand{\headrulewidth}{0pt}
\renewcommand{\footrulewidth}{0pt}
\lhead{\raisebox{-0.15ex}{\poplogo\fontsize{13}{13}\selectfont\color{popink}OnBuch\color{poporange}+}}
\rhead{\poppill{\DOCMATIERE\ \textbullet\ \DOCNIVEAU\ \textbullet\ Leçon \DOCLECON}{white}{popink}}
\lfoot{\popxbold\fontsize{7.6}{7.6}\selectfont\color{popmuted}\MakeUppercase{Cours signé OnBuch+}\ \textbullet\ {\popsemi\DOCTITRE}}
\rfoot{\tikz[baseline=-0.6ex]\node[rounded corners=8.5pt,fill=popink,minimum height=17pt,minimum width=17pt,inner xsep=5pt,inner sep=0]{\popxbold\fontsize{8}{8}\selectfont\color{popcream}\thepage\,/\,\pageref*{LastPage}};}

% ============================================================
% Titres
% ============================================================
\newcommand{\chaptitre}[2]{%
  \begin{tcolorbox}[enhanced,colback=poporange,colframe=popink,boxrule=2.6pt,arc=11pt,
      drop shadow=popink,left=15pt,right=15pt,top=12pt,bottom=11pt,before skip=3pt,after skip=18pt]
    \poppill{#2}{popink}{popcream}\par\vspace{9pt}
    {\raggedright\hyphenpenalty=10000\exhyphenpenalty=10000\popdisplay\fontsize{22}{24}\selectfont\color{popcream}\MakeUppercase{#1}\par}
  \end{tcolorbox}
}
% Section numérotée : pastille orange avec le numéro + titre Archivo
\newcounter{popsec}
\newcommand{\coursec}[1]{%
  \stepcounter{popsec}\setcounter{popsub}{0}%
  \par\vspace{16pt}\noindent
  \begin{minipage}{\linewidth}
  \tikz[baseline=-0.7ex]\node[draw=popink,line width=1.6pt,fill=poporange,rounded corners=4pt,minimum size=22pt,inner sep=0]{\popdisplay\fontsize{12}{12}\selectfont\color{popcream}\thepopsec};%
  \hspace{8pt}{\popdisplay\fontsize{15}{17}\selectfont\color{popink}\MakeUppercase{#1}}\par
  \vspace{4pt}\noindent{\color{poporange}\rule{36pt}{3.6pt}}
  \end{minipage}\par\nopagebreak\vspace{8pt}
}
\newcounter{popsub}
\newcommand{\courssub}[1]{%
  \stepcounter{popsub}%
  \par\vspace{9pt}\noindent{\popxbold\fontsize{11.5}{13}\selectfont\color{popink}{\color{poporange}\thepopsec.\thepopsub}\hspace{6pt}#1}\par\nopagebreak\vspace{4pt}
}

% ============================================================
% Boîtes pédagogiques — même recette : fond blanc, bordure épaisse colorée,
% ombre dure encre, badge plein. Titre optionnel [#1].
% ============================================================
\tcbset{popbase/.style={breakable,enhanced,colback=white,boxrule=2.3pt,arc=9pt,
  shadow={3pt}{-3pt}{0pt}{popink},left=14pt,right=14pt,top=11pt,bottom=11pt,before skip=11pt,after skip=15pt,
  fontupper=\popsemi\fontsize{10.6}{14.5}\selectfont\color{popink},
  fontlower=\popsemi\fontsize{10.6}{14.5}\selectfont\color{popink}}}
% Coche dessinée (le glyphe ✓ n'existe pas dans les polices du document)
\newcommand{\popcheck}[1]{\tikz[baseline=-0.5ex]\draw[#1,line width=1.6pt,line cap=round,line join=round](-0.13,0.0)--(-0.04,-0.09)--(0.13,0.1);}
\providecommand{\coche}{\popcheck{popgreen}}
\newcommand{\popboxhead}[3]{\poppill{#1}{#2}{white}\ifx\relax#3\relax\else\hspace{6pt}{\popxbold\fontsize{10.6}{12}\selectfont\color{popink}#3}\fi\par\vspace{7pt}}

\newtcolorbox{aretenir}[1][]{popbase,colframe=popgreen,before upper={\popboxhead{À retenir}{popgreen}{#1}}}
\newtcolorbox{exemplebox}[1][]{popbase,colframe=poporange,before upper={\popboxhead{Exemple}{poporange}{#1}}}
\newtcolorbox{definition}[1][]{popbase,colframe=popblue,before upper={\popboxhead{Définition}{popblue}{#1}}}
\newtcolorbox{propriete}[1][]{popbase,colframe=poppurple,before upper={\popboxhead{Propriété}{poppurple}{#1}}}
\newtcolorbox{methode}[1][]{popbase,colframe=popgold,before upper={\popboxhead{Méthode}{popgold}{#1}}}
\newtcolorbox{attention}[1][]{popbase,colframe=poppink,before upper={\popboxhead{Attention}{poppink}{#1}}}
\newtcolorbox{experience}[1][]{popbase,colframe=popblue,colback=popblueL,before upper={\popboxhead{Expérience}{popblue}{#1}}}
% « Le savais-tu ? » : carte violette pleine (contexte, culture, Cameroun)
\newtcolorbox{savaistu}[1][]{popbase,colback=poppurple,colframe=popink,
  fontupper=\popsemi\fontsize{10.4}{14.2}\selectfont\color{white},
  before upper={\poppill{Le savais-tu\,?}{white}{poppurple}\ifx\relax#1\relax\else\hspace{6pt}{\popxbold\color{white}#1}\fi\par\vspace{7pt}}}
% Exercice résolu : énoncé en haut, \tcblower puis la solution sur fond crème
\newcounter{popexo}\newcounter{popexr}
\newtcolorbox{exoresolu}[1][]{popbase,colframe=poporange,colbacklower=popcream,
  segmentation style={popink,line width=1.2pt,dashed},
  before upper={\stepcounter{popexr}\popboxhead{Exercice résolu \thepopexr}{poporange}{#1}},
  before lower={\poppill{Solution}{popink}{popcream}\par\vspace{6pt}}}
% Exercice (non résolu en place) — la correction est dans la partie Corrigés
\newtcolorbox{exercice}[1][]{popbase,colframe=popink,
  before upper={\stepcounter{popexo}\popboxhead{Exercice \thepopexo}{popink}{#1}}}
% Corrigé
\newtcolorbox{corrige}[1][]{popbase,colframe=popgreen,colback=popgreenL,
  before upper={\popboxhead{Corrigé}{popgreen}{#1}}}
% Objectifs / prérequis / bilan
\newtcolorbox{objectifs}{popbase,colframe=popink,colback=poporangeL,
  before upper={\popboxhead{Objectifs de la leçon}{popink}{}}}
\newtcolorbox{prerequis}{popbase,colframe=popink,colback=popblueL,
  before upper={\popboxhead{Prérequis}{popblue}{}}}
\newtcolorbox{bilan}{popbase,colframe=popink,colback=white,boxrule=2.8pt,
  before upper={\popboxhead{Fiche bilan}{poporange}{L'essentiel à connaître pour l'examen}}}
% Figure encadrée avec légende
\newtcolorbox{popfigure}[1][]{enhanced,breakable=false,colback=white,colframe=popline,boxrule=1.4pt,arc=8pt,
  left=10pt,right=10pt,top=10pt,bottom=8pt,before skip=10pt,after skip=12pt,halign=center,
  lower separated=false,halign lower=center,
  fontlower=\popsemi\fontsize{9}{11.5}\selectfont\color{popmuted},
  #1}
\newcommand{\legende}[1]{\par\vspace{6pt}{\popsemi\fontsize{9}{11.5}\selectfont\color{popmuted}#1}}

% ============================================================
% Signature OnBuch+
% ============================================================
\newcommand{\poplogoinline}[1]{{\poplogo\fontsize{#1}{#1}\selectfont\color{popink}OnBuch\color{poporange}+}}
\newcommand{\signatureonbuch}{%
  \begin{tcolorbox}[enhanced,colback=popink,colframe=popink,boxrule=2.2pt,arc=10pt,
      drop shadow=poporange,left=14pt,right=14pt,top=10pt,bottom=10pt,before skip=0pt,after skip=0pt]
    \tikz[baseline=-0.7ex]\node[circle,fill=popgreen,draw=popcream,line width=1.3pt,minimum size=22pt,inner sep=0]{\popcheck{white}};%
    \hspace{9pt}\begin{minipage}[c]{0.62\linewidth}
      {\popxbold\fontsize{12}{13}\selectfont\color{popcream}Signé\ \textbullet\ {\poplogo OnBuch\color{poporange}+}}\par\vspace{2pt}
      {\raggedright\popsemi\fontsize{8}{10.5}\selectfont\color{popline}\MakeUppercase{Équipe pédagogique OnBuch+}\\\MakeUppercase{Conforme au programme officiel MINESEC}\par}
    \end{minipage}\hfill
    {\poplogo\fontsize{22}{22}\selectfont\color{popcream}OnBuch\color{poporange}+}
  \end{tcolorbox}%
}

% ============================================================
% Page de garde
% #1 titre · #2 matière · #3 niveau · #4 module · #5 n° leçon · #6 durée ·
% #7 description / objectifs (texte ou itemize)
% ============================================================
\newcommand{\pagegarde}[7]{%
  \thispagestyle{empty}
  \newgeometry{a4paper,margin=1.7cm,top=1.6cm,bottom=1.4cm}%
  \begin{tikzpicture}[remember picture,overlay]
    \fill[poporange!18] ([xshift=-3.2cm,yshift=-3.6cm]current page.north east) circle (4.6cm);
    \fill[poppurple!14] ([xshift=2.4cm,yshift=7.6cm]current page.south west) circle (3.8cm);
  \end{tikzpicture}%
  {\poplogo\fontsize{30}{30}\selectfont\color{popink}OnBuch\color{poporange}+}\hfill
  \raisebox{0.6ex}{\poppill{Cours signé \textbullet\ Premium}{popink}{popcream}}\par
  \vspace{1pt}\popsquiggle{poppurple}\par
  \vspace{8pt}
  \begin{tcolorbox}[enhanced,colback=poporange,colframe=popink,boxrule=2.8pt,arc=13pt,
      drop shadow=popink,left=18pt,right=18pt,top=12pt,bottom=13pt,after skip=10pt]
    \poppill{#2 \textbullet\ #3}{popink}{popcream}\hspace{5pt}\poppill{#4}{white}{popink}\par\vspace{8pt}
    {\popdisplay\fontsize{14}{14}\selectfont\color{popink}LEÇON #5}\par\vspace{4pt}
    {\raggedright\hyphenpenalty=10000\exhyphenpenalty=10000\popdisplay\fontsize{25}{27}\selectfont\color{popcream}\MakeUppercase{#1}\par}
    \vspace{8pt}
    {\popxbold\fontsize{9.5}{9.5}\selectfont\color{popink}\MakeUppercase{Durée conseillée : #6}}
  \end{tcolorbox}
  \begin{tcolorbox}[enhanced,colback=white,colframe=popink,boxrule=2.2pt,arc=10pt,
      drop shadow=popink,left=16pt,right=16pt,top=10pt,bottom=10pt,after skip=8pt]
    \poppill{Dans cette leçon}{poppurple}{white}\par\vspace{5pt}
    \popsemi\fontsize{9.3}{12.6}\selectfont\color{popink}#7
  \end{tcolorbox}
  \noindent\begin{minipage}[t]{0.487\linewidth}
    \begin{tcolorbox}[enhanced,colback=popgreen,colframe=popink,boxrule=2.2pt,arc=10pt,
        drop shadow=popink,left=11pt,right=11pt,top=10pt,bottom=10pt,height=3.1cm,valign=center]
      \tcbox[colback=white,colframe=popink,boxrule=1.3pt,arc=5pt,boxsep=0pt,nobeforeafter,
        left=3pt,right=3pt,top=3pt,bottom=3pt,valign=center]{\qrcode[height=1.9cm]{https://play.google.com/store/apps/details?id=com.luvvix.onbuch&hl=fr}}%
      \hspace{8pt}\begin{minipage}[c]{0.5\linewidth}
        {\popdisplay\fontsize{11}{12.5}\selectfont\color{white}\MakeUppercase{Télécharge OnBuch}}\par\vspace{4pt}
        {\raggedright\popsemi\fontsize{8.4}{11}\selectfont\color{white}Gratuit \textbullet\ Google Play\\Tuteur IA, annales, concours, résultats\par}
      \end{minipage}
    \end{tcolorbox}
  \end{minipage}\hfill
  \begin{minipage}[t]{0.487\linewidth}
    \begin{tcolorbox}[enhanced,colback=poppurple,colframe=popink,boxrule=2.2pt,arc=10pt,
        drop shadow=popink,left=11pt,right=11pt,top=10pt,bottom=10pt,height=3.1cm,valign=center]
      \tikz[baseline=-0.7ex]\node[circle,fill=white,draw=popink,line width=1.3pt,minimum size=1.6cm,inner sep=0]{\popdisplay\fontsize{24}{24}\selectfont\color{poppurple}+};%
      \hspace{8pt}\begin{minipage}[c]{0.6\linewidth}
        {\popdisplay\fontsize{11}{12.5}\selectfont\color{white}\MakeUppercase{Passe à OnBuch+}}\par\vspace{4pt}
        {\raggedright\popsemi\fontsize{8.4}{11}\selectfont\color{white}Tous les cours signés, Tuteur IA illimité, fiches PDF — onglet OnBuch+ de l'app\par}
      \end{minipage}
    \end{tcolorbox}
  \end{minipage}
  \vspace{8pt}
  \popcontenu
  \vfill
  \signatureonbuch
  \restoregeometry
  \newpage
}

% Rangée « Ce cours contient » (page de garde)
\newcommand{\poptile}[3]{% #1 couleur #2 gros texte #3 légende
  \begin{tcolorbox}[enhanced,colback=#1,colframe=popink,boxrule=2pt,arc=9pt,shadow={2.5pt}{-2.5pt}{0pt}{popink},
      left=6pt,right=6pt,top=9pt,bottom=9pt,halign=center,valign=center,height=1.75cm,width=\linewidth,nobeforeafter]
    {\popdisplay\fontsize{12.5}{13}\selectfont\color{white}\MakeUppercase{#2}}\par\vspace{3pt}
    {\centering\popsemi\fontsize{7.8}{9.5}\selectfont\color{white}#3\par}
  \end{tcolorbox}}
\newcommand{\popcontenu}{%
  \par\noindent{\popxboldls\fontsize{7.5}{7.5}\selectfont\color{popmuted}CE COURS SIGNÉ CONTIENT}\par\vspace{7pt}
  \noindent\begin{minipage}[t]{0.235\linewidth}\poptile{popblue}{Cours}{complet, illustré, pas à pas}\end{minipage}\hfill
  \begin{minipage}[t]{0.235\linewidth}\poptile{poporange}{Méthodes}{et exercices résolus}\end{minipage}\hfill
  \begin{minipage}[t]{0.235\linewidth}\poptile{poppink}{Exercices}{type examen + corrigés}\end{minipage}\hfill
  \begin{minipage}[t]{0.235\linewidth}\poptile{popgreen}{Fiche bilan}{l'essentiel à retenir}\end{minipage}\par}

% Sommaire
\newtcolorbox{sommaire}{enhanced,colback=white,colframe=popink,boxrule=2.2pt,arc=10pt,
  drop shadow=popink,left=18pt,right=18pt,top=13pt,bottom=13pt,before skip=6pt,after skip=20pt,
  before upper={\poppill{Sommaire}{poppurple}{white}\par\vspace{9pt}\popsemi\fontsize{10.6}{14.5}\selectfont\color{popink}}}

% Signature de fin de cours — poussée tout en bas de la dernière page
\newcommand{\findecours}{%
  \par\vfill\nopagebreak
  \begin{center}\popsquiggle{poporange}\end{center}\vspace{4pt}
  \signatureonbuch
}

%% FIGFIX-BEGIN v3 — correctifs globaux des figures (docs/onbuch-generation/tools/figfix)
\usepackage{adjustbox}\usepackage{etoolbox}
% toute figure TikZ/pgfplots plus large que la ligne est réduite à la largeur disponible
\BeforeBeginEnvironment{tikzpicture}{\begin{adjustbox}{max width=\linewidth}}
\AfterEndEnvironment{tikzpicture}{\end{adjustbox}}
% légendes pgfplots lisibles par-dessus les courbes
\pgfplotsset{every axis/.append style={legend style={fill=white,fill opacity=0.92,text opacity=1,draw=black!15,inner sep=3pt}}}
% étiquettes d'arêtes (syntaxe edge["..."]) avec fond blanc
\tikzset{every edge quotes/.append style={fill=white,inner sep=1.2pt}}
% bulles de cartes mentales : texte à la ligne dans la bulle, bulle un peu plus grande
\tikzset{every concept/.append style={text width=2.0cm,align=center,minimum size=2.3cm,inner sep=2pt}}
%% FIGFIX-END
\begin{document}

\pagegarde{Mouvement des planètes et des satellites}{Physique}{Tle D}{Module 2 : Mouvements et interactions — évolutions temporelles des systèmes mécaniques}{P07}{5 h}{Cette leçon étudie le mouvement des planètes et des satellites dans le champ de gravitation d'un astre central. Elle établit les lois de Kepler, les expressions de la vitesse et de la période d'un satellite, la condition géostationnaire et l'énergie mécanique d'un satellite, outils indispensables pour comprendre l'exploration spatiale et réussir l'épreuve de physique au Baccalauréat C.
\vspace{4pt}
\begin{multicols}{2}\small
\begin{itemize}
\item À la fin de cette leçon, je sais définir les référentiels géocentrique et héliocentrique et choisir le référentiel adapté à l'étude d'un mouvement
\item je sais énoncer les trois lois de Kepler et les exploiter
\item je sais mettre en équation le mouvement circulaire d'un satellite sous l'action de la force de gravitation
\item je sais établir et exploiter les expressions de la vitesse et de la période d'un satellite en orbite circulaire
\item je sais démontrer la troisième loi de Kepler dans le cas d'une orbite circulaire
\item je sais calculer l'altitude d'un satellite géostationnaire et justifier ses conditions d'orbite
\end{itemize}\end{multicols}}

\begin{sommaire}
\begin{multicols}{2}
\begin{enumerate}
\item Référentiels d'étude et lois de Kepler
\item Mouvement circulaire uniforme d'un satellite
\item Le satellite géostationnaire
\item Énergie mécanique d'un satellite (complément Bac)
\item Synthèse et méthodes de résolution
\item Activité d'intégration
\item Méthodes et exercices résolus
\item Exercices
\item Corrigés des exercices
\item Fiche bilan
\end{enumerate}
\end{multicols}
\end{sommaire}

\begin{objectifs}
\begin{itemize}
\item À la fin de cette leçon, je sais définir les référentiels géocentrique et héliocentrique et choisir le référentiel adapté à l'étude d'un mouvement
\item je sais énoncer les trois lois de Kepler et les exploiter
\item je sais mettre en équation le mouvement circulaire d'un satellite sous l'action de la force de gravitation
\item je sais établir et exploiter les expressions de la vitesse et de la période d'un satellite en orbite circulaire
\item je sais démontrer la troisième loi de Kepler dans le cas d'une orbite circulaire
\item je sais calculer l'altitude d'un satellite géostationnaire et justifier ses conditions d'orbite
\item je sais déterminer la masse d'un astre à partir du mouvement d'un de ses satellites
\item je sais exprimer l'énergie mécanique d'un satellite et définir la vitesse de libération (complément Bac)
\end{itemize}
\end{objectifs}

\begin{prerequis}
\begin{itemize}
\item Force d'interaction gravitationnelle et champ de gravitation (Première C)
\item Mouvement circulaire uniforme : vitesse angulaire, période, accélération normale a = v²/r
\item Théorème du centre d'inertie (2e loi de Newton)
\item Énergie cinétique, énergie potentielle, énergie mécanique (Première C)
\item Notion de référentiel galiléen
\end{itemize}
\end{prerequis}

\newpage

\coursec{Référentiels d'étude et lois de Kepler}

Chaque soir, des millions de Camerounais regardent Canal+ ou suivent les matchs des Lions Indomptables grâce à des paraboles pointées vers un point fixe du ciel. Observe bien : la parabole, solidement fixée au mur de la maison, ne bouge jamais. Pourtant, elle capte un satellite qui tourne autour de la Terre ! Quelque part à environ \SI{36000}{km} au-dessus de l'équateur, ce satellite tourne exactement au même rythme que la Terre, de sorte qu'il semble \emph{immobile} vu depuis le sol. Comment les ingénieurs ont-ils calculé cette altitude ? Et comment les astronomes ont-ils \og pesé \fg{} le Soleil sans jamais y poser une balance ?

Toute la réponse repose sur deux piliers : le choix d'un bon \cle{référentiel} pour décrire le mouvement, et trois lois empiriques établies par Johannes Kepler au début du XVII\textsuperscript{e} siècle. C'est ce que nous allons découvrir dans cette première section. La problématique qui guidera toute la leçon est la suivante : \emph{comment décrire, prévoir et exploiter le mouvement des planètes et des satellites dans un champ de gravitation ?}

\courssub{La notion de référentiel : un rappel indispensable}

Imaginons un voyageur assis dans le bus Yaoundé--Douala roulant à \SI{80}{km.h^{-1}}. Pour son voisin de siège, il est immobile ; pour un piéton au bord de la route, il file à \SI{80}{km.h^{-1}}. Qui a raison ? Les deux ! La description d'un mouvement dépend de l'objet par rapport auquel on l'observe. Cet objet de référence s'appelle un \cle{référentiel}.

\begin{definition}[Référentiel]
Un \cle{référentiel} est un solide de référence par rapport auquel on décrit le mouvement d'un corps. On lui associe un repère d'espace (origine et trois axes) et une horloge pour mesurer le temps. Toute étude de mouvement doit commencer par la phrase : \og le mouvement est étudié dans le référentiel \ldots \fg{}.
\end{definition}

\begin{attention}[Réflexe obligatoire]
Au Baccalauréat, oublier de préciser le référentiel avant un bilan des forces ou l'application de la deuxième loi de Newton est une erreur qui coûte des points. Écris systématiquement : \og Dans le référentiel \ldots, supposé galiléen \ldots\fg{}.
\end{attention}

Un référentiel est dit \cle{galiléen} si le principe d'inertie s'y applique : un corps isolé ou pseudo-isolé y est soit au repos, soit animé d'un mouvement rectiligne uniforme. C'est uniquement dans un référentiel galiléen que l'on peut appliquer la deuxième loi de Newton sous la forme $\sum \vec{F} = m\,\vec{a}$.

\courssub{Le référentiel héliocentrique}

Pour étudier le mouvement des planètes autour du Soleil, il serait absurde de se placer dans un référentiel lié au sol camerounais : la Terre tourne sur elle-même et autour du Soleil, ce qui compliquerait énormément les trajectoires. Les astronomes utilisent le référentiel le plus adapté.

\begin{definition}[Référentiel héliocentrique]
Le \cle{référentiel héliocentrique} (ou référentiel de Copernic) a pour origine le centre d'inertie du Soleil. Ses trois axes pointent vers trois étoiles très lointaines, dites \og fixes \fg{} car leur direction apparente ne varie pratiquement pas à l'échelle humaine. Il est adapté à l'étude du mouvement des planètes et des comètes du système solaire.
\end{definition}

Pourquoi ce référentiel est-il galiléen (au moins de façon approchée) ? Le Soleil contient à lui seul environ \SI{99,86}{\%} de la masse du système solaire. Il est donc quasiment isolé des autres astres et presque immobile par rapport à l'ensemble du système solaire. Les étoiles \og fixes \fg{} étant à des distances immenses (des années-lumière), les axes gardent une direction constante. Le référentiel héliocentrique est donc un excellent référentiel galiléen pour les durées d'étude habituelles.

\courssub{Le référentiel géocentrique}

Pour étudier un satellite de télécommunications ou la Lune, le référentiel héliocentrique n'est pas commode : il faudrait y décrire à la fois la révolution de la Terre autour du Soleil et celle du satellite autour de la Terre. On préfère un référentiel \og centré sur la Terre \fg{} mais qui ne tourne pas avec elle.

\begin{definition}[Référentiel géocentrique]
Le \cle{référentiel géocentrique} a pour origine le centre d'inertie de la Terre. Ses trois axes restent constamment parallèles à ceux du référentiel héliocentrique (ils pointent donc vers les mêmes étoiles lointaines). Il est adapté à l'étude des satellites terrestres et de la Lune.
\end{definition}

\begin{attention}[Géocentrique ne veut pas dire terrestre !]
Le référentiel géocentrique \textbf{ne tourne pas} avec la Terre : ses axes gardent une direction fixe dans l'espace. Ne le confonds pas avec le référentiel terrestre (lié au sol, celui de ton laboratoire), qui tourne avec la Terre en \SI{24}{h}. C'est précisément parce que le référentiel géocentrique ne tourne pas qu'un satellite géostationnaire, immobile dans le référentiel terrestre, y décrit un cercle en \SI{23}{h} \SI{56}{min} \SI{4}{s} (jour sidéral).
\end{attention}

Le référentiel géocentrique n'est pas rigoureusement galiléen : la Terre est en mouvement accéléré autour du Soleil (accélération d'environ \SI{6e-3}{m.s^{-2}}). Mais pour un satellite proche de la Terre, sur une durée courte devant un an, l'attraction terrestre domine largement et l'approximation galiléenne est excellente. Nous l'admettrons.

\begin{popfigure}
\begin{center}
\begin{tikzpicture}[scale=0.9]
% Soleil
\shade[ball color=popgold] (0,0) circle (0.55);
\node[below=0.6cm of {(0,0)}, popdark] {\small Soleil};
% Axes héliocentriques
\draw[-{Stealth}, popink, thick] (0,0) -- (2.6,0) node[right] {\small $x_h$};
\draw[-{Stealth}, popink, thick] (0,0) -- (0,2.2) node[above] {\small $y_h$};
\draw[-{Stealth}, popink, thick] (0,0) -- (-1.3,-1.1) node[below left] {\small $z_h$};
% Orbite terrestre
\draw[popline, dashed] (0,0) circle (3.4);
% Terre
\shade[ball color=popblue] (3.4,0) circle (0.28);
\node[below=0.32cm of {(3.4,0)}, popdark] {\small Terre};
% Axes géocentriques
\draw[-{Stealth}, poporange, thick] (3.4,0) -- (5.6,0) node[right] {\small $x_g$};
\draw[-{Stealth}, poporange, thick] (3.4,0) -- (3.4,1.9) node[above] {\small $y_g$};
\draw[-{Stealth}, poporange, thick] (3.4,0) -- (2.4,-0.85) node[below left] {\small $z_g$};
% Parallélisme
\node[popmuted] at (1.7,-2.3) {\small $(x_g) \parallel (x_h)$,\quad $(y_g) \parallel (y_h)$,\quad $(z_g) \parallel (z_h)$};
\end{tikzpicture}
\end{center}
\legende{Le référentiel géocentrique (axes oranges) a son origine au centre de la Terre et ses axes restent parallèles à ceux du référentiel héliocentrique (axes bleus) au cours de la révolution de la Terre autour du Soleil.}
\end{popfigure}

\courssub{Première loi de Kepler : la loi des orbites}

\textbf{Une intuition d'abord.} Pendant des siècles, on a cru que les planètes décrivaient des cercles parfaits, le cercle étant la figure \og divine \fg{}. En analysant les observations extrêmement précises de l'astronome danois Tycho Brahe, Kepler s'aperçut que l'orbite de Mars ne fermait pas tout à fait : il manquait 8 minutes d'arc, l'épaisseur d'un ongle vu à bout de bras ! Plutôt que de rejeter les mesures, Kepler rejeta le cercle et adopta l'\cle{ellipse}.

\begin{propriete}[Première loi de Kepler — loi des orbites (admise)]
Dans le référentiel héliocentrique, chaque planète décrit une \cle{ellipse} dont le Soleil occupe l'un des deux \cle{foyers}. De même, dans le référentiel géocentrique, un satellite terrestre décrit une ellipse dont la Terre occupe un foyer. Le \cle{cercle} est un cas particulier de l'ellipse, lorsque les deux foyers sont confondus avec le centre.
\end{propriete}

Quelques éléments de vocabulaire indispensables :

\begin{itemize}
\item Le \cle{demi-grand axe}, noté $a$, est la moitié du plus grand diamètre de l'ellipse. Pour un cercle, $a$ se confond avec le rayon $r$.
\item Pour une planète : le point de l'orbite le plus proche du Soleil est le \cle{périhélie}, le plus éloigné est l'\cle{aphélie}.
\item Pour un satellite terrestre : le point le plus proche de la Terre est le \cle{périgée}, le plus éloigné est l'\cle{apogée}.
\item La \cle{période de révolution} $T$ est la durée d'un tour complet.
\end{itemize}

\courssub{Deuxième loi de Kepler : la loi des aires}

\textbf{Observation.} La comète de Halley file très vite quand elle frôle le Soleil, puis ralentit énormément quand elle s'en éloigne vers les confins du système solaire. Sa vitesse n'est donc pas constante. Kepler trouva la loi cachée derrière ce comportement : ce n'est pas la vitesse qui se conserve, mais l'\emph{aire balayée par unité de temps}.

\begin{propriete}[Deuxième loi de Kepler — loi des aires (admise)]
Le segment joignant le Soleil à la planète (le \cle{rayon vecteur}) balaie des \cle{aires égales pendant des durées égales}. Conséquence immédiate : sur une orbite elliptique, la vitesse de la planète n'est \textbf{pas constante} ; elle est maximale au périhélie et minimale à l'aphélie. Sur une orbite \textbf{circulaire}, la loi des aires impose en revanche une vitesse constante : le mouvement est circulaire \emph{uniforme}.
\end{propriete}

\begin{popfigure}
\begin{center}
\begin{tikzpicture}[scale=1.05]
% Ellipse
\draw[popblue, thick] (0,0) ellipse (3.2 and 1.9);
% Foyer (Soleil)
\coordinate (F) at (2.55,0);
\shade[ball color=popgold] (F) circle (0.16);
\node[below=0.2cm of F, popdark] {\small Soleil (foyer $F$)};
% Autre foyer
\fill[popmuted] (-2.55,0) circle (0.06);
\node[below, popmuted] at (-2.55,0) {\small $F'$};
% Demi-grand axe
\draw[{Stealth}-{Stealth}, popink] (0,-0.12) -- (3.2,-0.12) node[midway, below] {\small $a$};
% Planète proche du périhélie : deux positions
\coordinate (P1) at (3.05,0.55);
\coordinate (P2) at (2.75,0.98);
\fill[poppink!60] (F) -- (P1) arc[start angle=10, end angle=19.6, x radius=3.2, y radius=1.9] -- cycle;
\draw[popdark] (F) -- (P1);
\draw[popdark] (F) -- (P2);
\fill[popblue] (P1) circle (0.09);
\fill[popblue] (P2) circle (0.09);
\node[right, popdark] at (3.0,0.8) {\small $\mathcal{A}_1$};
% Planète près de l'aphélie
\coordinate (A1) at (-3.15,0.42);
\coordinate (A2) at (-3.19,-0.18);
\fill[popgreen!50] (F) -- (A1) arc[start angle=172.4, end angle=183.2, x radius=3.2, y radius=1.9] -- cycle;
\draw[popdark] (F) -- (A1);
\draw[popdark] (F) -- (A2);
\fill[popblue] (A1) circle (0.09);
\fill[popblue] (A2) circle (0.09);
\node[left, popdark] at (-3.1,0.35) {\small $\mathcal{A}_2$};
% Vitesses
\draw[-{Stealth}, poporange, very thick] (P1) -- ++(0.25,0.75) node[right] {\small $\vec{v}_{\text{périhélie}}$ (grande)};
\draw[-{Stealth}, poporange, very thick] (A1) -- ++(0.05,0.45) node[above left=-0.1cm] {\small $\vec{v}_{\text{aphélie}}$ (faible)};
\node[popdark] at (0,-2.5) {\small Même durée $\Delta t$ \; $\Rightarrow$ \; $\mathcal{A}_1 = \mathcal{A}_2$};
\end{tikzpicture}
\end{center}
\legende{Loi des aires : pendant la même durée $\Delta t$, le rayon vecteur balaie deux aires égales $\mathcal{A}_1 = \mathcal{A}_2$. Près du Soleil, la planète parcourt un arc plus long : elle va donc plus vite.}
\end{popfigure}

\courssub{Troisième loi de Kepler : la loi des périodes}

\textbf{Intuition.} Plus une planète est loin du Soleil, plus son orbite est grande et plus elle va lentement : son année est plus longue. Mercure boucle sa révolution en 88 jours, Neptune en 165 ans ! Kepler chercha pendant dix ans la relation mathématique exacte entre la période $T$ et la taille $a$ de l'orbite. Il la publia en 1619.

\begin{propriete}[Troisième loi de Kepler — loi des périodes (admise pour l'instant)]
Pour tous les satellites (ou planètes) gravitant autour d'un \textbf{même astre attracteur}, le rapport du carré de la période de révolution au cube du demi-grand axe de l'orbite est constant :
\[ \frac{T^2}{a^3} = k \]
où $k$ est une constante qui ne dépend \cle{que de l'astre attracteur} (sa masse), et non du satellite. Pour une orbite circulaire de rayon $r$, on écrit $\dfrac{T^2}{r^3} = k$. Cette loi sera \textbf{démontrée en section 2} dans le cas circulaire, à partir de la loi de gravitation de Newton ; on y montrera que $k = \dfrac{4\pi^2}{G\,M_A}$ où $M_A$ est la masse de l'astre attracteur.
\end{propriete}

\begin{attention}[Le piège classique de la constante $k$]
Dans $\dfrac{T^2}{a^3} = k$, la constante $k$ dépend \textbf{uniquement de l'astre attracteur}, jamais du satellite. Ainsi :
\begin{itemize}
\item on peut comparer la Terre et Mars (même attracteur : le Soleil) : $\dfrac{T_T^2}{a_T^3} = \dfrac{T_M^2}{a_M^3}$ ;
\item on peut comparer deux satellites terrestres (même attracteur : la Terre) ;
\item mais il est \textbf{interdit} d'écrire $\dfrac{T_{\text{Lune}}^2}{a_{\text{Lune}}^3} = \dfrac{T_{\text{Terre}}^2}{a_{\text{Terre}}^3}$, car la Lune tourne autour de la Terre et la Terre autour du Soleil : les attracteurs sont différents !
\end{itemize}
Autre piège : l'homogénéité. $T^2/a^3$ s'exprime en \si{s^2.m^{-3}} dans le système international ; si tu utilises les années et les unités astronomiques, la valeur numérique de $k$ change, mais la loi reste valable.
\end{attention}

\begin{exemplebox}[Vérification de la troisième loi avec la Terre et Mars]
Prenons comme unités l'année (an) et l'unité astronomique (u.a., distance moyenne Terre--Soleil).
\begin{itemize}
\item Terre : $a_T = 1{,}00$ u.a. et $T_T = 1{,}00$ an, donc $\dfrac{T_T^2}{a_T^3} = \dfrac{1{,}00^2}{1{,}00^3} = 1{,}00$ an$^2$/u.a.$^3$.
\item Mars : $a_M = 1{,}52$ u.a. et $T_M = 1{,}88$ an, donc $\dfrac{T_M^2}{a_M^3} = \dfrac{1{,}88^2}{1{,}52^3} = \dfrac{3{,}53}{3{,}51} \approx 1{,}01$ an$^2$/u.a.$^3$.
\end{itemize}
Les deux rapports sont égaux (à la précision des données près) : la troisième loi de Kepler est vérifiée. Réciproquement, connaissant $a_M = 1{,}52$ u.a., on peut \emph{prévoir} la période de Mars : $T_M = a_M^{3/2} = 1{,}52^{1{,}5} \approx 1{,}87$ an.
\end{exemplebox}

Le graphique ci-dessous, tracé pour huit planètes du système solaire, montre $T^2$ en fonction de $a^3$ : les points sont alignés sur une droite passant par l'origine, signature d'une proportionnalité.

\begin{popfigure}
\begin{center}
\begin{tikzpicture}
\begin{axis}[
width=0.85\linewidth, height=6.5cm,
xlabel={$a^3$ (u.a.$^3$)}, ylabel={$T^2$ (an$^2$)},
xmin=0, xmax=28000, ymin=0, ymax=28000,
grid=both, grid style={popline!40},
legend style={at={(0.03,0.97)}, anchor=north west, font=\small},
]
\addplot[popcurve, domain=0:28000] {x};
\addlegendentry{$T^2 = a^3$ (droite de Kepler)}
\addplot[only marks, mark=*, mark size=1.8pt, poporange] coordinates {
(0.058,0.058) (0.39,0.39) (1,1) (3.5,3.5) (140,140) (860,860) (19000,19000) (27000,27000)
};
\addlegendentry{Planètes (Mercure à Neptune)}
\end{axis}
\end{tikzpicture}
\end{center}
\legende{Pour les planètes du système solaire, $T^2$ est proportionnel à $a^3$ : les points s'alignent sur une droite passant par l'origine, dont la pente est la constante $k$ du Soleil.}
\end{popfigure}

\begin{savaistu}[Kepler, Tycho Brahe et Newton]
Johannes Kepler (1571--1630) établit ses deux premières lois en 1609 et la troisième en 1619, en dépouillant pendant des années les observations de Mars accumulées par l'astronome danois Tycho Brahe, qui mesurait les positions des astres \emph{à l'œil nu} avec une précision remarquable. Près de 70 ans plus tard, en 1687, Isaac Newton montra dans ses \emph{Principia} que ces trois lois empiriques découlent d'une seule loi fondamentale : la \cle{gravitation universelle}. C'est l'un des plus beaux exemples de l'histoire des sciences : l'observation patiente, puis la loi empirique, puis la théorie unificatrice. Aujourd'hui, la troisième loi de Kepler permet de \og peser \fg{} les astres : en mesurant la période et le rayon de l'orbite d'une lune de Jupiter, on calcule la masse de Jupiter — c'est ainsi que les astronomes déterminent la masse des étoiles et même des trous noirs !
\end{savaistu}

\begin{exoresolu}[Période d'un satellite de télédétection]
Un satellite camerounais d'observation de la Terre (projet spatial africain) décrit une orbite circulaire de rayon $r = \SI{7,0e3}{km}$ autour de la Terre. La Lune, satellite naturel de la Terre, a une orbite de rayon $r_L = \SI{3,84e5}{km}$ et une période $T_L = \SI{27,3}{jours}$. Calculer la période $T$ du satellite en heures.
\tcblower
\textbf{Méthode :} le satellite et la Lune gravitent autour du \emph{même} astre attracteur (la Terre) : on peut appliquer la troisième loi de Kepler.
\begin{align*}
\frac{T^2}{r^3} &= \frac{T_L^2}{r_L^3} \\
T^2 &= T_L^2 \times \frac{r^3}{r_L^3} \\
T &= T_L \times \left(\frac{r}{r_L}\right)^{3/2}
\end{align*}
Application numérique :
\[ T = 27{,}3 \times \left(\frac{7{,}0\times 10^{3}}{3{,}84\times 10^{5}}\right)^{3/2} = 27{,}3 \times (1{,}823\times 10^{-2})^{1{,}5} \approx 27{,}3 \times 2{,}46\times 10^{-3} \approx 6{,}72\times 10^{-2}\ \text{jour} \]
soit $T \approx 6{,}72\times 10^{-2} \times 24 \approx \SI{1,6}{h}$. Le satellite boucle son orbite en environ \SI{1,6}{h}, soit près de 15 tours par jour : c'est cohérent, car il est beaucoup plus proche de la Terre que la Lune.
\end{exoresolu}

\begin{aretenir}[L'essentiel de la section]
\begin{itemize}
\item Toute étude de mouvement commence par le choix du référentiel : \cle{héliocentrique} (centre du Soleil, axes vers des étoiles fixes) pour les planètes ; \cle{géocentrique} (centre de la Terre, axes parallèles à ceux de l'héliocentrique) pour les satellites terrestres et la Lune. Tous deux sont supposés galiléens.
\item \textbf{1\textsuperscript{re} loi de Kepler} : orbites elliptiques, l'astre attracteur occupant un foyer (le cercle en est un cas particulier).
\item \textbf{2\textsuperscript{e} loi de Kepler} : aires égales en durées égales ; la vitesse varie sur une ellipse, elle est constante sur un cercle.
\item \textbf{3\textsuperscript{e} loi de Kepler} : $\dfrac{T^2}{a^3} = k$, où $k$ ne dépend que de la masse de l'\cle{astre attracteur}.
\end{itemize}
\end{aretenir}

Dans la section suivante, nous quitterons le terrain de l'observation pour celui de la démonstration : armés de la loi de gravitation universelle et de la deuxième loi de Newton, nous établirons les expressions de la vitesse et de la période d'un satellite en orbite circulaire, et nous retrouverons la troisième loi de Kepler comme une conséquence de la physique newtonienne.

\coursec{Mouvement circulaire uniforme d'un satellite}

Dans la section précédente, nous avons vu que les planètes et les satellites obéissent aux lois de Kepler, découvertes par l'observation. Nous allons maintenant \emph{comprendre} l'origine physique de ces lois en modélisant le mouvement d'un satellite artificiel ou naturel en orbite circulaire autour d'un astre massif. C'est la modélisation la plus simple, mais elle rend compte de l'essentiel de la dynamique spatiale : du lancement de l'ISS à la mise en orbite des satellites de télécommunication camerounais (exploités via CAMTEL ou les opérateurs MTN/Orange), en passant par le mouvement de la Lune autour de la Terre.

\courssub{Modélisation et bilan des forces}

Nous considérons un satellite $S$ de masse $m$ (par exemple un satellite de communication, l'ISS ou la Lune) en orbite autour d'un astre central $A$ de masse $M$ (la Terre, le Soleil, etc.). Nous supposons $M \gg m$ : le centre d'inertie du système $\{A, S\}$ est confondu avec le centre de l'astre $A$. Le référentiel d'étude est le \textbf{référentiel astrocentrique} (lié au centre de l'astre $A$, sans rotation propre), que l'on suppose \textbf{galiléen} (les forces d'inertie sont négligeables).

\begin{definition}[Modèle du satellite en orbite circulaire]
Un satellite est modélisé par un point matériel de masse $m$ décrivant une trajectoire circulaire de rayon $r$ (distance centre-centre) autour d'un astre de masse $M$, sous l'action exclusive de la force de gravitation newtonienne.
\end{definition}

\begin{popfigure}
\begin{tikzpicture}[scale=1.2, >=Stealth]
    % Astre central
    \fill[popblue!20] (0,0) circle (1.2cm);
    \draw[popblue, thick] (0,0) circle (1.2cm);
    \node[popblue, font=\bfseries] at (0,0) {Astre $A$ ($M$)};
    \node[popblue, font=\small] at (0,-1.5) {Centre $O$};
    
    % Orbite
    \draw[popmuted, dashed, thick] (0,0) circle (3.5cm);
    \node[popmuted, font=\small] at (3.8, 2.5) {Orbite circulaire rayon $r$};
    
    % Satellite
    \coordinate (S) at (45:3.5cm);
    \fill[poporange!80] (S) circle (0.25cm);
    \draw[poporange, thick] (S) circle (0.25cm);
    \node[poporange, font=\bfseries, anchor=south west] at (S) {Satellite $S$ ($m$)};
    
    % Vecteur position
    \draw[popdark, ->, thick] (0,0) -- (S) node[midway, above left, font=\small] {$\vec{r}$};
    
    % Base de Frenet (u_r sortant, u_theta direct)
    \draw[popgreen, ->, thick] (S) -- ++(45:1.2) node[above right, popgreen, font=\small] {$\vec{u}_r$ (radial sortant)};
    \draw[poppurple, ->, thick] (S) -- ++(-45:1.2) node[below right, poppurple, font=\small] {$\vec{u}_\theta$ (transversal direct)};
    
    % Vitesse
    \draw[poppink, ->, very thick] (S) -- ++(-45:1.5) node[below right, poppink, font=\bfseries\small] {$\vec{v} = v\,\vec{u}_\theta$};
    
    % Force de gravitation
    \draw[popred, ->, very thick] (S) -- ++(225:1.8) node[below left, popred, font=\bfseries\small] {$\vec{F}_g = -\dfrac{GMm}{r^2}\,\vec{u}_r$};
    
    % Angle
    \draw[popmuted, ->] (1,0) arc (0:45:1) node[midway, right, font=\small] {$\theta$};
\end{tikzpicture}
\legende{Satellite $S$ de masse $m$ en orbite circulaire de rayon $r$ autour de l'astre $A$ de masse $M$. La force de gravitation $\vec{F}_g$ est centripète (opposée à $\vec{u}_r$) ; la vitesse $\vec{v}$ est tangentielle (along $\vec{u}_\theta$). Base de Frenet $(\vec{u}_r, \vec{u}_\theta)$ avec $\vec{u}_r$ radial sortant.}
\end{popfigure}

La seule force agissant sur le satellite est la force de gravitation exercée par l'astre $A$ :
\[
\vec{F}_g = -\frac{G M m}{r^2}\,\vec{u}_r
\]
où $G \approx \SI{6,67e-11}{N.m^2.kg^{-2}}$ est la constante de gravitation universelle, $r$ la distance entre les centres de masse, et $\vec{u}_r$ le vecteur unitaire radial \textbf{sortant}. Cette force est \textbf{centripète} (dirigée vers le centre $O$, donc opposée à $\vec{u}_r$) et sa norme est constante car $r$ est constant.

\courssub{Application du théorème du centre d'inertie en base de Frenet}

Le théorème du centre d'inertie (deuxième loi de Newton) s'écrit $\sum \vec{F}_{\text{ext}} = m \vec{a}$. En base de Frenet pour un mouvement circulaire de rayon $r$ constant, l'accélération s'écrit (avec $\vec{u}_r$ radial sortant) :
\[
\vec{a} = \underbrace{\frac{dv}{dt}}_{a_\tau}\,\vec{u}_\theta \underbrace{- \frac{v^2}{r}}_{a_n}\,\vec{u}_r
\]
La force de gravitation n'a \textbf{aucune composante tangentielle} : $\vec{F}_g \cdot \vec{u}_\theta = 0$. Le projeté du théorème du centre d'inertie sur l'axe tangentiel ($\vec{u}_\theta$) donne donc :
\[
m\,\frac{dv}{dt} = 0 \quad \Longrightarrow \quad \frac{dv}{dt} = 0 \quad \Longrightarrow \quad v = \text{constante}.
\]

\begin{propriete}[Mouvement circulaire uniforme imposé par la gravitation]
\emph{Démonstration exigible au Bac.} Un satellite soumis uniquement à la force de gravitation centrale d'un astre, et dont la trajectoire initiale est circulaire, est animé d'un \textbf{mouvement circulaire uniforme} (MCU). L'accélération tangentielle est nulle car la force est purement radiale ; la norme de la vitesse $v$ ne change pas.
\end{propriete}

Le projeté sur l'axe radial ($\vec{u}_r$, orienté vers l'extérieur) donne l'équation de la dynamique radiale :
\[
-m\,\frac{v^2}{r} = -\frac{G M m}{r^2}
\]
Les deux membres sont négatifs car l'accélération normale (centripète) et la force sont dirigées vers le centre (opposées à $\vec{u}_r$). En simplifiant par $m$ (la masse du satellite \emph{disparaît} !) et en multipliant par $r$, on obtient la \textbf{relation fondamentale} :
\[
v^2 = \frac{G M}{r}
\]

\begin{propriete}[Vitesse orbitale en orbite circulaire]
Pour un satellite de masse $m$ (négligeable devant $M$) en orbite circulaire de rayon $r$ autour d'un astre de masse $M$ :
\[
\boxed{v = \sqrt{\frac{G M}{r}}}
\]
\begin{itemize}
    \item $v$ ne dépend \textbf{que de $M$ et $r$}, pas de la masse $m$ du satellite (équivalence des masses inertielle et pesante).
    \item Plus l'orbite est basse ($r$ petit), plus la vitesse est grande ($v \propto 1/\sqrt{r}$).
\end{itemize}
\end{propriete}

\begin{attention}[Piège classique : $r$ vs $h$]
\emph{Erreur fréquente au Bac.} Le rayon de l'orbite $r$ est la distance \textbf{centre-centre} (centre de l'astre au centre du satellite). L'altitude $h$ est la distance \textbf{sol-satellite}. La relation est :
\[
r = R_{\text{astre}} + h
\]
où $R_{\text{astre}}$ est le rayon de l'astre (ex. $R_{\text{Terre}} \approx \SI{6370}{km}$). \textbf{Utilisez toujours $r$ dans les formules $v = \sqrt{GM/r}$ et $T = 2\pi\sqrt{r^3/GM}$.}
\end{attention}

\courssub{Période de révolution et troisième loi de Kepler (cas circulaire)}

La période $T$ est le temps pour effectuer un tour complet ($2\pi r$) à la vitesse constante $v$ :
\[
T = \frac{2\pi r}{v} = 2\pi r \sqrt{\frac{r}{G M}} = 2\pi \sqrt{\frac{r^3}{G M}}
\]

\begin{propriete}[Période de révolution et 3e loi de Kepler (cas circulaire)]
\[
\boxed{T = 2\pi \sqrt{\frac{r^3}{G M}}} \qquad \text{et} \qquad \boxed{\frac{T^2}{r^3} = \frac{4\pi^2}{G M} = \text{constante pour un astre central donné}}
\]
C'est la \textbf{troisième loi de Kepler} démontrée dans le cas particulier du mouvement circulaire. Elle montre que le rapport $T^2/r^3$ ne dépend \textbf{que de la masse $M$ de l'astre central}, pas du satellite.
\end{propriete}

\begin{aretenir}[Analyse dimensionnelle (vérification rapide)]
$[G] = \SI{}{N.m^2.kg^{-2}} = \SI{}{m^3.kg^{-1}.s^{-2}}$.
$[GM/r] = \SI{}{m^3.kg^{-1}.s^{-2} \cdot kg \cdot m^{-1}} = \SI{}{m^2.s^{-2}} \Rightarrow [\sqrt{GM/r}] = \SI{}{m.s^{-1}}$ (vitesse). ✓
$[r^3/GM] = \SI{}{m^3 \cdot kg \cdot s^2 \cdot m^{-3} \cdot kg^{-1}} = \SI{}{s^2} \Rightarrow [2\pi\sqrt{r^3/GM}] = \SI{}{s}$ (période). ✓
\end{aretenir}

\courssub{Application : « Peser » un astre — Détermination de la masse $M$}

La relation $T^2/r^3 = 4\pi^2/(GM)$ permet de déterminer la masse $M$ d'un astre \textbf{sans s'y rendre}, en observant simplement le mouvement d'un de ses satellites (naturel ou artificiel). C'est l'une des plus belles applications de la physique newtonienne.

\begin{methode}[Déterminer la masse d'un astre central]
\begin{enumerate}
    \item Identifier un satellite (lune naturelle, sonde artificielle, planète) en orbite \emph{quasi circulaire} autour de l'astre.
    \item Mesurer (ou relever) le rayon orbital $r$ (distance centre-centre) et la période de révolution $T$.
    \item Appliquer la formule inversée :
    \[
    M = \frac{4\pi^2 r^3}{G T^2}
    \]
    \item Vérifier l'homogénéité des unités (SI impératif : $r$ en m, $T$ en s, $G$ en $\SI{}{N.m^2.kg^{-2}}$) et donner le résultat en kg avec le bon nombre de chiffres significatifs.
\end{enumerate}
\end{methode}

\begin{exemplebox}[Masse du Soleil à partir du mouvement de la Terre]
La Terre tourne autour du Soleil sur une orbite quasi circulaire de rayon $r_{\text{Terre-Soleil}} = \SI{1,496e11}{m}$ (1 UA) avec une période $T = \SI{365,25}{jours} = \SI{3,15576e7}{s}$.
\[
M_{\text{Soleil}} = \frac{4\pi^2 (\SI{1,496e11}{m})^3}{(\SI{6,67e-11}{N.m^2.kg^{-2}}) \times (\SI{3,15576e7}{s})^2}
\approx \frac{39,48 \times \SI{3,348e33}{m^3}}{\SI{6,67e-11}{} \times \SI{9,959e14}{s^2}}
\approx \SI{1,99e30}{kg}
\]
C'est la valeur de référence (1 masse solaire $\approx \SI{1,99e30}{kg}$). \textbf{Remarque :} On a « pesé » le Soleil sans y aller, juste en observant la Terre !
\end{exemplebox}

\begin{savaistu}[Exoplanètes et trous noirs : la même méthode à l'échelle de la galaxie]
C'est exactement cette méthode (loi de Kepler / vitesse orbitale) qui permet aux astrophysiciens de :
\begin{itemize}
    \item \textbf{Détecter des exoplanètes :} en mesurant la vitesse radiale de l'étoile (effet Doppler) due à la planète qui l'orbite, on en déduit $M_{\text{planète}}$ et $r$.
    \item \textbf{« Peser » le trou noir supermassif Sagittarius A* au centre de notre galaxie :} en suivant sur 20 ans les orbites d'étoiles (comme $S_2$) qui passent très près de lui ($r \sim \SI{120}{UA}$, $T \sim \SI{16}{ans}$), on a obtenu $M \approx 4 \times 10^6$ masses solaires (${\sim}\SI{8e36}{kg}$). Prix Nobel de physique 2020 (Reinhard Genzel et Andrea Ghez).
\end{itemize}
La physique du Bac vous donne les outils pour comprendre l'Univers !
\end{savaistu}

\courssub{Ordres de grandeur : de l'ISS à la Lune}

Le graphique ci-dessous illustre la décroissance de la vitesse orbitale $v \propto 1/\sqrt{r}$ et l'augmentation de la période $T \propto r^{3/2}$ avec l'altitude.

\begin{popfigure}
\begin{tikzpicture}
\begin{semilogxaxis}[
    width=\linewidth, height=8cm,
    xlabel={Rayon orbital $r$ (km)}, ylabel={Vitesse orbitale $v$ (km/s)},
    xmin=6370, xmax=400000,
    ymin=0, ymax=8.5,
    grid=major, grid style={popline},
    log ticks with fixed point,
    xtick={6370, 6780, 26560, 42164, 384400},
    xticklabels={$R_\oplus$ (sol), ISS, GPS, Geo, Lune},
    legend style={at={(0.98,0.98)}, anchor=north east, fill=popcream, draw=popmuted},
    domain=6370:400000, samples=200,
]
% Courbe théorique v = sqrt(GM/r) avec GM = 3.986e5 km^3/s^2
\addplot[popcurve, thick, color=popblue] {sqrt(max(0,3.986e5/x))};
\addlegendentry{$v = \sqrt{GM/r}$ (Théorique)}

% Points caractéristiques
\addplot[only marks, mark=*, mark size=3, poporange] coordinates {(6780, 7.67)};
\addlegendentry{ISS ($h\approx 400$ km)}

\addplot[only marks, mark=square*, mark size=3, popgreen] coordinates {(26560, 3.87)};
\addlegendentry{GPS ($h\approx 20\,200$ km)}

\addplot[only marks, mark=triangle*, mark size=3, poppurple] coordinates {(42164, 3.07)};
\addlegendentry{Géostationnaire ($h\approx 35\,786$ km)}

\addplot[only marks, mark=diamond*, mark size=3, popred] coordinates {(384400, 1.02)};
\addlegendentry{Lune ($r\approx 384\,400$ km)}
\end{semilogxaxis}
\end{tikzpicture}
\legende{Vitesse orbitale $v$ en fonction du rayon $r$ (échelle logarithmique) pour les satellites de la Terre. $GM_\oplus \approx \SI{3,986e5}{km^3.s^{-2}}$. La vitesse chute rapidement au début, puis plus lentement.}
\end{popfigure}

\begin{exoresolu}[Caractéristiques de l'orbite de l'ISS (Station Spatiale Internationale)]
\textbf{Données :} Altitude $h = \SI{400}{km}$, $R_\oplus = \SI{6370}{km}$, $M_\oplus = \SI{5,972e24}{kg}$, $G = \SI{6,67e-11}{N.m^2.kg^{-2}}$.
\textbf{Questions :}
\begin{enumerate}
    \item Calculer le rayon orbital $r$ en mètres.
    \item Déterminer la vitesse orbitale $v$ en $\SI{}{km.s^{-1}}$.
    \item Calculer la période de révolution $T$ en heures et minutes.
    \item Combien d'orbites l'ISS effectue-t-elle par jour ?
\end{enumerate}
\tcblower
\textbf{Solution détaillée :}
\begin{enumerate}
    \item \textbf{Rayon orbital $r$ :}
    $r = R_\oplus + h = \SI{6370}{km} + \SI{400}{km} = \SI{6770}{km} = \SI{6,77e6}{m}$.
    \item \textbf{Vitesse orbitale $v$ :}
    $v = \sqrt{\dfrac{G M_\oplus}{r}} = \sqrt{\dfrac{(\SI{6,67e-11}{}) \times (\SI{5,972e24}{})}{\SI{6,77e6}{}}}
    = \sqrt{\SI{5,89e7}{}} \approx \SI{7670}{m.s^{-1}} = \SI{7,67}{km.s^{-1}}$.
    \item \textbf{Période de révolution $T$ :}
    $T = \dfrac{2\pi r}{v} = \dfrac{2\pi \times \SI{6,77e6}{}}{\SI{7670}{}} \approx \SI{5540}{s} \approx \SI{1,54}{h} \approx \textbf{1 h 32 min}$.
    \item \textbf{Nombre d'orbites par jour :} $N = \dfrac{\SI{24}{h}}{\SI{1,54}{h}} \approx 15,6$ orbites/jour.
\end{enumerate}
\textit{L'ISS fait le tour de la Terre environ 16 fois par jour ! Les astronautes voient 16 levers et 16 couchers de Soleil par 24 h.}
\end{exoresolu}

\begin{exoresolu}[Masse de la Terre à partir du mouvement de la Lune]
\textbf{Données :} Distance Terre-Lune $r = \SI{3,844e8}{m}$, Période sidérale de la Lune $T = \SI{27,32}{jours} = \SI{2,360e6}{s}$.
\textbf{Question :} En déduire la masse de la Terre $M_\oplus$.
\tcblower
\textbf{Solution détaillée :}
Application de la 3e loi de Kepler (cas circulaire) :
\[
M_\oplus = \frac{4\pi^2 r^3}{G T^2}
= \frac{4\pi^2 (\SI{3,844e8}{})^3}{(\SI{6,67e-11}{}) \times (\SI{2,360e6}{})^2}
\]
Numérateur : $4\pi^2 \times \SI{5,68e25}{} \approx \SI{2,24e27}{}$.
Dénominateur : $\SI{6,67e-11}{} \times \SI{5,57e12}{} \approx \SI{3,72e2}{}$.
$M_\oplus \approx \dfrac{\SI{2,24e27}{}}{\SI{3,72e2}{}} \approx \SI{6,02e24}{kg}$.
\textit{Valeur acceptée : $\SI{5,97e24}{kg}$. L'écart (moins de 1\%) vient de l'excentricité de l'orbite lunaire ($e \approx 0,055$) et des perturbations solaires.}
\end{exoresolu}

\courssub{Synthèse des formules clés (à connaître par cœur)}

\begin{center}
\begin{tabularx}{\linewidth}{|l|X|X|}\hline
\textbf{Grandeur} & \textbf{Formule} & \textbf{Dépendance} \\ \hline
Vitesse orbitale & $v = \sqrt{\dfrac{G M}{r}}$ & $\propto M^{1/2}\, r^{-1/2}$ (indép. de $m$) \\ \hline
Période & $T = \dfrac{2\pi r}{v} = 2\pi \sqrt{\dfrac{r^3}{G M}}$ & $\propto M^{-1/2}\, r^{3/2}$ \\ \hline
3e loi de Kepler & $\dfrac{T^2}{r^3} = \dfrac{4\pi^2}{G M}$ & Constante pour un astre $M$ donné \\ \hline
Masse de l'astre & $M = \dfrac{4\pi^2 r^3}{G T^2}$ & Se déduit de $r$ et $T$ d'un satellite \\ \hline
\end{tabularx}
\end{center}

\begin{attention}[Conditions d'application — Rappel pour le Bac]
Ces formules sont valides \textbf{uniquement} si :
\begin{enumerate}
    \item L'orbite est \textbf{circulaire} (excentricité $e=0$). Pour les orbites elliptiques, $v$ et $r$ varient ; on utilise alors la loi des aires (2e loi) et l'énergie mécanique.
    \item La masse du satellite $m$ est \textbf{négligeable} devant $M$ ($m \ll M$). Sinon, $M$ est remplacé par $M+m$ (problème des deux corps).
    \item Le référentiel astrocentrique est \textbf{galiléen} (astre non accéléré, pas de rotation propre significative).
    \item Aucune autre force ne s'exerce (pas de traînée atmosphérique, pas de poussée moteur, pas de pression de radiation significative).
\end{enumerate}
\end{attention}

\begin{exercice}[Entraînement : Satellite GPS]
Un satellite du système GPS (type Navstar) orbite à une altitude $h = \SI{20200}{km}$.
Données : $R_\oplus = \SI{6370}{km}$, $GM_\oplus = \SI{3,986e14}{m^3.s^{-2}}$.
\begin{enumerate}
    \item Calculer son rayon orbital $r$ en mètres.
    \item Déterminer sa vitesse orbitale $v$ en $\SI{}{km.s^{-1}}$.
    \item Calculer sa période de révolution $T$ en heures et minutes.
    \item Vérifier que $T \approx \SI{12}{h}$ (demi-journée sidérale), caractéristique des orbites GPS.
\end{enumerate}
\end{exercice}

\begin{corrige}[Exercice : Satellite GPS]
\begin{enumerate}
    \item $r = R_\oplus + h = \SI{6370}{km} + \SI{20200}{km} = \SI{26570}{km} = \SI{2,657e7}{m}$.
    \item $v = \sqrt{\dfrac{GM_\oplus}{r}} = \sqrt{\dfrac{\SI{3,986e14}{}}{\SI{2,657e7}{}}} = \sqrt{\SI{1,500e7}{}} \approx \SI{3873}{m.s^{-1}} = \SI{3,87}{km.s^{-1}}$.
    \item $T = \dfrac{2\pi r}{v} = \dfrac{2\pi \times \SI{2,657e7}{}}{\SI{3873}{}} \approx \SI{43100}{s} \approx \SI{11,97}{h} \approx \textbf{11 h 58 min}$.
    \item C'est bien une période de demi-journée sidérale ($\approx \SI{11h58min}{}$). Les satellites GPS effectuent 2 orbites par jour sidéral, ce qui garantit une répétition de la géométrie constellation/sol chaque jour (pratique pour les récepteurs au sol, y compris au Cameroun).
\end{enumerate}
\end{corrige}

Nous avons maintenant tous les outils pour étudier le cas particulier et crucial du \textbf{satellite géostationnaire} (section 3), puis aborder l'énergie mécanique et la vitesse de libération (section 4).

\coursec{Le satellite géostationnaire}

Dans la section précédente, nous avons établi que tout satellite en orbite circulaire autour de la Terre possède une vitesse $v = \sqrt{\dfrac{G\,M_T}{r}}$ et une période $T = 2\pi\sqrt{\dfrac{r^3}{G\,M_T}}$ qui ne dépendent que du rayon $r$ de son orbite. Autrement dit, \textbf{à chaque altitude correspond une période bien déterminée}. L'ISS, à \SI{400}{km} d'altitude, boucle son tour en environ \SI{92}{min} ; la Lune, à \SI{384000}{km}, met environ 27,3 jours. Une question naturelle se pose alors : existe-t-il une altitude pour laquelle la période du satellite est \emph{exactement égale} à la période de rotation de la Terre sur elle-même ? Si oui, un tel satellite tournerait « en même temps » que la Terre et paraîtrait \textbf{immobile} vu depuis le sol. C'est précisément le satellite géostationnaire, pilier invisible de nos télécommunications quotidiennes : chaque fois que tu regardes une chaîne reçue par parabole à Yaoundé, le signal provient d'un satellite fixe dans le ciel, à près de \SI{36000}{km} d'altitude.

\courssub{Définition et conditions d'un satellite géostationnaire}

\begin{definition}[Satellite géostationnaire]
Un \cle{satellite géostationnaire} est un satellite qui paraît \textbf{immobile} par rapport à un point fixe de la surface de la Terre. Pour cela, trois conditions doivent être \emph{simultanément} réalisées :
\begin{enumerate}
\item \textbf{Condition de plan} : son orbite circulaire est située dans le \cle{plan équatorial} terrestre ;
\item \textbf{Condition de sens} : il tourne dans le \textbf{même sens} que la rotation propre de la Terre, c'est-à-dire d'ouest en est ;
\item \textbf{Condition de période} : sa période de révolution est égale à la période de rotation propre de la Terre :
\[ T = T_0 = \SI{86164}{s} \quad (\text{soit } 23~\text{h}~56~\text{min}~4~\text{s}). \]
\end{enumerate}
On dit alors que le satellite est \cle{géosynchrone} (même période que la Terre) \emph{et} équatorial : c'est la conjonction des trois conditions qui le rend géostationnaire.
\end{definition}

\begin{attention}[Jour sidéral ou jour solaire ?]
La période à utiliser est le \cle{jour sidéral} $T_0 = \SI{86164}{s}$, et \textbf{non} \SI{24}{h} ! Le jour solaire de \SI{24}{h} est le temps qui sépare deux passages du Soleil au méridien d'un lieu ; mais pendant que la Terre tourne sur elle-même, elle avance aussi sur son orbite autour du Soleil (environ $1^{\circ}$ par jour). Elle doit donc tourner d'un peu plus de $360^{\circ}$ pour « rattraper » le Soleil : cela prend environ 4 minutes de plus. Le jour sidéral, lui, est la durée d'un tour complet ($360^{\circ}$ exactement) par rapport aux étoiles lointaines, donc par rapport au référentiel géocentrique : c'est \textbf{lui} qui intervient dans la troisième loi de Kepler. Utiliser $T = \SI{24}{h} = \SI{86400}{s}$ fausse le calcul de l'altitude d'environ \SI{200}{km}.
\end{attention}

\begin{attention}[Un géostationnaire ne peut pas « survoler » Douala]
Le centre de l'orbite d'un satellite est nécessairement le \textbf{centre de la Terre} (la force de gravitation est centrale). Si l'orbite n'est pas dans le plan équatorial, le satellite oscille en latitude au cours de sa révolution : il passe alternativement au-dessus de l'hémisphère Nord puis de l'hémisphère Sud, et ne peut donc \emph{jamais} paraître fixe. Douala étant à la latitude $4^{\circ}$ Nord, aucun satellite géostationnaire ne peut se trouver à sa verticale : tous sont à la verticale d'un point de l'\textbf{équateur}. C'est pourquoi, à Douala comme à Yaoundé, les paraboles ne pointent pas vers le zénith mais vers un point du ciel situé au-dessus de l'équateur.
\end{attention}

\begin{popfigure}
\begin{center}
\begin{tikzpicture}[scale=1.0]
% Terre vue du pôle Nord
\fill[popblue!25] (0,0) circle (1.0);
\draw[popblue, thick] (0,0) circle (1.0);
\fill[popgreen!60] (-0.3,0.2) circle (0.25);
\fill[popgreen!60] (0.4,-0.3) circle (0.18);
\fill[popgreen!60] (0.15,0.45) circle (0.12);
\node[popblue!80!black] at (0,-0.05) {\small Terre};
% Orbite géostationnaire
\draw[poporange, very thick] (0,0) circle (3.2);
% Satellite
\fill[poporange] (3.2,0) circle (0.10);
\node[poporange!80!black, right] at (3.3,0) {\small satellite};
% Vecteur vitesse
\draw[-{Stealth}, poporange!80!black, very thick] (3.2,0) -- (3.2,1.3) node[midway, right] {\small $\vec{v}$};
% Force gravitationnelle
\draw[-{Stealth}, poppurple, very thick] (3.2,0) -- (1.7,0) node[midway, above] {\small $\vec{F}_{T/S}$};
% Rayon coté
\draw[-{Stealth}, popdark, thick] (0,-1.15) -- (3.2,-1.15);
\node[popdark, below] at (1.6,-1.2) {\small $r = R_T + h \approx \SI{42160}{km}$};
\draw[popdark, dashed] (0,-1.0) -- (0,-1.3);
\draw[popdark, dashed] (3.2,-0.1) -- (3.2,-1.3);
% Sens de rotation de la Terre
\draw[-{Stealth}, popblue!80!black, thick] (0.75,0.65) arc (40:140:0.99);
\node[popblue!80!black] at (-1.5,1.35) {\small rotation Terre};
% Sens de révolution du satellite
\draw[-{Stealth}, poporange!80!black, thick] (2.26,2.26) arc (45:110:3.2);
\node[poporange!80!black] at (-2.6,2.6) {\small révolution (vers l'est)};
% Point fixe au sol
\fill[popink] (1.0,0) circle (0.06);
\node[popink, above right] at (0.75,-0.35) {\small point $P$};
\draw[popink, dashed, thick] (0,0) -- (3.2,0);
\end{tikzpicture}
\end{center}
\legende{La Terre vue du pôle Nord. Le satellite géostationnaire évolue dans le plan équatorial, dans le même sens que la rotation terrestre, avec la même période $T_0$ : il reste en permanence à la verticale du même point $P$ de l'équateur. La force de gravitation $\vec{F}_{T/S}$ joue le rôle de force centripète.}
\end{popfigure}

\courssub{Calcul de l'altitude : une orbite unique}

Puisque la période $T$ d'un satellite circulaire ne dépend que du rayon $r$ de son orbite (troisième loi de Kepler), imposer $T = T_0$ fixe \textbf{une seule valeur possible} de $r$, donc une seule altitude. Tous les satellites géostationnaires du monde se trouvent ainsi sur un même cercle imaginaire autour de la Terre.

\begin{propriete}[Altitude du satellite géostationnaire — démonstration exigible]
Un satellite géostationnaire évolue sur une orbite circulaire de rayon
\[ r = \left(\dfrac{G\,M_T\,T_0^{\,2}}{4\pi^2}\right)^{1/3} \approx \SI{42160}{km}, \]
soit une altitude
\[ h = r - R_T \approx \SI{35786}{km} \approx \SI{36000}{km}. \]
\textbf{Démonstration.} Dans le référentiel géocentrique, supposé galiléen, le satellite $S$ de masse $m$ n'est soumis qu'à la force de gravitation terrestre $\vec{F} = -\dfrac{G\,M_T\,m}{r^2}\,\vec{u}_r$. La deuxième loi de Newton, projetée sur la direction normale (centripète) du repère de Frenet, s'écrit :
\[ \dfrac{G\,M_T\,m}{r^2} = m\,\dfrac{v^2}{r} \quad\Longrightarrow\quad v = \sqrt{\dfrac{G\,M_T}{r}}. \]
Le mouvement est uniforme, donc $v = \dfrac{2\pi r}{T_0}$. En égalisant :
\[ \dfrac{4\pi^2 r^2}{T_0^{\,2}} = \dfrac{G\,M_T}{r} \quad\Longrightarrow\quad \dfrac{T_0^{\,2}}{r^3} = \dfrac{4\pi^2}{G\,M_T} \quad\text{(3\textsuperscript{e} loi de Kepler)}. \]
On isole $r$ : $r^3 = \dfrac{G\,M_T\,T_0^{\,2}}{4\pi^2}$, d'où le résultat annoncé. \hfill$\blacksquare$
\end{propriete}

\begin{methode}[Déterminer l'altitude d'un satellite géostationnaire]
\begin{enumerate}
\item \textbf{Écrire la troisième loi de Kepler} pour le satellite autour de la Terre : $\dfrac{T^2}{r^3} = \dfrac{4\pi^2}{G\,M_T}$.
\item \textbf{Remplacer $T$ par le jour sidéral} $T_0 = \SI{86164}{s}$ (condition de géostationnarité).
\item \textbf{Isoler le rayon} : $r = \left(\dfrac{G\,M_T\,T_0^{\,2}}{4\pi^2}\right)^{1/3}$.
\item \textbf{Retrancher le rayon terrestre} : $h = r - R_T$, avec $R_T \approx \SI{6370}{km}$.
\item \textbf{Vérifier la cohérence} : on doit trouver $h \approx \SI{36000}{km}$ ; toute valeur très différente signale une erreur (souvent $T = \SI{24}{h}$ ou $r$ confondu avec $h$).
\end{enumerate}
\end{methode}

\begin{exoresolu}[Altitude et vitesse d'un satellite géostationnaire]
\textbf{Énoncé.} On donne $G = \num{6,67e-11}~\text{N}\cdot\text{m}^2\cdot\text{kg}^{-2}$, $M_T = \num{5,97e24}{kg}$, $R_T = \SI{6370}{km}$ et $T_0 = \SI{86164}{s}$. Déterminer le rayon $r$ de l'orbite géostationnaire, l'altitude $h$ du satellite, puis sa vitesse $v$ sur orbite.
\tcblower
\textbf{Solution.} Appliquons la troisième loi de Kepler avec $T = T_0$ :
\begin{align*}
r^3 &= \dfrac{G\,M_T\,T_0^{\,2}}{4\pi^2} = \dfrac{\num{6,67e-11}\times\num{5,97e24}\times(\num{86164})^2}{4\pi^2} \\
&= \dfrac{\num{6,67e-11}\times\num{5,97e24}\times\num{7,424e9}}{39,48} \approx \num{7,49e22}~\text{m}^3.
\end{align*}
D'où $r \approx (\num{7,49e22})^{1/3} \approx \num{4,216e7}{m}$, soit $\boxed{r \approx \SI{42160}{km}}$.
L'altitude est alors :
\[ h = r - R_T = 42160 - 6370 \approx \boxed{\SI{35790}{km}} \approx \SI{36000}{km}. \]
La vitesse sur orbite vaut :
\[ v = \sqrt{\dfrac{G\,M_T}{r}} = \sqrt{\dfrac{\num{6,67e-11}\times\num{5,97e24}}{\num{4,216e7}}} \approx \sqrt{\num{9,44e6}} \approx \boxed{\num{3,07e3}{m/s}}, \]
soit environ $\SI{3,1}{km/s}$, c'est-à-dire près de $\SI{11000}{km/h}$. Remarquons que c'est nettement moins que la vitesse de l'ISS ($\SI{7,7}{km/s}$) : plus l'orbite est haute, plus le satellite est lent, conformément à $v = \sqrt{G M_T/r}$.
\end{exoresolu}

\begin{attention}[Pièges de rédaction au Bac]
\begin{itemize}
\item Ne jamais confondre le \textbf{rayon de l'orbite} $r$ (mesuré depuis le \emph{centre} de la Terre) et l'\textbf{altitude} $h$ (mesurée depuis la surface) : $r = R_T + h$. La loi de Kepler porte sur $r$, jamais sur $h$.
\item Convertir systématiquement les grandeurs en unités SI : $T_0$ en secondes, $r$ en mètres.
\item Un satellite géostationnaire paraît immobile mais il n'est \textbf{pas au repos} : il file à $\SI{3,1}{km/s}$ dans le référentiel géocentrique. « Immobile » signifie seulement : immobile \emph{par rapport au sol terrestre}.
\end{itemize}
\end{attention}

\courssub{Applications et place parmi les orbites terrestres}

L'immobilité apparente du satellite géostationnaire est une aubaine technologique : une antenne de réception n'a besoin d'être pointée qu'\textbf{une seule fois}, puis reste fixe.

\begin{itemize}
\item \textbf{Télécommunications et télévision} : les signaux des bouquets satellite (Canal+, DStv...) sont relayés par des satellites géostationnaires ; trois satellites bien répartis suffisent à couvrir quasiment toute la planète (hors pôles).
\item \textbf{Météorologie} : les satellites Meteosat photographient en continu le même disque terrestre, permettant de suivre l'évolution des nuages et des cyclones.
\item \textbf{Internet et téléphonie} dans les zones rurales non couvertes par la fibre ou le réseau mobile.
\item \textbf{Nuance importante} : les satellites du système \textbf{GPS} ne sont \emph{pas} géostationnaires ! Ils évoluent à environ $\SI{20200}{km}$ d'altitude avec une période d'environ 12 h, afin qu'au moins quatre d'entre eux soient visibles de tout point du globe à tout instant.
\end{itemize}

\begin{savaistu}[L'orbite de Clarke]
L'orbite géostationnaire est aussi appelée \cle{orbite de Clarke}. En 1945, bien avant le premier satellite artificiel (Spoutnik, 1957), l'écrivain et ingénieur britannique Arthur C. Clarke — futur auteur de \emph{2001, l'Odyssée de l'espace} — publia un article visionnaire intitulé \emph{Extra-Terrestrial Relays}, dans lequel il calcula qu'un satellite placé à environ $\SI{36000}{km}$ au-dessus de l'équateur paraîtrait fixe et pourrait servir de relais mondial de télécommunications. Dix-huit ans plus tard, en 1963, le satellite Syncom 2 réalisa exactement cette prédiction. Aujourd'hui, des centaines de satellites occupent cette « ceinture » unique, et les positions angulaires (les « slots » orbitaux) y sont attribuées par l'Union internationale des télécommunications comme une ressource rare.
\end{savaistu}

\begin{popfigure}
\begin{center}
\begin{tikzpicture}[scale=1.0]
% Terre
\fill[popblue!30] (0,0) circle (0.637);
\draw[popblue, thick] (0,0) circle (0.637);
\node[popblue!80!black] at (0,0) {\scriptsize Terre};
% ISS : r = 6770 km -> 0.677
\draw[popgreen, thick] (0,0) circle (0.677);
\fill[popgreen!80!black] (0.677,0) circle (0.05);
\node[popgreen!60!black, right] at (0.72,0.12) {\scriptsize ISS ($h \approx 400$ km)};
% GPS : r = 26570 km -> 2.657
\draw[poppurple, thick] (0,0) circle (2.657);
\fill[poppurple] (-2.657,0) circle (0.05);
\node[poppurple!80!black, left] at (-2.7,0.15) {\scriptsize GPS ($h \approx 20200$ km)};
% Géostationnaire : r = 42160 km -> 4.216
\draw[poporange, very thick] (0,0) circle (4.216);
\fill[poporange] (0,4.216) circle (0.06);
\node[poporange!80!black, above] at (0,4.28) {\scriptsize géostationnaire ($h \approx 35786$ km)};
% Lune : r = 384400 km -> 38.44 : on représente la direction et on note hors échelle
\draw[popmuted, dashed, thick] (0,0) -- (4.6,-2.6);
\fill[popmuted] (4.6,-2.6) circle (0.12);
\node[popmuted!80!black, below right] at (4.0,-2.75) {\scriptsize Lune (non à l'échelle : $d \approx 384400$ km, soit 9 fois plus loin)};
% Échelle
\draw[popdark, thick] (-4.6,-3.4) -- (-3.6,-3.4);
\draw[popdark, thick] (-4.6,-3.5) -- (-4.6,-3.3);
\draw[popdark, thick] (-3.6,-3.5) -- (-3.6,-3.3);
\node[popdark, below] at (-4.1,-3.45) {\scriptsize 10000 km};
\end{tikzpicture}
\end{center}
\legende{Comparaison à l'échelle des principales orbites terrestres (rayons depuis le centre de la Terre) : ISS ($r \approx \SI{6770}{km}$), GPS ($r \approx \SI{26570}{km}$), géostationnaire ($r \approx \SI{42160}{km}$). La Lune, à $\SI{384400}{km}$, sortirait largement du cadre : sa distance est indiquée hors échelle.}
\end{popfigure}

\begin{exoresolu}[Pourquoi toutes les paraboles de Yaoundé pointent-elles dans la même direction ?]
\textbf{Énoncé.} Un technicien installe des paraboles à Yaoundé (latitude $\varphi \approx 3,9^{\circ}$ Nord) pour recevoir un bouquet diffusé par un satellite géostationnaire positionné à la longitude de Yaoundé. (a) Expliquer pourquoi toutes les paraboles de la ville pointent dans la même direction. (b) Vers quel point cardinal sont-elles orientées ? (c) Calculer la durée $\Delta t$ mise par le signal pour faire l'aller-retour parabole--satellite--parabole (on prendra $h \approx \SI{35786}{km}$ et $c = \num{3,00e8}{m/s}$).
\tcblower
\textbf{Solution.} (a) Le satellite est géostationnaire : il paraît \textbf{immobile} dans le ciel, à la verticale d'un point fixe de l'équateur. La direction dans laquelle il faut pointer l'antenne est donc la même pour tous les habitants de la ville, et elle ne change jamais : une fois la parabole réglée, on n'y touche plus.
(b) Le satellite se trouve à la verticale de l'équateur, donc au \textbf{sud} de Yaoundé (qui est dans l'hémisphère Nord) : les paraboles pointent vers le sud, inclinées vers le haut d'un angle voisin de $90^{\circ} - 3,9^{\circ} \approx 86^{\circ}$ par rapport à l'horizontale, c'est-à-dire presque à la verticale, légèrement penchées vers le sud.
(c) Le signal parcourt deux fois la distance $h$ (montée et descente) :
\[ \Delta t = \dfrac{2h}{c} = \dfrac{2\times\num{3,5786e7}}{\num{3,00e8}} \approx \boxed{\num{0,24}{s}}. \]
C'est ce léger retard, perceptible dans les directs télévisés par satellite, qui trahit l'immense distance parcourue par l'onde.
\end{exoresolu}

\begin{aretenir}[L'essentiel sur le satellite géostationnaire]
\begin{itemize}
\item Un satellite géostationnaire paraît \textbf{fixe} vu du sol : orbite \textbf{équatoriale}, sens de révolution \textbf{identique} à celui de la Terre (vers l'est), période égale au \textbf{jour sidéral} $T_0 = \SI{86164}{s}$.
\item La troisième loi de Kepler impose un rayon \textbf{unique} : $r \approx \SI{42160}{km}$, soit une altitude $h \approx \SI{35786}{km} \approx \SI{36000}{km}$, et une vitesse $v \approx \SI{3,1}{km/s}$.
\item Attention aux deux pièges classiques : utiliser \SI{24}{h} au lieu du jour sidéral, et confondre $r$ (depuis le centre) avec $h$ (depuis la surface).
\item Applications : télévision par satellite, télécommunications, météorologie — mais pas le GPS, dont les satellites évoluent plus bas.
\end{itemize}
\end{aretenir}

\coursec{Énergie mécanique d'un satellite (complément Bac)}

Dans la section précédente, nous avons étudié le satellite géostationnaire et découvert que son altitude n'est pas choisie au hasard : elle résulte d'une condition précise sur sa période de révolution. Mais une question pratique demeure : comment placer un satellite sur cette orbite ? Quelle énergie faut-il lui fournir ? Et surtout, quelle vitesse faut-il communiquer à une fusée pour qu'elle quitte définitivement la Terre, comme les sondes spatiales ? Ces questions, qui passionnent les ingénieurs du spatial, trouvent leur réponse dans un outil puissant que tu connais déjà : l'\cle{énergie mécanique}. Cette section, bien que présentée comme complément, est un classique des sujets de Baccalauréat C.

\courssub{Énergie potentielle de gravitation}

\textbf{Pourquoi une nouvelle expression de l'énergie potentielle ?}\par

Au chapitre de la mécanique, tu as appris que l'énergie potentielle de pesanteur s'écrit $E_p = mgz$ pour un objet de masse $m$ situé à l'altitude $z$. Mais cette expression repose sur une hypothèse forte : le champ de pesanteur $g$ est \emph{uniforme}. Or, pour un satellite qui évolue à des centaines ou des milliers de kilomètres d'altitude, le champ de gravitation varie sensiblement : $g(r) = \dfrac{GM}{r^2}$ dépend de la distance $r$ au centre de l'astre. Il faut donc une expression plus générale.

\textbf{Établissement de l'expression.} L'énergie potentielle associée à une force conservative est définie à partir du travail de cette force. Considérons un satellite de masse $m$ attiré par un astre de masse $M$ (dont le centre $O$ est l'origine des positions). La force de gravitation, radiale et attractive, s'écrit :
\[ \vec{F} = -\,\dfrac{GMm}{r^2}\,\vec{u}_r \]
où $\vec{u}_r$ est le vecteur unitaire dirigé de $O$ vers le satellite. Le travail élémentaire de cette force lors d'un déplacement radial $dr$ est $\delta W = \vec{F}\cdot d\vec{r} = -\dfrac{GMm}{r^2}\,dr$. Par définition de l'énergie potentielle, $\delta W = -dE_p$, donc :
\[ dE_p = \dfrac{GMm}{r^2}\,dr \quad\Longrightarrow\quad E_p(r) = -\dfrac{GMm}{r} + \text{constante} \]
La constante se fixe en choisissant une \cle{origine des énergies potentielles}. Par convention universelle, on choisit $E_p = 0$ lorsque le satellite est \emph{infiniment éloigné} de l'astre ($r \to \infty$), c'est-à-dire lorsqu'il n'interagit plus avec lui. La constante est alors nulle.

\begin{definition}[Énergie potentielle de gravitation]
L'énergie potentielle de gravitation d'un système \{astre de masse $M$ — satellite de masse $m$\} dont les centres sont distants de $r$ est :
\[ \boxed{E_p = -\,\dfrac{G\,M\,m}{r}} \]
avec $G = \num{6,67}\times 10^{-11}\ \mathrm{N\,m^2\,kg^{-2}}$ la constante de gravitation universelle. L'origine de $E_p$ est prise à l'infini : $E_p \to 0$ quand $r \to \infty$.
\end{definition}

\textbf{Vérification dimensionnelle.} $[G] = \mathrm{m^3\,kg^{-1}\,s^{-2}}$, donc $\left[\dfrac{GMm}{r}\right] = \dfrac{\mathrm{m^3\,kg^{-1}\,s^{-2}}\times \mathrm{kg}\times \mathrm{kg}}{\mathrm{m}} = \mathrm{kg\,m^2\,s^{-2}} = \mathrm{J}$. L'expression est bien homogène à une énergie.

\textbf{Pourquoi ce signe négatif ?} C'est une question que se posent tous les élèves, et elle mérite une réponse claire. L'énergie potentielle mesure le « réservoir » d'énergie lié à l'interaction. Ici, l'interaction est \emph{attractive} : pour arracher le satellite à l'astre et l'emmener jusqu'à l'infini (où $E_p = 0$), un opérateur extérieur doit \emph{fournir} du travail. Partant de $E_p = 0$ à l'infini, l'énergie potentielle ne peut donc que \emph{diminuer} quand on rapproche le satellite : elle devient négative. Le signe négatif traduit le fait que le satellite est dans un \cle{puits de gravitation} : il est \textbf{lié} à l'astre, exactement comme une bille au fond d'une cuvette est liée à la cuvette.

\begin{attention}[Origine de l'énergie potentielle : le piège classique]
Deux erreurs reviennent sans cesse dans les copies :
\begin{itemize}
\item \textbf{Oublier le signe négatif} de $E_p$ (et donc de $E_m$). Une énergie mécanique positive pour un satellite en orbite fermée est le signe certain d'une erreur de signe.
\item \textbf{Prendre l'origine de $E_p$ à la surface de l'astre}, par habitude de $E_p = mgz$. C'est faux ici : l'origine est à l'\emph{infini}. À la surface de l'astre ($r = R$), l'énergie potentielle vaut $E_p(R) = -\dfrac{GMm}{R}$, qui est sa valeur \emph{la plus négative} (le fond du puits).
\end{itemize}
Retiens : $r$ est toujours mesuré \emph{depuis le centre de l'astre}, jamais depuis sa surface. Si l'énoncé donne l'altitude $h$, alors $r = R + h$.
\end{attention}

\courssub{Énergie cinétique et énergie mécanique sur orbite circulaire}

\textbf{Énergie cinétique du satellite}\par

Nous avons établi dans la section 2 que, pour un satellite en mouvement circulaire uniforme de rayon $r$, la vitesse est :
\[ v = \sqrt{\dfrac{GM}{r}} \quad\Longrightarrow\quad v^2 = \dfrac{GM}{r} \]
L'énergie cinétique s'écrit donc :
\[ E_c = \dfrac{1}{2}mv^2 = \dfrac{1}{2}m\,\dfrac{GM}{r} \]
\[ \boxed{E_c = \dfrac{GMm}{2r}} \]
Remarque remarquable : en valeur absolue, $E_c$ est exactement la \emph{moitié} de $E_p$ : $E_c = -\dfrac{E_p}{2}$.

\textbf{Énergie mécanique : démonstration complète}\par

\begin{propriete}[Énergie mécanique d'un satellite en orbite circulaire — démonstration exigible]
Pour un satellite de masse $m$ en orbite circulaire de rayon $r$ autour d'un astre de masse $M$, l'énergie mécanique est :
\[ \boxed{E_m = E_c + E_p = -\,\dfrac{GMm}{2r}} \]
\textbf{Démonstration.} L'énergie mécanique est la somme des énergies cinétique et potentielle :
\begin{align*}
E_m &= E_c + E_p \\
&= \dfrac{GMm}{2r} + \left(-\dfrac{GMm}{r}\right) \\
&= \dfrac{GMm}{2r} - \dfrac{2\,GMm}{2r} \\
&= -\,\dfrac{GMm}{2r} \qquad \blacksquare
\end{align*}
Cette démonstration, courte mais fondamentale, est régulièrement demandée au Baccalauréat : sache la rédiger proprement en citant $v^2 = \dfrac{GM}{r}$ (issue de la deuxième loi de Newton appliquée au satellite).
\end{propriete}

\begin{aretenir}[Les trois énergies en un coup d'œil]
Pour une orbite circulaire de rayon $r$ :
\[ E_p = -\dfrac{GMm}{r} \qquad E_c = +\dfrac{GMm}{2r} \qquad E_m = -\dfrac{GMm}{2r} \]
Relations remarquables : $E_m = -E_c$ et $E_m = \dfrac{E_p}{2}$.
\end{aretenir}

\textbf{Commentaires physiques essentiels :}
\begin{itemize}
\item $E_m < 0$ : l'énergie mécanique négative est la signature d'un \cle{état lié}. Le satellite ne peut pas s'échapper seul de l'attraction de l'astre ; il est prisonnier du puits de gravitation, comme un électron autour du noyau dans l'atome.
\item Plus $r$ est grand, plus $E_m$ \emph{augmente} (elle se rapproche de $0$ par valeurs négatives). Passer sur une orbite plus haute demande donc de \emph{fournir} de l'énergie au satellite : c'est le rôle des moteurs de la fusée ou des propulseurs d'apogée.
\item Paradoxe apparent : sur une orbite plus haute, le satellite va \emph{moins vite} ($v = \sqrt{GM/r}$ diminue) mais possède \emph{plus} d'énergie mécanique. L'augmentation de $E_p$ l'emporte sur la diminution de $E_c$.
\end{itemize}

\begin{popfigure}
\begin{center}
\begin{tikzpicture}
\begin{axis}[
  width=0.95\linewidth, height=7.2cm,
  xlabel={$r$ (en unités de $R_T$)}, ylabel={Énergie (unités arbitraires)},
  xmin=1, xmax=8, ymin=-1.15, ymax=0.75,
  axis lines=left, grid=major, grid style={popline!40},
  legend style={at={(0.97,0.97)}, anchor=north east, font=\small},
  samples=100, domain=1:8]
  \addplot[popcurve, popblue, very thick] {-1/x};
  \addlegendentry{$E_p(r) = -\dfrac{GMm}{r}$}
  \addplot[popcurve, popgreen, very thick] {0.5/x};
  \addlegendentry{$E_c(r) = +\dfrac{GMm}{2r}$}
  \addplot[popcurve, poporange, very thick] {-0.5/x};
  \addlegendentry{$E_m(r) = -\dfrac{GMm}{2r}$}
  \addplot[popline, dashed] {0};
\end{axis}
\end{tikzpicture}
\end{center}
\legende{Évolutions de $E_p$, $E_c$ et $E_m$ en fonction du rayon $r$ de l'orbite. On vérifie graphiquement que $E_m = E_c + E_p = -E_c$ : la courbe orange est la symétrique de la courbe verte par rapport à l'axe des abscisses. Quand $r$ augmente, les trois énergies tendent vers $0$.}
\end{popfigure}

\courssub{La vitesse de libération}

\textbf{Position du problème}\par

Lance une balle verticalement : elle retombe. Lance-la plus fort : elle retombe plus loin, plus tard. Newton lui-même imaginait un canon au sommet d'une montagne tirant des boulets de plus en plus vite : le boulet finit par « retomber au-delà de l'horizon » et se mettre en orbite. Existe-t-il une vitesse à partir de laquelle le projectile \emph{ne retombe plus du tout} et s'éloigne indéfiniment de l'astre ? Cette vitesse existe : c'est la \cle{vitesse de libération}.

\begin{definition}[Vitesse de libération — démonstration par conservation de l'énergie]
La \cle{vitesse de libération} $v_{\text{lib}}$ est la vitesse \emph{minimale} qu'il faut communiquer à un objet de masse $m$, au départ de la \emph{surface} d'un astre de masse $M$ et de rayon $R$, pour qu'il échappe définitivement à l'attraction de cet astre (sans propulsion ultérieure) :
\[ \boxed{v_{\text{lib}} = \sqrt{\dfrac{2GM}{R}}} \]
\textbf{Démonstration.} On néglige les frottements de l'atmosphère ; l'énergie mécanique du projectile se conserve. Au départ, à la surface ($r = R$), lancé à la vitesse $v$ :
\[ E_m^{\text{départ}} = \dfrac{1}{2}mv^2 - \dfrac{GMm}{R} \]
« Échapper à l'attraction » signifie atteindre l'infini ($r \to \infty$, donc $E_p \to 0$). Le cas limite — la vitesse \emph{minimale} — correspond à un projectile qui arrive à l'infini avec une vitesse nulle ($E_c \to 0$). Son énergie mécanique à l'infini est alors :
\[ E_m^{\infty} = 0 + 0 = 0 \]
La conservation de l'énergie mécanique s'écrit :
\begin{align*}
\dfrac{1}{2}mv_{\text{lib}}^2 - \dfrac{GMm}{R} &= 0 \\
\dfrac{1}{2}v_{\text{lib}}^2 &= \dfrac{GM}{R} \\
v_{\text{lib}} &= \sqrt{\dfrac{2GM}{R}} \qquad \blacksquare
\end{align*}
La masse $m$ du projectile se simplifie : la vitesse de libération \emph{ne dépend que de l'astre} ($M$ et $R$), pas de l'objet lancé.
\end{definition}

\begin{popfigure}
\begin{center}
\begin{tikzpicture}[scale=1.0]
  % Terre
  \shade[ball color=popblue!60] (0,0) circle (1.1);
  \node[white, font=\small\bfseries] at (0,0) {Terre};
  % Pas de tir
  \fill[popdark] (0,1.1) circle (0.07);
  \node[above, font=\small] at (0,1.15) {pas de tir};
  % Trajectoire v < v1 : retombe
  \draw[popgreen, very thick, -{Stealth[length=3mm]}] (0,1.1) .. controls (0.9,1.9) and (2.1,1.5) .. (2.6,0.35);
  \node[popgreen, font=\small, align=center] at (3.6,1.35) {$v < \sqrt{\dfrac{GM}{R}}$ :\\ le projectile retombe};
  % Trajectoire orbitale
  \draw[poporange, very thick] (0,0) circle (2.1);
  \draw[poporange, -{Stealth[length=3mm]}] (2.1,0) arc (0:35:2.1);
  \node[poporange, font=\small, align=center] at (-3.9,-1.7) {$v = \sqrt{\dfrac{GM}{R}}$ :\\ orbite circulaire};
  % Trajectoire d'échappement
  \draw[poppurple, very thick, -{Stealth[length=3mm]}] (0,1.1) .. controls (0.35,2.6) and (1.5,3.6) .. (3.6,4.3);
  \node[poppurple, font=\small, align=center] at (4.1,3.3) {$v \geq v_{\text{lib}}$ :\\ échappement\\ (trajectoire ouverte)};
  % Vecteur vitesse initiale
  \draw[-{Stealth[length=3mm]}, very thick, popdark] (0,1.1) -- (1.1,1.1) node[midway, above, font=\small] {$\vec{v}$};
\end{tikzpicture}
\end{center}
\legende{L'expérience de pensée du « canon de Newton ». Selon la vitesse initiale communiquée au projectile, la trajectoire est un arc qui retombe au sol (vert), une orbite fermée (orange) ou une trajectoire ouverte d'échappement (violet) dès que $v \geq v_{\text{lib}} = \sqrt{2GM/R}$.}
\end{popfigure}

\textbf{Comparaison avec la vitesse de satellisation}\par

Il est très instructif de comparer la vitesse de libération à la vitesse de satellisation « au ras du sol », c'est-à-dire la vitesse d'une orbite circulaire de rayon $r = R$ :
\[ v_{\text{sat}} = \sqrt{\dfrac{GM}{R}} \qquad\text{et}\qquad v_{\text{lib}} = \sqrt{\dfrac{2GM}{R}} = \sqrt{2}\;v_{\text{sat}} \]
\[ \boxed{v_{\text{lib}} = \sqrt{2}\;v_{\text{sat}}} \]
Le facteur $\sqrt{2} \approx \num{1,41}$ entre les deux est une conséquence directe du facteur $1/2$ dans l'expression de $E_c$. Pour la Terre : $v_{\text{sat}} \approx \num{7,9}\ \mathrm{km/s}$ et $v_{\text{lib}} \approx \num{11,2}\ \mathrm{km/s}$.

\begin{exemplebox}[Vitesses de libération de la Terre et de la Lune]
\textbf{Données :} $G = \num{6,67}\times 10^{-11}\ \mathrm{N\,m^2\,kg^{-2}}$ ; Terre : $M_T = \num{5,97}\times 10^{24}\ \mathrm{kg}$, $R_T = \num{6,37}\times 10^{6}\ \mathrm{m}$ ; Lune : $M_L = \num{7,35}\times 10^{22}\ \mathrm{kg}$, $R_L = \num{1,74}\times 10^{6}\ \mathrm{m}$.

\textbf{Pour la Terre :}
\[ v_{\text{lib}}^{T} = \sqrt{\dfrac{2\times \num{6,67}\times 10^{-11}\times \num{5,97}\times 10^{24}}{\num{6,37}\times 10^{6}}} = \sqrt{\num{1,25}\times 10^{8}} \approx \num{1,12}\times 10^{4}\ \mathrm{m/s} \]
soit $v_{\text{lib}}^{T} \approx \num{11,2}\ \mathrm{km/s}$, environ 33 fois la vitesse du son dans l'air !

\textbf{Pour la Lune :}
\[ v_{\text{lib}}^{L} = \sqrt{\dfrac{2\times \num{6,67}\times 10^{-11}\times \num{7,35}\times 10^{22}}{\num{1,74}\times 10^{6}}} = \sqrt{\num{5,63}\times 10^{6}} \approx \num{2,37}\times 10^{3}\ \mathrm{m/s} \]
soit $v_{\text{lib}}^{L} \approx \num{2,4}\ \mathrm{km/s}$.

\textbf{Conséquence spectaculaire :} les molécules d'un gaz, agitées thermiquement, ont des vitesses de l'ordre de quelques centaines de mètres par seconde, voire davantage pour les plus légères (dihydrogène, hélium). Sur la Lune, $v_{\text{lib}}$ est si faible qu'une fraction non négligeable des molécules dépasse cette vitesse et s'échappe dans l'espace : c'est l'une des raisons pour lesquelles \emph{la Lune n'a pratiquement pas d'atmosphère}. La Terre, avec ses $\num{11,2}\ \mathrm{km/s}$, retient bien mieux la sienne.
\end{exemplebox}

\courssub{Coût énergétique d'un changement d'orbite}

Puisque $E_m = -\dfrac{GMm}{2r}$ augmente avec $r$, faire passer un satellite d'une orbite basse ($r_1$) à une orbite plus haute ($r_2 > r_1$) exige de lui fournir l'énergie :
\[ \Delta E_m = E_m(r_2) - E_m(r_1) = \dfrac{GMm}{2}\left(\dfrac{1}{r_1} - \dfrac{1}{r_2}\right) > 0 \]
C'est précisément l'énergie que doivent fournir les moteurs. C'est pourquoi placer un satellite géostationnaire (altitude $\approx \num{35800}\ \mathrm{km}$) coûte beaucoup plus cher qu'un satellite d'observation en orbite basse ($\approx 600\ \mathrm{km}$) : la masse de propergol embarquée croît considérablement, ce qui alourdit le lanceur.

\begin{exoresolu}[Énergie nécessaire pour élever un satellite]
Un satellite de masse $m = 800\ \mathrm{kg}$ est en orbite circulaire basse à l'altitude $h_1 = 400\ \mathrm{km}$. On veut le transférer sur une orbite circulaire à l'altitude $h_2 = 800\ \mathrm{km}$. Calculer l'énergie que doivent fournir les propulseurs. On donne $GM_T = \num{3,98}\times 10^{14}\ \mathrm{m^3\,s^{-2}}$ et $R_T = \num{6,37}\times 10^{6}\ \mathrm{m}$.
\tcblower
\textbf{Solution.} Les rayons des orbites sont :
\[ r_1 = R_T + h_1 = \num{6,37}\times 10^{6} + \num{0,40}\times 10^{6} = \num{6,77}\times 10^{6}\ \mathrm{m} \]
\[ r_2 = R_T + h_2 = \num{6,37}\times 10^{6} + \num{0,80}\times 10^{6} = \num{7,17}\times 10^{6}\ \mathrm{m} \]
L'énergie à fournir est la variation d'énergie mécanique :
\begin{align*}
\Delta E_m &= \dfrac{GM_T\,m}{2}\left(\dfrac{1}{r_1} - \dfrac{1}{r_2}\right) \\
&= \dfrac{\num{3,98}\times 10^{14}\times 800}{2}\left(\dfrac{1}{\num{6,77}\times 10^{6}} - \dfrac{1}{\num{7,17}\times 10^{6}}\right) \\
&= \num{1,592}\times 10^{17}\times\left(\num{1,477}\times 10^{-7} - \num{1,395}\times 10^{-7}\right) \\
&= \num{1,592}\times 10^{17}\times \num{8,24}\times 10^{-9} \\
&\approx \num{1,31}\times 10^{9}\ \mathrm{J}
\end{align*}
Il faut fournir environ $\num{1,3}\ \mathrm{GJ}$, soit l'équivalent de l'énergie électrique consommée par un quartier de Yaoundé pendant plusieurs heures. Vérification de cohérence : $\Delta E_m > 0$, ce qui est normal puisque l'orbite haute possède une énergie mécanique plus grande (moins négative).
\end{exoresolu}

\begin{attention}[Deux pièges supplémentaires à éviter]
\begin{itemize}
\item Ne confonds pas $v_{\text{lib}} = \sqrt{\dfrac{2GM}{R}}$ (quitter l'astre, départ de la \emph{surface}) et $v = \sqrt{\dfrac{GM}{r}}$ (rester en \emph{orbite} de rayon $r$). Le facteur $2$ sous le radical fait toute la différence.
\item Dans le calcul de $\Delta E_m$, n'oublie pas que $r = R + h$ : utiliser directement l'altitude $h$ à la place de $r$ est une erreur fréquente qui fausse tout le résultat.
\end{itemize}
\end{attention}

\begin{savaistu}[Un complément... qui tombe au Bac !]
Bien que présenté comme « complément » dans le programme officiel, le bilan énergétique d'un satellite apparaît régulièrement dans les sujets de Baccalauréat C, souvent sous forme de question bonus en fin d'exercice de mécanique : « Montrer que l'énergie mécanique du satellite s'écrit $E_m = -\dfrac{GMm}{2r}$ » ou « Calculer la vitesse minimale de lancement ». Maîtriser cette section, c'est donc s'offrir des points précieux. Sache aussi que la vitesse de libération terrestre de $\num{11,2}\ \mathrm{km/s}$ explique pourquoi les lanceurs comme Ariane sont des fusées à plusieurs étages : aucune structure ne supporterait de donner toute cette vitesse en une seule impulsion, et larguer les réservoirs vides en vol allège progressivement l'ensemble. Les sondes Voyager, lancées en 1977 à plus de $\num{11,2}\ \mathrm{km/s}$, ont quitté le Système solaire et continuent d'envoyer des signaux... captés par des antennes géantes, cousines de celles des télécommunications qui relient le Cameroun au reste du monde.
\end{savaistu}

\begin{exercice}[À toi de jouer]
La planète Mars a une masse $M = \num{6,42}\times 10^{23}\ \mathrm{kg}$ et un rayon $R = \num{3,39}\times 10^{6}\ \mathrm{m}$.
\begin{enumerate}
\item Calculer la vitesse de libération à la surface de Mars.
\item Un satellite de $500\ \mathrm{kg}$ est en orbite circulaire de rayon $r = \num{9,4}\times 10^{6}\ \mathrm{m}$ autour de Mars. Calculer ses énergies cinétique, potentielle et mécanique.
\item Le signe de $E_m$ était-il prévisible ? Justifier.
\end{enumerate}
\end{exercice}

\begin{corrige}[Exercice — Vitesse de libération et énergies autour de Mars]
\textbf{1.} $v_{\text{lib}} = \sqrt{\dfrac{2GM}{R}} = \sqrt{\dfrac{2\times \num{6,67}\times 10^{-11}\times \num{6,42}\times 10^{23}}{\num{3,39}\times 10^{6}}} = \sqrt{\num{2,53}\times 10^{7}} \approx \num{5,0}\times 10^{3}\ \mathrm{m/s}$, soit $\num{5,0}\ \mathrm{km/s}$.

\textbf{2.} $E_c = \dfrac{GMm}{2r} = \dfrac{\num{6,67}\times 10^{-11}\times \num{6,42}\times 10^{23}\times 500}{2\times \num{9,4}\times 10^{6}} \approx \num{1,14}\times 10^{9}\ \mathrm{J}$ ; $E_p = -2E_c \approx -\num{2,28}\times 10^{9}\ \mathrm{J}$ ; $E_m = -E_c \approx -\num{1,14}\times 10^{9}\ \mathrm{J}$.

\textbf{3.} Oui : le satellite est en orbite fermée, donc dans un état lié à Mars. Son énergie mécanique doit être strictement négative, ce que confirme le calcul.
\end{corrige}

\coursec{Synthèse et méthodes de résolution}

Nous avons découvert, dans la section précédente, que l'énergie mécanique d'un satellite en orbite circulaire s'écrit $E_m = -\dfrac{GMm}{2r}$ : une formule élégante qui résume à elle seule le lien entre énergie cinétique et énergie potentielle dans un champ de gravitation. Il est temps maintenant de prendre de la hauteur — comme un satellite, justement ! — et d'organiser tout ce que nous avons appris en une véritable \cle{boîte à outils} pour résoudre, méthodiquement et sans faute, n'importe quel exercice du Baccalauréat sur les planètes et les satellites.

\courssub{Bilan des formules clés}

Rassemblons d'abord l'ensemble des relations établies dans cette leçon. Chacune possède des \cle{conditions d'application} précises qu'il faut connaître par cœur : c'est là que se jouent la plupart des points au Bac.

\begin{aretenir}[Formulaire de la leçon]
Pour un satellite de masse $m$ en orbite \textbf{circulaire} de rayon $r$ autour d'un astre de masse $M$ (référentiel astrocentrique galiléen) :
\begin{itemize}
\item Force de gravitation : $F = \dfrac{GMm}{r^2}$ (toujours vraie, orbite quelconque) ;
\item Vitesse : $v = \sqrt{\dfrac{GM}{r}}$ (orbite circulaire uniquement) ;
\item Période : $T = 2\pi\sqrt{\dfrac{r^3}{GM}}$ (orbite circulaire) ;
\item Troisième loi de Kepler : $\dfrac{T^2}{r^3} = \dfrac{4\pi^2}{GM}$ (orbite circulaire ; pour une ellipse, remplacer $r$ par le demi-grand axe $a$) ;
\item Satellite géostationnaire : $T = T_s = \SI{86164}{s}$ (jour sidéral), orbite dans le plan équatorial, altitude $h \approx \SI{35786}{km}$ ;
\item Énergie mécanique : $E_m = -\dfrac{GMm}{2r}$ (orbite circulaire), avec $E_c = \dfrac{GMm}{2r}$ et $E_p = -\dfrac{GMm}{r}$.
\end{itemize}
\end{aretenir}

\begin{attention}[Le rayon $r$ n'est pas l'altitude $h$]
L'erreur la plus fréquente au Baccalauréat consiste à remplacer $r$ par l'altitude $h$ dans les formules. Or $r = R + h$, où $R$ est le rayon de l'astre. Pour un satellite géostationnaire : $r = \SI{6378}{km} + \SI{35786}{km} = \SI{42164}{km}$. Écris systématiquement $r = R + h$ en tête de ta copie avant tout calcul.
\end{attention}

\courssub{Méthode générale de mise en équation}

Tout exercice de mécanique céleste se résout selon le même schéma, hérité de la méthode générale de la dynamique. Le voici, décliné pour les satellites.

\begin{methode}[Fiche de résolution type : un exercice de satellite en 5 étapes]
\begin{enumerate}
\item \textbf{Référentiel.} Préciser le référentiel d'étude (géocentrique pour un satellite terrestre, héliocentrique pour une planète) et le supposer galiléen.
\item \textbf{Système et bilan des forces.} Système : le satellite (masse $m$). Une seule force : l'attraction gravitationnelle $\vec{F} = -\dfrac{GMm}{r^2}\,\vec{u}_r$, dirigée vers le centre de l'astre. Faire un schéma avec le vecteur force et le repère de Frenet $(\vec{u}_t, \vec{u}_n)$.
\item \textbf{Théorème du centre d'inertie (TCI).} Projeter $\vec{F} = m\vec{a}$ dans la base de Frenet :
\[ \dfrac{GMm}{r^2} = m\,\dfrac{v^2}{r} \quad (\text{composante normale}), \qquad m\dfrac{dv}{dt} = 0 \quad (\text{composante tangentielle}). \]
La composante tangentielle nulle prouve que le mouvement circulaire est \textbf{uniforme}.
\item \textbf{Exploitation.} Simplifier par $m$ (la trajectoire ne dépend pas de la masse du satellite !) et tirer la grandeur demandée : $v$, $T$, $r$, $M$\dots
\item \textbf{Vérification.} Contrôler l'homogénéité (analyse dimensionnelle), convertir toutes les données en unités SI (mètres, secondes, kilogrammes) et donner le résultat avec un nombre raisonnable de chiffres significatifs et son unité.
\end{enumerate}
\end{methode}

\begin{attention}[Erreur fréquente : mélanger les unités]
La constante gravitationnelle $G = \num{6,67e-11}~\mathrm{N\,m^2\,kg^{-2}}$ est exprimée en unités SI. Si tu injectes $r$ en \textbf{kilomètres} ou $T$ en \textbf{heures}, le résultat sera faux de plusieurs ordres de grandeur. Règle d'or : \textbf{tout convertir en mètres et en secondes avant le calcul}, puis reconvertir le résultat si l'énoncé le demande. Exemple : $r = \SI{42164}{km} = \num{4,2164e7}{m}$ ; $T = \SI{86164}{s}$ (et non $24 \times 3600 = \SI{86400}{s}$ pour un géostationnaire !).
\end{attention}

\courssub{Exploiter la troisième loi de Kepler par comparaison}

Voici une technique redoutable, très appréciée des correcteurs : pour comparer deux satellites d'un \textbf{même astre}, on n'a besoin ni de $G$ ni de $M$ ! En effet, la troisième loi de Kepler donne $\dfrac{T^2}{r^3} = \dfrac{4\pi^2}{GM}$, constante identique pour les deux satellites.

\begin{methode}[Comparaison de deux orbites par le rapport $T^2/r^3$]
\begin{enumerate}
\item Identifier les deux satellites (1) et (2) gravitant autour du \textbf{même} astre.
\item Écrire l'égalité des constantes de Kepler :
\[ \dfrac{T_1^2}{r_1^3} = \dfrac{T_2^2}{r_2^3} \qquad \text{donc} \qquad T_2 = T_1 \left(\dfrac{r_2}{r_1}\right)^{3/2}. \]
\item Vérifier que les deux rayons sont dans la \textbf{même unité} (le rapport $r_2/r_1$ est sans dimension, n'importe quelle unité convient tant qu'elle est commune).
\item Calculer, puis reconvertir le résultat dans l'unité demandée.
\end{enumerate}
Cette méthode s'applique aussi à deux planètes du système solaire (avec $T$ en années et $r$ en u.a., le rapport reste valable).
\end{methode}

\begin{exoresolu}[Période d'un satellite GPS à partir du géostationnaire]
\textbf{Énoncé.} Les satellites du système GPS évoluent sur des orbites circulaires d'altitude $h \approx \SI{20200}{km}$. Un satellite géostationnaire a une période $T_g = \SI{86164}{s}$ et une altitude $h_g = \SI{35786}{km}$. On donne $R_T = \SI{6378}{km}$. Déterminer la période $T_{GPS}$ d'un satellite GPS sans utiliser $G$ ni $M_T$.
\tcblower
\textbf{Solution.} Les deux satellites gravitent autour de la Terre : on applique la troisième loi de Kepler par comparaison.

Rayons des orbites :
\[ r_g = R_T + h_g = 6378 + 35786 = \SI{42164}{km}, \qquad r_{GPS} = R_T + h_{GPS} = 6378 + 20200 = \SI{26578}{km}. \]

Égalité des constantes de Kepler :
\[ \dfrac{T_{GPS}^2}{r_{GPS}^3} = \dfrac{T_g^2}{r_g^3} \quad \Longrightarrow \quad T_{GPS} = T_g \left(\dfrac{r_{GPS}}{r_g}\right)^{3/2}. \]

Application numérique :
\begin{align*}
T_{GPS} &= \num{86164} \times \left(\dfrac{\num{26578}}{\num{42164}}\right)^{3/2} = \num{86164} \times (\num{0,6304})^{3/2} \\
&= \num{86164} \times \num{0,5005} \approx \num{4,31e4}{s}.
\end{align*}

Conversion : $T_{GPS} \approx \dfrac{\num{43125}}{3600} \approx \SI{11,98}{h}$, soit environ \textbf{11 h 58 min}.

Vérification de cohérence : le satellite GPS est plus bas que le géostationnaire, donc plus rapide (troisième loi de Kepler : $T$ croît avec $r$) : sa période doit être inférieure à $\SI{24}{h}$. C'est bien le cas. Un satellite GPS effectue ainsi presque exactement deux révolutions par jour sidéral.
\end{exoresolu}

\courssub{Déterminer la masse d'un astre à partir d'un satellite}

C'est l'une des plus belles applications de la leçon : « peser » un astre sans jamais le toucher, simplement en observant le mouvement d'un de ses satellites.

\begin{methode}[Masse d'un astre par la troisième loi de Kepler]
\begin{enumerate}
\item Relever le rayon $r$ de l'orbite du satellite (attention : $r = R + h$, en mètres) et sa période $T$ (en secondes).
\item Isoler $M$ dans la troisième loi de Kepler :
\[ \dfrac{T^2}{r^3} = \dfrac{4\pi^2}{GM} \quad \Longrightarrow \quad \boxed{M = \dfrac{4\pi^2 r^3}{G T^2}}. \]
\item Vérifier l'homogénéité : $G$ en $\mathrm{m^3\,kg^{-1}\,s^{-2}}$, $r^3$ en $\mathrm{m^3}$, $T^2$ en $\mathrm{s^2}$.
\[ [M] = \dfrac{\mathrm{m^3}}{\mathrm{m^3\,kg^{-1}\,s^{-2}} \cdot \mathrm{s^2}} = \mathrm{kg}. \quad \text{La formule est homogène.} \]
\item Conclure avec le nombre de chiffres significatifs adapté aux données.
\end{enumerate}
\end{methode}

\begin{exemplebox}[« Peser » Jupiter grâce à Io]
Io, satellite de Jupiter, décrit une orbite quasi circulaire de rayon $r = \num{4,22e8}{m}$ en $T = \SI{1,53e5}{s}$ (environ 1,77 jour). Alors :
\[ M_J = \dfrac{4\pi^2 r^3}{G T^2} = \dfrac{4\pi^2 \times (\num{4,22e8})^3}{\num{6,67e-11} \times (\num{1,53e5})^2} \approx \num{1,9e27}{kg}, \]
soit environ 318 fois la masse de la Terre. C'est exactement ainsi que les astronomes déterminent les masses des planètes, des étoiles et même des trous noirs.
\end{exemplebox}

\begin{savaistu}[Le Cameroun et les satellites géostationnaires]
Le Cameroun, comme de nombreux pays africains, dépend fortement des satellites géostationnaires pour ses télécommunications (téléphonie mobile MTN/Orange, internet, télévision par satellite). Le satellite \textbf{RASCOM-QAF1} (Regional African Satellite Communication Organization), positionné à \SI{2,85}{\degree} Est sur l'orbite géostationnaire, assure une couverture du continent africain. Les antennes paraboliques installées sur les toits de Yaoundé, Douala ou dans les zones rurales (électrification rurale, télémédecine) pointent toutes vers cette orbite unique à \SI{35786}{km} d'altitude. La précision de l'altitude et de la période (jour sidéral) est donc une nécessité économique et stratégique, pas seulement un exercice de physique !
\end{savaistu}

\courssub{Arbre méthodologique : quelle formule choisir ?}

Face à un énoncé, la première question à se poser est : « Que demande-t-on ? » L'arbre ci-dessous te guide vers la bonne relation.

\begin{popfigure}
\begin{center}
\begin{tikzpicture}[
  font=\small,
  question/.style={rectangle, rounded corners, draw=popblue, fill=popblueL, text width=4.8cm, align=center, minimum height=1.1cm},
  meth/.style={rectangle, rounded corners, draw=poporange, fill=poporangeL, text width=4.0cm, align=center, minimum height=0.9cm},
  lien/.style={-{Stealth[length=2.5mm]}, thick, popdark},
  node distance=1.2cm and 1.5cm
]
\node[question, align=center] (q) {\textbf{Que demande l'énoncé ?}\\ (satellite en orbite circulaire)};
\node[meth, align=center] (v) [below left=of q, xshift=-1.5cm] {Vitesse $v$ ?\\ $v = \sqrt{GM/r}$};
\node[meth, align=center] (t) [below left=of q, xshift=1.5cm] {Période $T$ ou rayon $r$ ?\\ $T = 2\pi\sqrt{r^3/GM}$};
\node[meth, align=center] (m) [below right=of q, xshift=-1.5cm] {Masse $M$ de l'astre ?\\ $M = 4\pi^2 r^3/(GT^2)$};
\node[meth, align=center] (c) [below right=of q, xshift=1.5cm] {Comparer 2 satellites ?\\ $T_1^2/r_1^3 = T_2^2/r_2^3$};
\node[meth, align=center] (e) [below=of t] {Énergie, libération ?\\ $E_m = -GMm/(2r)$};
\node[meth, align=center] (g) [below=of m] {Géostationnaire ?\\ $T = T_s$, plan équatorial, $r = R_T + h$};
\draw[lien] (q) -- (v);
\draw[lien] (q) -- (t);
\draw[lien] (q) -- (m);
\draw[lien] (q) -- (c);
\draw[lien] (t) -- (e);
\draw[lien] (m) -- (g);
\end{tikzpicture}
\end{center}
\legende{Arbre méthodologique : à chaque question de l'énoncé correspond une relation de la leçon. Avant tout calcul : référentiel, schéma, conversion en unités SI.}
\end{popfigure}

\courssub{Tableau récapitulatif : grandeurs, unités et pièges}

\begin{center}
\rowcolors{2}{popcream}{white}
\begin{tabularx}{\linewidth}{|l|c|c|X|}
\hline
\rowcolor{popblueL} \textbf{Grandeur} & \textbf{Symbole} & \textbf{Unité SI} & \textbf{Pièges fréquents} \\
\hline
Rayon de l'orbite & $r$ & m & $r = R + h$, jamais $h$ seul ; convertir les km en m \\
\hline
Altitude & $h$ & m & Ne figure dans aucune formule directement \\
\hline
Période & $T$ & s & Jour sidéral $T_s = \SI{86164}{s} \neq$ jour solaire $\SI{86400}{s}$ \\
\hline
Vitesse orbitale & $v$ & $\mathrm{m\,s^{-1}}$ & Convertir les km/h ou km/s : $1~\mathrm{km/s} = 10^3~\mathrm{m/s}$ \\
\hline
Masse de l'astre & $M$ & kg & Ne pas confondre avec $m$, masse du satellite (qui se simplifie) \\
\hline
Constante de gravitation & $G$ & $\mathrm{m^3\,kg^{-1}\,s^{-2}}$ & Impose le système SI complet pour tout le calcul \\
\hline
Unité astronomique & u.a. & m & $1~\text{u.a.} = \num{1,496e11}{m}$ (distance Terre--Soleil) \\
\hline
Énergie mécanique & $E_m$ & J & $E_m < 0$ pour un satellite lié ; $E_m = -GMm/(2r)$ \\
\hline
\end{tabularx}
\end{center}

\begin{attention}[Jour sidéral contre jour solaire]
Le satellite géostationnaire doit tourner \textbf{avec la Terre par rapport aux étoiles}, pas par rapport au Soleil. Sa période est donc le jour sidéral $T_s = \SI{86164}{s} \approx 23~\text{h}~56~\text{min}~4~\text{s}$, et non 24 h. Utiliser $\SI{86400}{s}$ dans le calcul de l'altitude donne $h \approx \SI{35860}{km}$ au lieu de $\SI{35786}{km}$ : une erreur de près de 75 km, sanctionnée au Bac.
\end{attention}

\courssub{Entraînement aux conversions}

Les conversions sont le nerf de la guerre. Voici les réflexes à acquérir.

\begin{exemplebox}[Conversions usuelles en mécanique céleste]
\begin{itemize}
\item \textbf{Unité astronomique} : $1~\text{u.a.} = \num{1,496e11}{m}$. Ainsi Mars, à $1,52$~u.a. du Soleil, orbite à $r = 1,52 \times \num{1,496e11} = \num{2,27e11}{m}$.
\item \textbf{Année} : $1~\text{an} = 365,25 \times 86400 \approx \num{3,156e7}{s}$.
\item \textbf{Vitesse} : la Station spatiale file à $\SI{7,7}{km/s} = \num{7,7e3}{m/s}$.
\item \textbf{Jour sidéral} : $T_s = 23~\text{h}~56~\text{min}~4~\text{s} = 23 \times 3600 + 56 \times 60 + 4 = \SI{86164}{s}$.
\item \textbf{Altitude géostationnaire} : $h = \SI{35786}{km} = \num{3,5786e7}{m}$ ; rayon orbital $r = \num{4,2164e7}{m}$.
\end{itemize}
\end{exemplebox}

\begin{exercice}[Pour s'entraîner : la Lune et la masse de la Terre]
La Lune décrit autour de la Terre une orbite quasi circulaire de rayon $r = \num{3,84e8}{m}$ en $T = 27,3$ jours.
\begin{enumerate}
\item Convertir la période en secondes.
\item En déduire la masse de la Terre. On donne $G = \num{6,67e-11}~\mathrm{m^3\,kg^{-1}\,s^{-2}}$.
\item Calculer la vitesse orbitale de la Lune en $\mathrm{km/s}$.
\end{enumerate}
\end{exercice}

\begin{corrige}[Exercice : la Lune et la masse de la Terre]
\begin{enumerate}
\item $T = 27,3 \times 86400 = \num{2,36e6}{s}$.
\item Troisième loi de Kepler : $M_T = \dfrac{4\pi^2 r^3}{G T^2} = \dfrac{4\pi^2 \times (\num{3,84e8})^3}{\num{6,67e-11} \times (\num{2,36e6})^2} \approx \num{6,0e24}{kg}$. C'est bien la valeur connue de la masse terrestre.
\item $v = \sqrt{\dfrac{GM_T}{r}} = \sqrt{\dfrac{\num{6,67e-11} \times \num{6,0e24}}{\num{3,84e8}}} \approx \num{1,02e3}{m/s} \approx \SI{1,0}{km/s}$.
\end{enumerate}
\end{corrige}

\begin{aretenir}[L'essentiel de la synthèse]
\begin{itemize}
\item Une seule force, une seule méthode : référentiel galiléen, TCI projeté dans Frenet, exploitation.
\item La trajectoire d'un satellite ne dépend pas de sa masse $m$ : seuls comptent $M$ et $r$.
\item Pour comparer deux satellites d'un même astre, utilise $\dfrac{T_1^2}{r_1^3} = \dfrac{T_2^2}{r_2^3}$ : ni $G$ ni $M$ ne sont nécessaires.
\item Pour « peser » un astre : $M = \dfrac{4\pi^2 r^3}{G T^2}$.
\item Toujours : $r = R + h$, tout en mètres et en secondes, jour sidéral pour le géostationnaire.
\end{itemize}
Tu disposes maintenant de toutes les clés pour aborder sereinement les exercices du Baccalauréat sur les mouvements des planètes et des satellites. Entraîne-toi régulièrement : la rigueur dans la rédaction et les conversions fait toute la différence !
\end{aretenir}

\newpage

\coursec{Activité d'intégration}

\coursec{Mouvement des planètes et des satellites}

\courssub{Objectifs de l'activité}

À l'issue de cette activité d'intégration, tu seras capable de :
\begin{itemize}
\item justifier, par le bilan des forces et la nature du mouvement, les conditions géométriques imposées à une orbite géostationnaire (plan équatorial, sens de rotation, période égale au jour sidéral) ;
\item mettre en équation le mouvement circulaire uniforme d'un satellite dans le champ de gravitation terrestre et en déduire le rayon de l'orbite, l'altitude et la vitesse ;
\item exploiter la troisième loi de Kepler sous forme de rapport pour comparer deux satellites sans connaître ni $G$ ni la masse de la Terre ;
\item calculer l'énergie mécanique d'un satellite et la vitesse de libération, puis en tirer des conclusions concrètes sur un lancement ;
\item rédiger une note technique argumentée, avec des résultats numériques assortis d'unités et de chiffres significatifs cohérents.
\end{itemize}

\courssub{Situation}

\textbf{Mission Nkongsamba-Sat : le Cameroun lance CamSat-1.}

Dans le cadre de son programme de souveraineté numérique, le ministère des Postes et Télécommunications (MINPOSTEL) confie à ton équipe d'ingénieurs-physiciens une étude de faisabilité. Le Cameroun souhaite lancer \textbf{CamSat-1}, un satellite de télécommunications de masse $m = \SI{1500}{kg}$, destiné à assurer la couverture télévisuelle et Internet des zones rurales (de l'Extrême-Nord au golfe de Guinée) sans interruption. Pour cela, le satellite doit rester \emph{fixe} par rapport au sol, au-dessus d'un point de l'équateur situé dans le golfe de Guinée, face au Cameroun.

Un second satellite, \textbf{CamSat-2}, de surveillance forestière (suivi de l'exploitation illégale dans le bassin du Congo), doit quant à lui effectuer un tour complet de la Terre en exactement $T_2 = \SI{2,0}{h}$.

\textbf{Données :}
\begin{itemize}
\item constante de gravitation universelle : $G = \SI{6,67e-11}{m^3.kg^{-1}.s^{-2}}$ ;
\item masse de la Terre : $M_T = \SI{5,97e24}{kg}$ ;
\item rayon terrestre (à l'équateur) : $R_T = \SI{6,38e3}{km}$ ;
\item jour sidéral (période de rotation propre de la Terre) : $T_{sid} = \SI{86164}{s} \approx \SI{23,93}{h}$ ;
\item on suppose la Terre sphérique, à répartition de masse à symétrie sphérique, et les orbites circulaires ; le référentiel géocentrique est supposé galiléen.
\end{itemize}

\textbf{Problème à résoudre :} ton équipe doit dimensionner les deux orbites et rédiger une note technique au MINPOSTEL précisant l'altitude et la vitesse de CamSat-1, l'altitude de CamSat-2, ainsi qu'une première estimation énergétique du lancement.

\begin{popfigure}
\begin{tikzpicture}[scale=0.8, every node/.style={font=\small}]
% Terre
\fill[popblueL!50] (0,0) circle (2cm);
\draw[thick, popblue] (0,0) circle (2cm);
\node at (0,0) {\textbf{Terre}};
% Axe de rotation
\draw[->, thick, popdark] (0,-2.5) -- (0,2.5) node[above] {$\vec{\omega}_T$};
\draw[dashed, popmuted] (-2.5,0) -- (2.5,0) node[right] {Plan équatorial};
% Satellite géostationnaire
\coordinate (Sat) at (0,5);
\fill[poporange] (Sat) circle (0.15cm);
\node[above right] at (Sat) {\textbf{CamSat-1}};
% Liaison Terre-Satellite
\draw[thick, poporange] (0,2) -- (Sat);
% Position au sol (Golfe de Guinée)
\coordinate (Sol) at (0,2);
\fill[popgreen] (Sol) circle (0.08cm);
\node[left] at (Sol) {Golfe de Guinée};
% Flèche rotation satellite
\draw[->, thick, poporange] (1,5) arc (0:90:1cm) node[midway, right] {Même sens};
% Altitude
\draw[<->, popmuted] (2.2,2) -- (2.2,5) node[midway, right] {$h_1 \approx \SI{35800}{km}$};
\draw[<->, popmuted] (0,0) -- (0,2) node[midway, left] {$R_T$};
\draw[<->, popmuted] (0,0) -- (0,5) node[midway, left] {$r_1$};
\end{tikzpicture}
\legende{Orbite géostationnaire de CamSat-1 : le satellite reste fixe au-dessus du golfe de Guinée. L'orbite est circulaire, équatoriale, de période $T_{sid}$ et de rayon $r_1 = R_T + h_1$.}
\end{popfigure}

\courssub{Consignes}

\textbf{Partie A — Orbite de CamSat-1 (géostationnaire).}
\begin{enumerate}
\item Rappelle les trois conditions qu'une orbite géostationnaire doit vérifier. Justifie, à l'aide du bilan des forces (la force de gravitation est dirigée vers le centre de la Terre), que le plan de l'orbite contient nécessairement le centre de la Terre, puis qu'une orbite fixe au-dessus du golfe de Guinée doit être \emph{équatoriale}.
\item Dans le référentiel géocentrique, fais le bilan des forces appliquées à CamSat-1 et représente-les sur un schéma. En appliquant la deuxième loi de Newton dans la base de Frenet, établis l'expression de la vitesse $v$ du satellite en fonction de $G$, $M_T$ et du rayon $r$ de l'orbite. Vérifie l'homogénéité de ton expression.
\item Montre que le mouvement est uniforme, puis établis la troisième loi de Kepler : $\dfrac{T^2}{r^3} = \dfrac{4\pi^2}{G M_T}$.
\item En déduis la valeur numérique du rayon $r_1$ de l'orbite géostationnaire, puis l'altitude $h_1$ de CamSat-1. Compare avec la valeur usuelle d'environ \SI{35800}{km}.
\item Calcule la vitesse $v_1$ de CamSat-1 sur son orbite, en \si{km.s^{-1}} puis en \si{km.h^{-1}}.
\end{enumerate}

\textbf{Partie B — Orbite de CamSat-2 (surveillance, $T_2 = \SI{2,0}{h}$).}
\begin{enumerate}
\setcounter{enumi}{5}
\item En exploitant la troisième loi de Kepler sous forme du rapport $\dfrac{T_2^2}{r_2^3} = \dfrac{T_1^2}{r_1^3}$ (sans utiliser $G$ ni $M_T$), détermine le rayon $r_2$ de l'orbite de CamSat-2, puis son altitude $h_2$.
\item Vérifie que cette orbite est compatible avec un satellite (c'est-à-dire $r_2 > R_T$). Calcule la vitesse $v_2$ de CamSat-2 et compare-la à $v_1$ : commente la phrase « plus un satellite est proche de la Terre, plus il file vite ».
\end{enumerate}

\textbf{Partie C — Aspects énergétiques (complément Bac).}
\begin{enumerate}
\setcounter{enumi}{7}
\item Calcule l'énergie mécanique $E_m$ de CamSat-1 sur son orbite. Que signifie son signe ?
\item Calcule la vitesse de libération $v_{lib}$ depuis la surface de la Terre. Sachant qu'un lanceur doit déjà fournir environ \SI{8}{km.s^{-1}} pour placer une charge en orbite basse, discute brièvement de l'effort énergétique nécessaire pour atteindre l'orbite géostationnaire.
\end{enumerate}

\textbf{Restitution :} rédigez, par groupe, une note technique d'une page au MINPOSTEL : schéma de l'orbite, tableau des grandeurs calculées (avec unités et 3 chiffres significatifs), et une phrase de conclusion argumentée sur la faisabilité.

\courssub{Ressources à mobiliser}

\begin{aretenir}[Boîte à outils de l'activité]
\begin{itemize}
\item Force de gravitation : $\vec{F} = -\dfrac{G M_T m}{r^2}\,\vec{u}_r$, dirigée vers le centre de la Terre.
\item Deuxième loi de Newton dans la base de Frenet, pour un mouvement circulaire : $a_N = \dfrac{v^2}{r} = \dfrac{G M_T}{r^2}$, $a_T = \dot{v} = 0$ (mouvement uniforme).
\item Vitesse orbitale : $v = \sqrt{\dfrac{G M_T}{r}}$ ; période : $T = 2\pi\sqrt{\dfrac{r^3}{G M_T}}$.
\item Troisième loi de Kepler : $\dfrac{T^2}{r^3} = \dfrac{4\pi^2}{G M_T}$ = constante pour tous les satellites d'un même astre.
\item Condition géostationnaire : orbite circulaire, dans le plan équatorial, dans le sens de rotation de la Terre, avec $T = T_{sid}$.
\item Énergie mécanique : $E_m = E_c + E_p = -\dfrac{G M_T m}{2r}$ ; vitesse de libération : $v_{lib} = \sqrt{\dfrac{2 G M_T}{R_T}}$.
\end{itemize}
\end{aretenir}

\begin{methode}[Dimensionner une orbite géostationnaire]
\begin{enumerate}
\item Identifier la période imposée : $T = T_{sid}$ (jour sidéral, \SI{86164}{s}, \textbf{pas} le jour solaire de 24 h).
\item Appliquer la 3e loi de Kepler : $r^3 = \dfrac{G M_T T^2}{4\pi^2}$.
\item Calculer $r$, puis l'altitude $h = r - R_T$.
\item Vérifier $h > 0$ et l'ordre de grandeur ($\sim \SI{3,6e4}{km}$).
\item Calculer la vitesse $v = \sqrt{GM_T/r} = 2\pi r/T$.
\end{enumerate}
\end{methode}

\begin{methode}[Comparer deux satellites par la 3e loi de Kepler (sans $G$, ni $M_T$)]
\begin{enumerate}
\item Écrire l'égalité des rapports : $\dfrac{T_1^2}{r_1^3} = \dfrac{T_2^2}{r_2^3}$.
\item Isoler $r_2 = r_1 \left(\dfrac{T_2}{T_1}\right)^{2/3}$.
\item Calculer $r_2$ et $h_2 = r_2 - R_T$.
\item Vérifier $r_2 > R_T$ (sinon l'orbite est impossible : collision ou atmosphère).
\item Calculer $v_2 = \dfrac{2\pi r_2}{T_2}$ ou $v_2 = \sqrt{GM_T/r_2}$ et comparer à $v_1$.
\end{enumerate}
\end{methode}

\begin{attention}[Pièges fréquents au Bac]
\begin{itemize}
\item \textbf{Jour solaire vs jour sidéral :} La période de rotation de la Terre par rapport aux étoiles fixes est $T_{sid} = \SI{86164}{s}$ (\SI{23,93}{h}), \textbf{pas} \SI{86400}{s} (24 h). Utiliser 24 h décale le satellite de \SI{2600}{km} par jour !
\item \textbf{Altitude vs Rayon :} L'altitude $h$ se compte depuis la surface ($h = r - R_T$). Ne pas confondre $r$ (distance au centre) et $h$.
\item \textbf{Unités :} $R_T$ est donné en \si{km}, $G$ en \si{m^3.kg^{-1}.s^{-2}}. Convertir \textbf{tout} en SI (m, kg, s) avant de calculer, puis reconvertir le résultat final en \si{km} si demandé.
\item \textbf{Signe de l'énergie :} $E_m = -\dfrac{GMm}{2r} < 0$ pour un satellite lié. Oublier le signe moins est une erreur grave d'interprétation physique.
\item \textbf{Vitesse de libération :} $v_{lib} = \sqrt{2}\,v_{circulaire}(R_T) \approx \SI{11,2}{km.s^{-1}}$. Ne pas confondre avec la vitesse orbitale basse ($\approx \SI{7,9}{km.s^{-1}}$).
\end{itemize}
\end{attention}

\begin{savaistu}[Le Cameroun et l'espace]
Le Cameroun ne possède pas encore de lanceur ni de base de lancement, mais il est utilisateur de l'espace via les opérateurs MTN, Orange, CAMTEL qui louent des capacités sur des satellites géostationnaires (ex. RASCOM-QAF1, Intelsat, Eutelsat) positionnés au-dessus de l'Afrique (souvent à \SI{3}{\degree} Est ou \SI{11}{\degree} Ouest). L'Agence Nationale des Technologies de l'Information et de la Communication (ANTIC) pilote la stratégie numérique. Un satellite propre (CamSat) permettrait de réduire la facture de location (plusieurs milliards de FCFA/an) et de maîtriser la couverture du territoire, notamment pour la télémédecine, l'éducation à distance et la surveillance des frontières/forêts. Le défi énergétique (vitesse de libération \SI{11,2}{km.s^{-1}}) explique pourquoi les pays émergents passent souvent par des lancements partagés (rideshare) sur Ariane, Falcon 9, Soyouz ou Longue Marche.
\end{savaistu}

\courssub{Démarche et corrigé commenté}

\begin{exoresolu}[Note technique — Mission Nkongsamba-Sat (corrigé complet)]
Énoncé : dimensionner l'orbite géostationnaire de CamSat-1 ($m = \SI{1500}{kg}$), déterminer l'altitude de CamSat-2 ($T_2 = \SI{2,0}{h}$) par comparaison, et estimer les grandeurs énergétiques du lancement.

\tcblower

\textbf{A.1. Conditions géostationnaires.} La seule force subie par le satellite est la gravitation, centrale : $\vec{F} = -\dfrac{G M_T m}{r^2}\vec{u}_r$ pointe vers le centre $O$ de la Terre. D'après la deuxième loi de Newton, l'accélération $\vec{a} = \vec{F}/m$ est aussi dirigée vers $O$ ; le mouvement est donc plan et le plan de la trajectoire contient $O$. Pour que le satellite paraisse \emph{fixe} au-dessus d'un point du sol (golfe de Guinée), trois conditions sont \textbf{nécessaires et suffisantes} :
\begin{itemize}
\item une orbite \textbf{circulaire} (vitesse angulaire constante, sinon le satellite oscillerait en longitude) ;
\item un plan orbital \textbf{équatorial} (contenant l'axe de rotation) : si le plan était incliné, le satellite oscillerait en latitude au nord et au sud de l'équateur (trajectoire en « 8 » vue du sol) et ne resterait pas fixe au-dessus du golfe de Guinée ;
\item une période égale au \textbf{jour sidéral} $T_{sid} = \SI{86164}{s}$, dans le \textbf{sens de rotation} de la Terre (ouest $\to$ est).
\end{itemize}

\begin{popfigure}
\begin{tikzpicture}[scale=1, every node/.style={font=\small}]
% Satellite
\fill[poporange] (0,0) circle (0.2cm);
\node[below] at (0,-0.4) {Satellite $m$};
% Vecteur vitesse (tangentiel)
\draw[->, very thick, popgreen] (0,0) -- (1.5,0) node[midway, above] {$\vec{v} = \vec{\tau}\,v$};
% Vecteur force/accélération (normal entrant)
\draw[->, very thick, popred] (0,0) -- (0,-1.5) node[midway, right] {$\vec{F} = m\vec{a}_N = -m\dfrac{v^2}{r}\vec{n}$};
% Base de Frenet
\node[left] at (-0.3,0) {$\vec{\tau}$};
\node[above] at (0,-0.3) {$\vec{n}$};
% Centre de la Terre
\node[below=2.5cm] at (0,0) {Centre Terre $O$};
\draw[dashed, popmuted] (0,-1.5) -- (0,-3);
\fill[popblue] (0,-3) circle (0.1cm);
% Coordonnées polaires
\draw[<->, popmuted] (0.3,0) arc (0:-90:0.3cm);
\node at (0.4,-0.4) {$\pi/2$};
\end{tikzpicture}
\legende{Bilan des forces et base de Frenet pour un satellite en orbite circulaire. La force de gravitation fournit l'accélération normale $a_N = v^2/r$. L'accélération tangentielle $a_T = \dot{v}$ est nulle : le mouvement est uniforme.}
\end{popfigure}

\textbf{A.2. Mise en équation.} Dans le référentiel géocentrique supposé galiléen, le bilan des forces se réduit à la gravitation :
\[
\vec{F} = -\frac{G M_T m}{r^2}\,\vec{u}_r.
\]
La deuxième loi de Newton, $\vec{F} = m\vec{a}$, projetée dans la base de Frenet ($\vec{\tau}$ tangent, $\vec{n}$ normal entrant) donne :
\begin{align*}
\text{Composante normale :}\quad & m\,\frac{v^2}{r} = \frac{G M_T m}{r^2} \\
\text{Composante tangentielle :}\quad & m\,\dot{v} = 0
\end{align*}
D'où l'expression de la vitesse orbitale :
\[
\boxed{v = \sqrt{\frac{G M_T}{r}}}.
\]
\textbf{Analyse dimensionnelle :} $[G] = \si{m^3.kg^{-1}.s^{-2}}$, $[M_T] = \si{kg}$, $[r] = \si{m}$.
$[GM_T/r] = \si{m^3.kg^{-1}.s^{-2}} \times \si{kg} \times \si{m^{-1}} = \si{m^2.s^{-2}}$.
La racine carrée donne $\si{m.s^{-1}}$, dimension d'une vitesse : \textbf{l'expression est homogène}.
L'équation tangentielle $\dot{v}=0$ prouve que le mouvement est \textbf{circulaire uniforme} (vitesse scalaire constante).

\textbf{A.3. Troisième loi de Kepler.} Pour un mouvement circulaire uniforme, $v = \dfrac{2\pi r}{T}$. En substituant dans l'expression de $v$ :
\[
\frac{4\pi^2 r^2}{T^2} = \frac{G M_T}{r}
\quad\Longrightarrow\quad
\boxed{\frac{T^2}{r^3} = \frac{4\pi^2}{G M_T}}.
\]
Ce rapport ne dépend \textbf{que de l'astre central} (la Terre ici), pas du satellite : c'est la troisième loi de Kepler, valable pour toute orbite circulaire (et même elliptique avec $a$ au lieu de $r$).

\textbf{A.4. Rayon et altitude de CamSat-1.} On isole $r_1$ à partir de $T_1 = T_{sid}$ :
\[
r_1 = \left(\frac{G M_T T_{sid}^2}{4\pi^2}\right)^{1/3}.
\]
Calculs intermédiaires (valeurs arrondies à 3 chiffres significatifs) :
\begin{align*}
G M_T &= 6,67\times 10^{-11} \times 5,97\times 10^{24} = \SI{3,98e14}{m^3.s^{-2}} \\
T_{sid}^2 &= (8,6164\times 10^{4})^2 = \SI{7,42e9}{s^2} \\
r_1^3 &= \frac{3,98\times 10^{14} \times 7,42\times 10^{9}}{4\pi^2} \approx \frac{2,95\times 10^{24}}{39,5} \approx \SI{7,47e22}{m^3} \\
r_1 &= \sqrt[3]{7,47\times 10^{22}} \approx \SI{4,21e7}{m} = \SI{4,21e4}{km}.
\end{align*}
Altitude :
\[
h_1 = r_1 - R_T = 4,21\times 10^{4} - 6,38\times 10^{3} = \SI{3,57e4}{km} \approx \SI{35700}{km}.
\]
La valeur usuelle est \SI{35786}{km} (avec $T_{sid} = \SI{86164,0905}{s}$ et $GM_T = \SI{3,986004e14}{m^3.s^{-2}}$). Notre résultat \SI{35700}{km} (3 chiffres significatifs) est en \textbf{parfait accord}.

\textbf{A.5. Vitesse de CamSat-1.}
\[
v_1 = \sqrt{\frac{G M_T}{r_1}} = \sqrt{\frac{3,98\times 10^{14}}{4,21\times 10^{7}}} \approx \SI{3,08e3}{m.s^{-1}} = \SI{3,08}{km.s^{-1}}.
\]
En \si{km.h^{-1}} : $v_1 = 3,08 \times 3600 \approx \SI{1,11e4}{km.h^{-1}}$.
C'est près de 10 fois la vitesse de croisière d'un avion de ligne (\SI{900}{km.h^{-1}}).

\begin{popfigure}
\begin{tikzpicture}[scale=0.6, every node/.style={font=\small}]
% Terre
\fill[popblueL!30] (0,0) circle (3cm);
\draw[thick, popblue] (0,0) circle (3cm);
\node at (0,0) {Terre ($R_T$)};
% Orbite CamSat-2 (basse)
\draw[thick, poporange, dashed] (0,0) circle (3.8cm);
\node[poporange] at (3.8, 0.5) {CamSat-2 : $r_2 \approx \SI{8060}{km}$};
\node[poporange] at (3.8, -0.5) {$h_2 \approx \SI{1680}{km}$};
% Orbite CamSat-1 (géostationnaire)
\draw[thick, popgreen, dashed] (0,0) circle (7.5cm);
\node[popgreen] at (7.5, 0.5) {CamSat-1 : $r_1 \approx \SI{42100}{km}$};
\node[popgreen] at (7.5, -0.5) {$h_1 \approx \SI{35700}{km}$};
% Légende vitesses
\node[align=center] at (0, -5) {
$v_2 \approx \SI{7,03}{km.s^{-1}}$ \quad (orbite basse) \\
$v_1 \approx \SI{3,08}{km.s^{-1}}$ \quad (géostationnaire)
};
\draw[->, thick, popmuted] (0,3) -- (0,3.8) node[midway, right] {$h_2$};
\draw[->, thick, popmuted] (0,3.8) -- (0,7.5) node[midway, right] {$h_1 - h_2$};
\end{tikzpicture}
\legende{Comparaison des orbites de CamSat-1 (géostationnaire, haute, lente) et CamSat-2 (basse, rapide). Plus l'orbite est basse, plus la gravitation est forte et plus le satellite doit aller vite pour ne pas tomber.}
\end{popfigure}

\textbf{B.6. Orbite de CamSat-2 par la 3e loi de Kepler.} Les deux satellites orbitent autour du même astre (la Terre), donc le rapport $T^2/r^3$ est identique :
\[
\frac{T_2^2}{r_2^3} = \frac{T_1^2}{r_1^3}
\quad\Longrightarrow\quad
r_2 = r_1\left(\frac{T_2}{T_1}\right)^{2/3}.
\]
On utilise $T_1 = T_{sid} = \SI{23,93}{h}$ et $T_2 = \SI{2,0}{h}$ (mêmes unités, le rapport est sans dimension) :
\[
\frac{T_2}{T_1} = \frac{2,0}{23,93} \approx 0,0836,
\qquad
\left(\frac{T_2}{T_1}\right)^{2/3} \approx (0,0836)^{0,6667} \approx 0,191.
\]
D'où :
\[
r_2 \approx 4,21\times 10^{4} \times 0,191 \approx \SI{8,04e3}{km},
\qquad
h_2 = r_2 - R_T \approx 8040 - 6380 = \SI{1,66e3}{km}.
\]

\textbf{B.7. Cohérence et vitesse.} $r_2 = \SI{8040}{km} > R_T = \SI{6380}{km}$ : l'orbite est possible (orbite basse, LEO), bien au-dessus de l'atmosphère dense ($\sim \SI{200}{km}$).
Vitesse de CamSat-2 :
\[
v_2 = \frac{2\pi r_2}{T_2} = \frac{2\pi \times 8,04\times 10^{6}}{7200} \approx \SI{7,02e3}{m.s^{-1}} = \SI{7,02}{km.s^{-1}}.
\]
(Vérification par $v_2 = \sqrt{GM_T/r_2} = \sqrt{3,98\times 10^{14} / 8,04\times 10^{6}} \approx \SI{7,03}{km.s^{-1}}$ : cohérent).
On a bien $v_2 > v_1$ (\SI{7,02}{km.s^{-1}} > \SI{3,08}{km.s^{-1}}). Comme $v = \sqrt{GM_T/r}$, la vitesse orbitale \textbf{décroît} quand le rayon \textbf{croît}. La phrase « plus un satellite est proche de la Terre, plus il file vite » est \textbf{physiquement exacte} : pour « contrer » une gravitation plus intense à basse altitude, la force centripète requise ($mv^2/r$) impose une vitesse plus grande. CamSat-2 effectue $24/2 = 12$ tours par jour : idéal pour balayer la forêt du bassin du Congo.

\textbf{C.8. Énergie mécanique de CamSat-1.}
L'énergie potentielle de gravitation (zéro à l'infini) est $E_p = -\dfrac{G M_T m}{r}$.
L'énergie cinétique est $E_c = \dfrac{1}{2}m v^2 = \dfrac{1}{2}\dfrac{G M_T m}{r}$.
L'énergie mécanique totale :
\[
E_m = E_c + E_p = \frac{1}{2}\frac{G M_T m}{r} - \frac{G M_T m}{r} = -\frac{G M_T m}{2r}.
\]
Numériquement pour CamSat-1 ($m = \SI{1500}{kg}$, $r_1 = \SI{4,21e7}{m}$) :
\[
E_m = -\frac{3,98\times 10^{14} \times 1500}{2 \times 4,21\times 10^{7}}
\approx -\frac{5,97\times 10^{17}}{8,42\times 10^{7}}
\approx \SI{-7,09e9}{J}.
\]
$E_m < 0$ : le satellite est dans un \textbf{état lié} (orbite fermée). Il ne peut pas s'échapper seul de l'attraction terrestre ; il lui faudrait une énergie supplémentaire d'au moins $|E_m| = \SI{7,1}{GJ}$ pour atteindre $E_m = 0$ (parabole, évasion).

\textbf{C.9. Vitesse de libération.}
C'est la vitesse minimale à communiquer au sol ($r = R_T$) pour que $E_m = 0$ :
\[
\frac{1}{2}m v_{lib}^2 - \frac{G M_T m}{R_T} = 0
\quad\Longrightarrow\quad
\boxed{v_{lib} = \sqrt{\frac{2 G M_T}{R_T}}}.
\]
Numériquement :
\[
v_{lib} = \sqrt{\frac{2 \times 3,98\times 10^{14}}{6,38\times 10^{6}}}
= \sqrt{1,247\times 10^{8}} \approx \SI{1,117e4}{m.s^{-1}} = \SI{11,2}{km.s^{-1}}.
\]
\textbf{Discussion :} Un lanceur doit déjà fournir $\sim \SI{8}{km.s^{-1}}$ pour atteindre une orbite basse circulaire (vitesse orbitale $v_{orb} = \sqrt{GM_T/R_T} \approx \SI{7,9}{km.s^{-1}}$ + pertes atmosphériques/gravité). Pour aller de l'orbite basse à l'orbite géostationnaire, il faut un transfert de Hohmann (deux impulsions) qui demande un $\Delta v$ supplémentaire d'environ \SI{3,9}{km.s^{-1}} au total. L'énergie totale à fournir est donc de l'ordre de plusieurs dizaines de GJ pour \SI{1500}{kg}, ce qui explique le coût élevé (dizaines de millions de dollars) et l'intérêt pour le Cameroun de mutualiser les lancements ou de louer des capacités existantes (MTN, Orange, CAMTEL) en attendant un programme spatial national mature.
\end{exoresolu}

\courssub{Bilan de l'activité}

Cette mission a permis de mettre en évidence plusieurs idées fondamentales de la leçon :
\begin{itemize}
\item le caractère \textbf{central} de la force de gravitation impose la géométrie de l'orbite : c'est la physique, et non la géographie, qui exige le plan équatorial pour un satellite géostationnaire ;
\item la mise en équation par la deuxième loi de Newton dans la base de Frenet fournit \emph{toutes} les grandeurs orbitales ($v$, $T$, $r$) à partir de deux seules données, $G$ et $M_T$ ;
\item la troisième loi de Kepler est un outil de \textbf{comparaison} remarquable : le rapport $T^2/r^3$ étant le même pour tous les satellites d'un astre, on dimensionne une nouvelle orbite sans refaire aucun calcul de gravitation ;
\item l'altitude géostationnaire ($\approx \SI{35800}{km}$) n'est pas un choix mais une \emph{conséquence} de la période imposée : il n'existe qu'\textbf{une seule} orbite géostationnaire, d'où les négociations internationales (UIT) pour l'attribution des positions orbitales (slots) et fréquences ;
\item l'énergie mécanique négative traduit l'état lié du satellite, et la vitesse de libération ($\SI{11,2}{km.s^{-1}}$) fixe l'échelle de l'effort qu'exige toute conquête spatiale — un défi que le programme spatial camerounais devra relever avec rigueur et patience.
\end{itemize}

\begin{exercice}[Entraînement Bac — Satellite d'observation]
Un satellite d'observation terrestre (type Pléiades) de masse $m = \SI{1000}{kg}$ est placé sur une orbite circulaire héliosynchrone à l'altitude $h = \SI{694}{km}$. Période $T = \SI{98,5}{min}$.
\begin{enumerate}
\item Vérifier que cette période est cohérente avec la 3e loi de Kepler (utiliser $r_1, T_1$ de CamSat-1 comme référence).
\item Calculer sa vitesse orbitale $v$ et son énergie mécanique $E_m$.
\item On veut le transférer sur une orbite circulaire plus haute à $h' = \SI{800}{km}$. Quelle variation d'énergie mécanique $\Delta E_m$ faut-il lui communiquer ? Le satellite doit-il accélérer ou freiner pour monter ?
\end{enumerate}
\end{exercice}

\begin{corrige}[Exercice — Satellite d'observation]
\begin{enumerate}
\item $r = R_T + h = 6380 + 694 = \SI{7074}{km} = 7,074\times 10^6\,\si{m}$.
$T = 98,5 \times 60 = \SI{5910}{s}$.
Rapport $T^2/r^3 = 5910^2 / (7,074\times 10^6)^3 = 3,49\times 10^7 / 3,54\times 10^{20} = 9,86\times 10^{-14}\,\si{s^2.m^{-3}}$.
Pour CamSat-1 : $T_1^2/r_1^3 = (8,6164\times 10^4)^2 / (4,21\times 10^7)^3 = 7,42\times 10^9 / 7,46\times 10^{22} = 9,95\times 10^{-14}\,\si{s^2.m^{-3}}$.
Les deux valeurs sont égales à $\approx 1\%$ près (arrondis) : \textbf{cohérence validée}.
\item $v = 2\pi r/T = 2\pi \times 7,074\times 10^6 / 5910 \approx \SI{7,52e3}{m.s^{-1}} = \SI{7,52}{km.s^{-1}}$.
$E_m = -GM_T m / (2r) = -3,98\times 10^{14} \times 1000 / (2 \times 7,074\times 10^6) \approx \SI{-2,81e10}{J}$.
\item $r' = 6380 + 800 = \SI{7180}{km} = 7,18\times 10^6\,\si{m}$.
$E_m' = -GM_T m / (2r') \approx -3,98\times 10^{17} / (1,436\times 10^7) \approx \SI{-2,77e10}{J}$.
$\Delta E_m = E_m' - E_m = (-2,77 + 2,81)\times 10^{10} = \SI{+4e8}{J} > 0$.
L'énergie mécanique doit \textbf{augmenter} (devenir moins négative). Pour monter, le satellite doit \textbf{accélérer} (fournir une impulsion prograde) : paradoxalement, après le transfert, sa vitesse finale sur l'orbite haute sera \textbf{plus faible} que sur l'orbite basse ($v' = \sqrt{GM_T/r'} < v$), car une partie de l'énergie cinétique initiale a été convertie en énergie potentielle.
\end{enumerate}
\end{corrige}

\newpage

\coursec{Méthodes et exercices résolus}

\coursec{Référentiels d'étude et lois de Kepler}

\courssub{Choix du référentiel : géocentrique ou héliocentrique}

\begin{definition}[Référentiels d'étude du mouvement des astres]
Pour étudier le mouvement d'un satellite autour de la Terre, on utilise le \textbf{référentiel géocentrique} (origine $O$ au centre de la Terre, axes pointant vers des étoiles fixes). Pour une planète autour du Soleil, on utilise le \textbf{référentiel héliocentrique} (origine au centre du Soleil, axes fixes par rapport aux étoiles). Dans les deux cas, ces référentiels sont \textbf{supposés galiléens} (on néglige l'accélération du centre de l'astre attracteur et les perturbations des autres corps). Cette approximation est excellente pour un premier modèle.
\end{definition}

\begin{aretenir}[Vocabulaire du Bac]
\begin{itemize}
\item \cle{Référentiel géocentrique} : lié à la Terre, axes fixes par rapport aux étoiles (non tournant avec la Terre).
\item \cle{Référentiel héliocentrique} : lié au Soleil, axes fixes par rapport aux étoiles.
\item \cle{Galiléen} : référentiel dans lequel le principe d'inertie et la 2e loi de Newton s'appliquent sans forces d'inertie.
\end{itemize}
Toute application de la 2e loi de Newton doit débuter par : « Dans le référentiel géocentrique (ou héliocentrique), supposé galiléen... ».
\end{aretenir}

\courssub{Les trois lois de Kepler (énoncés)}

\begin{propriete}[Première loi de Kepler (1609) -- Loi des orbites]
Les planètes décrivent des \textbf{ellipses} dont le Soleil occupe un foyer. Pour un satellite artificiel terrestre, l'orbite est une ellipse dont la Terre occupe un foyer. \emph{Au programme, on se restreint souvent au cas particulier de l'orbite \textbf{circulaire} (excentricité nulle).}
\end{propriete}

\begin{propriete}[Deuxième loi de Kepler (1609) -- Loi des aires]
Le rayon vecteur reliant la planète au Soleil (ou le satellite à la Terre) \textbf{balaye des aires égales en temps égaux}. Conséquence : la vitesse angulaire n'est pas constante sur une ellipse ; le corps va plus vite au périastre (point le plus proche) qu'à l'apoastre (point le plus lointain). Sur une orbite circulaire, la vitesse est constante.
\end{propriete}

\begin{propriete}[Troisième loi de Kepler (1619) -- Loi des périodes]
Pour tous les satellites d'un \textbf{même astre attracteur}, le rapport du carré de la période $T$ au cube du demi-grand axe $a$ (ou du rayon $r$ pour un cercle) est une \textbf{constante} :
\[ \frac{T^2}{a^3} = \frac{4\pi^2}{GM} = \text{constante} \]
où $M$ est la masse de l'astre central et $G$ la constante de gravitation universelle. \emph{Démonstration exigible au Bac pour le cas circulaire (voir Méthode 1).}
\end{propriete}

\begin{savaistu}[Histoire des sciences : Tycho Brahe, Kepler et Newton]
C'est grâce aux observations d'une précision inégalée de l'astronome danois \textbf{Tycho Brahe} (fin XVIe siècle), notamment sur la position de Mars, que \textbf{Johannes Kepler} a pu formuler ses trois lois après des décennies de calculs acharnés. Elles sont restées empiriques jusqu'à ce qu'\textbf{Isaac Newton} (1687) les déduise de sa loi de la gravitation universelle, prouvant que la même force fait tomber la pomme et maintient la Lune en orbite.
\end{savaistu}

\coursec{Mouvement circulaire uniforme d'un satellite}

\courssub{Modélisation : bilan des forces et deuxième loi de Newton}

\begin{methode}[Établir l'équation du mouvement d'un satellite]
Pour mettre en équation le mouvement d'un satellite de masse $m$ autour d'un astre de masse $M$ :
\begin{enumerate}
\item \textbf{Choisir le référentiel} : référentiel astrocentrique (géocentrique pour un satellite terrestre, héliocentrique pour une planète), supposé galiléen. Le préciser explicitement.
\item \textbf{Faire le bilan des forces} : le satellite n'est soumis qu'à la force de gravitation $\vec{F} = -\dfrac{GMm}{r^2}\,\vec{u}_r$ où $\vec{u}_r$ est le vecteur unitaire radial \textbf{sortant} (on néglige les autres astres et les frottements atmosphériques résiduels).
\item \textbf{Appliquer la deuxième loi de Newton} : $m\vec{a} = \vec{F}$, soit $\vec{a} = -\dfrac{GM}{r^2}\,\vec{u}_r$. L'accélération est \cle{centripète} (dirigée vers le centre de l'astre).
\item \textbf{Projeter dans la base de Frenet} $(\vec{u}_t, \vec{u}_n)$ : pour un cercle, $\vec{u}_n = -\vec{u}_r$ (normal pointant vers le centre). On a $a_t = \dfrac{dv}{dt}$ et $a_n = \dfrac{v^2}{r}$. L'équation vectorielle donne $a_t = 0$ et $a_n = \dfrac{GM}{r^2}$.
\item \textbf{Conclure} : $a_t = 0 \implies v = \text{constante}$ : le mouvement circulaire est nécessairement \cle{uniforme}. La relation $a_n$ donne la vitesse orbitale.
\end{enumerate}
\textbf{Réflexe} : vérifier que la masse $m$ du satellite disparaît : le mouvement ne dépend que de $M$ et $r$.
\end{methode}

\begin{exoresolu}[Mise en équation du satellite SPOT]
Le satellite d'observation SPOT, de masse $m$, décrit une orbite circulaire de rayon $r$ autour de la Terre (masse $M_T$) dans le référentiel géocentrique.
\begin{enumerate}
\item Représenter sur un schéma la force exercée par la Terre sur le satellite et les vecteurs de la base de Frenet.
\item Établir l'expression du vecteur accélération du satellite.
\item Montrer que le mouvement circulaire du satellite est uniforme.
\end{enumerate}
\tcblower
\textbf{1. Schéma.} On place le satellite $S$ à la distance $r$ du centre $O$ de la Terre. La force $\vec{F}_{T/S}$ est attractive, dirigée de $S$ vers $O$. En base de Frenet, le vecteur normal $\vec{u}_n$ pointe vers le centre $O$ (centripète), le vecteur tangentiel $\vec{u}_t$ est porté par la tangente à la trajectoire dans le sens du mouvement.

\begin{popfigure}
\begin{tikzpicture}[scale=1.2]
% Terre
\fill[popblue] (0,0) circle (0.4);
\node[below=0.2cm] at (0,-0.4) {Terre ($O$)};
% Satellite
\coordinate (S) at (2.5,1.8);
\fill[poporange] (S) circle (0.08) node[above right] {$S$};
% Orbite
\draw[popline, dashed] (0,0) circle (3.2);
% Rayon
\draw[popmuted] (0,0) -- (S) node[midway, above left] {$r$};
% Force gravitation (vers le centre)
\draw[-{Stealth[length=3mm]}, poporange, very thick] (S) -- ($(0,0)!0.5!(S)$) node[midway, below right] {$\vec{F}_{T/S}$};
% Vecteur normal (vers le centre)
\draw[-{Stealth[length=3mm]}, popgreen, thick] (S) -- ($(0,0)!0.35!(S)$) node[midway, below left] {$\vec{u}_n$};
% Vecteur tangent
\draw[-{Stealth[length=3mm]}, poppurple, thick] (S) -- ++(-0.9,-1.25) node[below] {$\vec{u}_t$};
% Base radiale sortante (optionnel pour clarification)
\draw[-{Stealth[length=2mm]}, popgray, thin] (S) -- ++(0.6,0.4) node[above right] {$\vec{u}_r$};
\end{tikzpicture}
\legende{Satellite SPOT en orbite circulaire : forces et base de Frenet. $\vec{u}_n = -\vec{u}_r$.}
\end{popfigure}

\textbf{2. Vecteur accélération.} Système étudié : le satellite de masse $m$. Référentiel géocentrique, supposé galiléen. Bilan des forces : seule la force gravitationnelle terrestre s'exerce (action du Soleil, de la Lune et frottements négligés) :
\[ \vec{F}_{T/S} = -\frac{G M_T m}{r^2}\,\vec{u}_r \quad \text{($\vec{u}_r$ radial sortant)}. \]
La deuxième loi de Newton s'écrit $m\vec{a} = \vec{F}_{T/S}$, d'où :
\[ \vec{a} = -\frac{G M_T}{r^2}\,\vec{u}_r. \]
En base de Frenet, $\vec{u}_n = -\vec{u}_r$ (normal orienté vers le centre), donc :
\[ \boxed{\vec{a} = \frac{G M_T}{r^2}\,\vec{u}_n} \]
L'accélération est centripète, de norme $a = \dfrac{GM_T}{r^2}$. Elle ne dépend pas de la masse $m$ du satellite.

\textbf{3. Uniformité du mouvement.} Dans la base de Frenet, pour un mouvement circulaire de rayon $r$ :
\[ \vec{a} = \frac{dv}{dt}\,\vec{u}_t + \frac{v^2}{r}\,\vec{u}_n. \]
Par identification avec l'expression précédente ($\vec{a} = \frac{GM_T}{r^2}\vec{u}_n$), la composante tangentielle est nulle :
\[ a_t = \frac{dv}{dt} = 0 \quad \Longrightarrow \quad v = \text{constante}. \]
\textbf{Conclusion} : tout satellite soumis à la seule force de gravitation et décrivant une orbite circulaire a un mouvement circulaire \textbf{uniforme}. C'est une conséquence directe du caractère radial (central) de la force de gravitation.
\end{exoresolu}

\courssub{Vitesse orbitale et période de révolution}

\begin{propriete}[Vitesse et période sur une orbite circulaire]
Pour un satellite de masse $m$ sur une orbite circulaire de rayon $r$ autour d'un astre de masse $M$ :
\begin{align*}
\text{Vitesse orbitale :} \quad & v = \sqrt{\frac{GM}{r}} \quad (r = R + h) \\
\text{Période de révolution :} \quad & T = \frac{2\pi r}{v} = 2\pi\sqrt{\frac{r^3}{GM}}
\end{align*}
\textbf{Analyse dimensionnelle} : $[GM] = \mathrm{m^3\,s^{-2}}$, donc $\sqrt{r^3/(GM)}$ est homogène à des secondes. \emph{Plus l'orbite est haute ($r$ grand), plus le satellite est lent ($v$ petit) et plus sa période $T$ est grande.}
\end{propriete}

\begin{methode}[Calculer la vitesse et la période d'un satellite]
\begin{enumerate}
\item Identifier $M$ (masse de l'astre central) et $r = R + h$ (distance au centre, en mètres).
\item Appliquer $v = \sqrt{GM/r}$.
\item Calculer $T = 2\pi r / v$ ou directement $T = 2\pi\sqrt{r^3/(GM)}$.
\item Convertir le résultat dans l'unité demandée (s, min, h, jours).
\item \textbf{Contrôle d'ordre de grandeur} : $v_{\text{orbite basse}} \sim \SI{8}{km/s}$, $T \sim \SI{90}{min}$ ; $v_{\text{geo}} \sim \SI{3}{km/s}$, $T = \SI{23h56min}{}$.
\end{enumerate}
\end{methode}

\begin{exoresolu}[Vitesse et période de la Station Spatiale Internationale (ISS)]
L'ISS évolue sur une orbite circulaire à l'altitude $h = \SI{400}{km}$. Données : $G = \num{6,67e-11}\ \mathrm{m^3\,kg^{-1}\,s^{-2}}$, $M_T = \num{5,97e24}\ \mathrm{kg}$, $R_T = \SI{6370}{km}$.
\begin{enumerate}
\item Calculer la vitesse de l'ISS.
\item En déduire sa période de révolution en minutes.
\end{enumerate}
\tcblower
\textbf{1. Vitesse orbitale.} Rayon de l'orbite :
\[ r = R_T + h = 6370 + 400 = \SI{6770}{km} = \num{6,770e6}\ \mathrm{m}. \]
Dans le référentiel géocentrique galiléen, projection de la 2e loi de Newton sur $\vec{u}_n$ :
\[ \frac{v^2}{r} = \frac{GM_T}{r^2} \quad \Longrightarrow \quad v = \sqrt{\frac{GM_T}{r}}. \]
Application numérique :
\[ v = \sqrt{\frac{\num{6,67e-11} \times \num{5,97e24}}{\num{6,770e6}}} = \sqrt{\frac{\num{3,982e14}}{\num{6,770e6}}} = \sqrt{\num{5,88e7}} \approx \SI{7,67e3}{m.s^{-1}}. \]
Soit $v \approx \SI{7,7}{km.s^{-1}}$ ($\approx \SI{28000}{km.h^{-1}}$).

\textbf{2. Période de révolution.}
\[ T = \frac{2\pi r}{v} = \frac{2\pi \times \num{6,770e6}}{\num{7,67e3}} \approx \SI{5,54e3}{s}. \]
\[ T \approx \frac{5540}{60} \approx \SI{92}{min} \quad (\text{1 h 32 min}). \]
\textbf{Conclusion} : l'ISS effectue environ 16 révolutions par jour ; ses occupants voient 16 levers de Soleil quotidiens.
\end{exoresolu}

\courssub{Exploiter la troisième loi de Kepler}

\begin{propriete}[Troisième loi de Kepler -- Forme newtonienne (cas circulaire)]
Pour un satellite en orbite circulaire de rayon $r$ autour d'un astre de masse $M$ :
\[ \frac{T^2}{r^3} = \frac{4\pi^2}{GM} \]
Cette relation permet de \textbf{determiner la masse $M$ de l'astre central} à partir du mouvement d'un de ses satellites : $M = \dfrac{4\pi^2 r^3}{G T^2}$.
\end{propriete}

\begin{methode}[Utiliser la troisième loi de Kepler]
\begin{enumerate}
\item \textbf{Énoncé} : pour tous les satellites d'un même astre attracteur, $\dfrac{T^2}{r^3} = \dfrac{4\pi^2}{GM} = \text{constante}$ (pour une ellipse, $r$ est remplacé par le demi-grand axe $a$).
\item \textbf{Identifier l'astre central} : la constante est la même uniquement pour des corps gravitant autour du \emph{même} astre.
\item \textbf{Deux usages types} :
  \begin{itemize}
  \item Comparaison de deux satellites : $\dfrac{T_1^2}{r_1^3} = \dfrac{T_2^2}{r_2^3}$ (permet de trouver une inconnue sans $G$ ni $M$).
  \item Détermination de la masse de l'astre : $M = \dfrac{4\pi^2 r^3}{G T^2}$.
  \end{itemize}
\item Convertir toutes les grandeurs en unités SI (m, s, kg) avant calcul, sauf pour une comparaison de rapports (unités identiques pour les deux corps).
\end{enumerate}
\end{methode}

\begin{exoresolu}[La masse de la Terre grâce à la Lune]
La Lune décrit une orbite quasi circulaire de rayon moyen $r = \num{3,84e8}\ \mathrm{m}$ en $T = \num{27,3}$ jours. $G = \num{6,67e-11}\ \mathrm{m^3\,kg^{-1}\,s^{-2}}$.
\begin{enumerate}
\item Exprimer la masse $M_T$ de la Terre en fonction de $G$, $r$ et $T$.
\item Calculer $M_T$.
\end{enumerate}
\tcblower
\textbf{1. Expression.} Référentiel géocentrique galiléen. Troisième loi de Kepler :
\[ \frac{T^2}{r^3} = \frac{4\pi^2}{G M_T} \quad \Longrightarrow \quad \boxed{M_T = \frac{4\pi^2\, r^3}{G\, T^2}} \]
Vérification dimensionnelle : $\left[\dfrac{r^3}{GT^2}\right] = \dfrac{\mathrm{m^3}}{\mathrm{m^3\,kg^{-1}\,s^{-2}} \times \mathrm{s^2}} = \mathrm{kg}$.

\textbf{2. Application numérique.}
\[ T = 27,3 \times 24 \times 3600 = \num{2,359e6}\ \mathrm{s}. \]
\begin{align*}
r^3 &= (\num{3,84e8})^3 = \num{5,66e25}\ \mathrm{m^3}, \\
T^2 &= (\num{2,359e6})^2 = \num{5,56e12}\ \mathrm{s^2}.
\end{align*}
\[ M_T = \frac{4\pi^2 \times \num{5,66e25}}{\num{6,67e-11} \times \num{5,56e12}} = \frac{\num{2,233e27}}{\num{3,71e2}} \approx \num{6,0e24}\ \mathrm{kg}. \]
\textbf{Conclusion} : l'observation de la Lune permet de « peser » la Terre : $M_T \approx \num{6,0e24}\ \mathrm{kg}$ (valeur admise $\num{5,97e24}\ \mathrm{kg}$). Même méthode pour la masse du Soleil via les planètes.
\end{exoresolu}

\begin{methode}[Comparer deux satellites sans connaître $G$ ni $M$]
\begin{enumerate}
\item Vérifier que les deux corps gravitent autour du \textbf{même astre attracteur}.
\item Écrire $\dfrac{T_1^2}{r_1^3} = \dfrac{T_2^2}{r_2^3}$.
\item Isoler l'inconnue : $T_2 = T_1\left(\dfrac{r_2}{r_1}\right)^{3/2}$ ou $r_2 = r_1\left(\dfrac{T_2}{T_1}\right)^{2/3}$.
\item Les rapports étant sans dimension, on peut garder des unités pratiques (km, jours) si elles sont \emph{identiques} pour les deux corps.
\end{enumerate}
\textbf{Réflexe} : si $r_2 > r_1$, alors $T_2 > T_1$ et $v_2 < v_1$ : contrôler le sens physique.
\end{methode}

\begin{exoresolu}[Orbite de transfert vers l'orbite géostationnaire]
Un satellite est sur une orbite basse de stationnement $r_1 = \SI{6600}{km}$ (période $T_1$). Orbite géostationnaire : $r_2 = \SI{42160}{km}$, $T_2 = \SI{86164}{s}$.
\begin{enumerate}
\item Calculer $T_1$ sans utiliser $G$ ni $M_T$.
\item Exprimer et calculer le rapport des vitesses $\dfrac{v_1}{v_2}$.
\end{enumerate}
\tcblower
\textbf{1. Période sur l'orbite basse.}
\[ \frac{T_1^2}{r_1^3} = \frac{T_2^2}{r_2^3} \quad \Longrightarrow \quad T_1 = T_2\left(\frac{r_1}{r_2}\right)^{3/2}. \]
\[ \frac{r_1}{r_2} = \frac{6600}{42160} = \num{0,1565} \quad ; \quad (\num{0,1565})^{3/2} = \num{0,1565} \times \sqrt{\num{0,1565}} = \num{0,1565} \times \num{0,3956} = \num{0,0619}. \]
\[ T_1 = 86164 \times \num{0,0619} \approx \SI{5,33e3}{s} \approx \SI{89}{min}. \]
Contrôle : $r_1 < r_2 \implies T_1 < T_2$. Cohérent.

\textbf{2. Rapport des vitesses.}
\[ v = \frac{2\pi r}{T} \implies \frac{v_1}{v_2} = \frac{r_1}{r_2}\frac{T_2}{T_1} = \frac{r_1}{r_2}\left(\frac{r_2}{r_1}\right)^{3/2} = \sqrt{\frac{r_2}{r_1}}. \]
(On retrouve $v \propto 1/\sqrt{r}$.)
\[ \frac{v_1}{v_2} = \sqrt{\frac{42160}{6600}} = \sqrt{\num{6,39}} \approx \num{2,5}. \]
\textbf{Conclusion} : sur l'orbite basse, le satellite tourne en $\approx \SI{89}{min}$ et va 2,5 fois plus vite qu'en orbite géostationnaire. Plus on monte, plus on ralentit.
\end{exoresolu}

\coursec{Le satellite géostationnaire}

\courssub{Conditions et altitude}

\begin{definition}[Satellite géostationnaire]
Un satellite artificiel est \textbf{géostationnaire} s'il paraît immobile pour un observateur au sol. Cela impose trois conditions simultanées :
\begin{enumerate}
\item Orbite \textbf{circulaire} contenue dans le \textbf{plan équatorial}.
\item Mouvement dans le \textbf{même sens} que la rotation de la Terre (d'ouest en est).
\item Période de révolution égale à la période de rotation propre de la Terre : le \textbf{jour sidéral} $T_{\text{sid}} = \SI{86164}{s}$ (\SI{23}{h} \SI{56}{min} \SI{4}{s}), \textbf{et non} \SI{24}{h} (jour solaire moyen).
\end{enumerate}
\end{definition}

\begin{methode}[Déterminer l'altitude d'un satellite géostationnaire]
\begin{enumerate}
\item Écrire la 3e loi de Kepler avec $T = T_{\text{sid}}$ : $\dfrac{T_{\text{sid}}^2}{r^3} = \dfrac{4\pi^2}{GM_T}$.
\item Isoler le rayon orbital : $r = \left(\dfrac{GM_T\,T_{\text{sid}}^2}{4\pi^2}\right)^{1/3}$.
\item Calculer l'\cle{altitude} : $h = r - R_T$. \textbf{Ne jamais confondre $r$ et $h$.}
\item \textbf{Valeurs de référence à connaître} : $r \approx \SI{42160}{km}$, $h \approx \SI{35800}{km}$, $v \approx \SI{3,07}{km.s^{-1}}$.
\end{enumerate}
\end{methode}

\begin{exoresolu}[Un satellite de télécommunications au-dessus du Cameroun]
Pour la diffusion CRTV et la téléphonie (MTN, Orange) en Afrique centrale, on place un satellite géostationnaire à la verticale de l'équateur. Données : $G = \num{6,67e-11}$, $M_T = \num{5,97e24}\ \mathrm{kg}$, $R_T = \SI{6370}{km}$, $T_{\text{sid}} = \SI{86164}{s}$.
\begin{enumerate}
\item Rappeler les conditions d'un satellite géostationnaire.
\item Calculer le rayon $r$ et l'altitude $h$.
\item Calculer sa vitesse orbitale.
\end{enumerate}
\tcblower
\textbf{1. Conditions.} (Voir définition ci-dessus) : orbite circulaire équatoriale, sens direct, $T = T_{\text{sid}} = \SI{86164}{s}$.

\textbf{2. Rayon et altitude.}
\[ r^3 = \frac{GM_T\,T_{\text{sid}}^2}{4\pi^2} = \frac{\num{6,67e-11} \times \num{5,97e24} \times (\num{8,6164e4})^2}{4\pi^2} \]
\[ GM_T = \num{3,982e14}\ \mathrm{m^3.s^{-2}} \quad ; \quad T_{\text{sid}}^2 = \num{7,424e9}\ \mathrm{s^2} \]
\[ r^3 = \frac{\num{3,982e14} \times \num{7,424e9}}{39,48} = \frac{\num{2,957e24}}{39,48} = \num{7,49e22}\ \mathrm{m^3} \]
\[ r = (\num{7,49e22})^{1/3} \approx \num{4,216e7}\ \mathrm{m} = \SI{42160}{km} \]
\[ h = r - R_T = 42160 - 6370 \approx \SI{35790}{km} \approx \SI{35800}{km}. \]

\textbf{3. Vitesse orbitale.}
\[ v = \frac{2\pi r}{T_{\text{sid}}} = \frac{2\pi \times \num{4,216e7}}{\num{8,6164e4}} \approx \SI{3,07e3}{m.s^{-1}} \approx \SI{3,1}{km.s^{-1}}. \]
\textbf{Conclusion} : le satellite doit être à $\approx \SI{35800}{km}$ d'altitude, sur l'équateur, à $\SI{3,1}{km.s^{-1}}$.
\end{exoresolu}

\begin{savaistu}[Au Cameroun : télécommunications et satellites]
Le Cameroun utilise intensivement les satellites géostationnaires pour la télévision (CRTV via \textbf{Eutelsat} ou \textbf{Intelsat}), la téléphonie mobile (backhaul pour \textbf{MTN}, \textbf{Orange} via \textbf{RASCOM} ou satellites privés) et l'internet rural (VSAT). L'altitude de \SI{35800}{km} implique un délai de propagation aller-retour d'environ \SI{240}{ms}, perceptible en téléphonie. Les nouveaux projets (Starlink, OneWeb) utilisent des constellations en orbite basse (\SI{550}{km}) pour réduire cette latence.
\end{savaistu}

\coursec{Énergie mécanique d'un satellite (complément Bac)}

\courssub{Énergie mécanique sur orbite circulaire}

\begin{propriete}[Énergies d'un satellite sur orbite circulaire]
Pour un satellite de masse $m$ en orbite circulaire de rayon $r$ autour d'un astre de masse $M$ (référentiel galiléen, énergie potentielle nulle à l'infini) :
\begin{align*}
\text{Énergie cinétique :} \quad & E_c = \frac{1}{2}mv^2 = \frac{GMm}{2r} > 0 \\
\text{Énergie potentielle :} \quad & E_p = -\frac{GMm}{r} < 0 \\
\text{Énergie mécanique :} \quad & E_m = E_c + E_p = -\frac{GMm}{2r} < 0
\end{align*}
\textbf{Relations utiles (théorème du viriel pour $1/r^2$) :} $E_c = -E_m$ et $E_p = 2E_m$. L'énergie mécanique négative caractérise un \textbf{état lié} : le satellite ne peut pas s'échapper à l'infini sans apport d'énergie.
\end{propriete}

\begin{experience}[Démonstration guidée : relations entre $E_c$, $E_p$, $E_m$]
\textbf{Objectif :} Établir $E_c = -E_m$ et $E_p = 2E_m$ à partir de la dynamique.
\begin{enumerate}
\item Écrire la 2e loi de Newton projetée sur $\vec{u}_n$ : $m\frac{v^2}{r} = \frac{GMm}{r^2}$.
\item Multiplier par $r/2$ : $\frac{1}{2}mv^2 = \frac{GMm}{2r}$. C'est $E_c$.
\item Rappeler $E_p = -\frac{GMm}{r} = -2 \times \frac{GMm}{2r} = -2E_c$.
\item Calculer $E_m = E_c + E_p = E_c - 2E_c = -E_c$.
\item Conclure : $E_c = -E_m > 0$, $E_p = 2E_m < 0$, $E_m < 0$.
\end{enumerate}
\end{experience}

\courssub{Vitesse de libération}

\begin{propriete}[Vitesse de libération (ou vitesse de fuite)]
C'est la vitesse minimale à communiquer à un corps au départ de la surface d'un astre (rayon $R$, masse $M$) pour qu'il s'en échappe définitivement (vitesse nulle à l'infini), en négligeant l'atmosphère.
Par conservation de l'énergie mécanique entre la surface ($r=R$, $v=v_\ell$) et l'infini ($r\to\infty$, $v=0$, $E_p=0$) :
\[ \frac{1}{2}mv_\ell^2 - \frac{GMm}{R} = 0 \quad \Longrightarrow \quad \boxed{v_\ell = \sqrt{\frac{2GM}{R}}} \]
Relation avec la vitesse de satellisation rasante $v_1 = \sqrt{GM/R}$ : $v_\ell = \sqrt{2}\, v_1 \approx 1,41\, v_1$.
\end{propriete}

\begin{methode}[Exploiter l'énergie mécanique et la vitesse de libération]
\begin{enumerate}
\item \textbf{Énergie sur orbite} : $E_m = -\dfrac{GMm}{2r}$. Utiliser $E_c = -E_m$, $E_p = 2E_m$ pour aller vite.
\item \textbf{Satellisation} : énergie à fournir pour passer du sol (immobile) à l'orbite $r$ : $\Delta E = E_m(r) - E_m(R) = GMm\left(\dfrac{1}{R} - \dfrac{1}{2r}\right)$.
\item \textbf{Vitesse de libération} : $v_\ell = \sqrt{\dfrac{2GM}{R}}$. Ne pas confondre avec $v_1 = \sqrt{GM/R}$.
\item \textbf{Vérification} : $E_m < 0$ pour un satellite (état lié) ; $E_m \ge 0$ pour une trajectoire non liée (parabole, hyperbole).
\end{enumerate}
\end{methode}

\begin{exoresolu}[Énergie d'un satellite et vitesse de libération terrestre]
Satellite $m = \SI{800}{kg}$, orbite circulaire $h = \SI{600}{km}$. $G = \num{6,67e-11}$, $M_T = \num{5,97e24}\ \mathrm{kg}$, $R_T = \num{6,37e6}\ \mathrm{m}$.
\begin{enumerate}
\item Calculer $E_c$, $E_p$, $E_m$ du satellite.
\item Calculer la vitesse de libération $v_\ell$ depuis la surface.
\end{enumerate}
\tcblower
\textbf{1. Énergies du satellite.}
\[ r = R_T + h = \num{6,37e6} + \num{0,600e6} = \num{6,97e6}\ \mathrm{m}. \]
\[ GM_T m = \num{6,67e-11} \times \num{5,97e24} \times 800 = \num{3,186e17}\ \mathrm{J.m} \]
\[ E_c = \frac{GM_T m}{2r} = \frac{\num{3,186e17}}{2 \times \num{6,97e6}} = \frac{\num{3,186e17}}{\num{1,394e7}} \approx \num{2,29e10}\ \mathrm{J}. \]
\[ E_p = -2E_c \approx \num{-4,57e10}\ \mathrm{J}. \]
\[ E_m = E_c + E_p = -E_c \approx \num{-2,29e10}\ \mathrm{J}. \]
\textbf{Interprétation} : $E_m < 0$, le satellite est lié à la Terre.

\textbf{2. Vitesse de libération.}
\[ v_\ell = \sqrt{\frac{2GM_T}{R_T}} = \sqrt{\frac{2 \times \num{6,67e-11} \times \num{5,97e24}}{\num{6,37e6}}} = \sqrt{\num{1,250e8}} \approx \SI{1,12e4}{m.s^{-1}}. \]
\textbf{Conclusion} : $v_\ell \approx \SI{11,2}{km.s^{-1}}$ (environ $\sqrt{2} \times \SI{7,9}{km.s^{-1}}$). C'est la vitesse nécessaire pour quitter le champ gravitationnel terrestre (sans atmosphère).
\end{exoresolu}

\coursec{Synthèse et méthodes de résolution}

\courssub{Tableau récapitulatif des formules clés}

\begin{center}
\begin{tabularx}{\linewidth}{|l|X|X|}\hline
\rowcolor{popblueL}
\textbf{Grandeurs} & \textbf{Expression (orbite circulaire, rayon $r$)} & \textbf{Unité SI} \\ \hline
Vitesse orbitale & $v = \sqrt{\dfrac{GM}{r}}$ & $\mathrm{m.s^{-1}}$ \\ \hline
Période & $T = 2\pi\sqrt{\dfrac{r^3}{GM}} = \dfrac{2\pi r}{v}$ & $\mathrm{s}$ \\ \hline
3e loi de Kepler & $\dfrac{T^2}{r^3} = \dfrac{4\pi^2}{GM}$ & $\mathrm{s^2.m^{-3}}$ \\ \hline
Masse de l'astre & $M = \dfrac{4\pi^2 r^3}{G T^2}$ & $\mathrm{kg}$ \\ \hline
Énergie cinétique & $E_c = \dfrac{GMm}{2r}$ & $\mathrm{J}$ \\ \hline
Énergie potentielle & $E_p = -\dfrac{GMm}{r}$ & $\mathrm{J}$ \\ \hline
Énergie mécanique & $E_m = -\dfrac{GMm}{2r}$ & $\mathrm{J}$ \\ \hline
Vitesse de libération & $v_\ell = \sqrt{\dfrac{2GM}{R}}$ & $\mathrm{m.s^{-1}}$ \\ \hline
Géostationnaire & $T = T_{\text{sid}} = \SI{86164}{s}$, $h \approx \SI{35800}{km}$ & -- \\ \hline
\end{tabularx}
\end{center}

\courssub{Guide de résolution type Bac}

\begin{methode}[Résoudre un problème de mécanique spatiale (check-list)]
\begin{enumerate}
\item \textbf{Lire et annoter} : identifier l'astre central ($M$, $R$), le satellite ($m$), l'orbite ($r$ ou $h$, circulaire/elliptique), le référentiel demandé.
\item \textbf{Poser le référentiel} : écrire explicitement « Référentiel géocentrique (ou héliocentrique) supposé galiléen ».
\item \textbf{Bilan des forces} : dessiner le schéma (force gravitationnelle seule, centripète). Nommer les vecteurs ($\vec{F}$, $\vec{u}_n$, $\vec{u}_t$).
\item \textbf{Équation du mouvement} : projeter la 2e loi de Newton sur $\vec{u}_n$ : $m v^2/r = GMm/r^2$.
\item \textbf{Exploiter} :
  \begin{itemize}
  \item Vitesse / Période : formules directes ou 3e loi de Kepler.
  \item Masse de l'astre : $M = 4\pi^2 r^3/(GT^2)$.
  \item Comparaison deux satellites : $T_1^2/r_1^3 = T_2^2/r_2^3$.
  \item Géostationnaire : imposer $T = \SI{86164}{s}$, plan équatorial, sens direct.
  \item Énergie : $E_m = -GMm/(2r)$, $v_\ell = \sqrt{2GM/R}$.
  \end{itemize}
\item \textbf{Calculer} : convertir en SI (m, kg, s), calculer, arrondir au nombre de chiffres significatifs des données (généralement 2 ou 3), \textbf{écrire l'unité}.
\item \textbf{Contrôler} : ordre de grandeur, signe des énergies ($E_m<0$ pour lié), cohérence physique ($r$ croissant $\implies v$ décroissant, $T$ croissant).
\end{enumerate}
\end{methode}

\courssub{Erreurs fréquentes et pièges à éviter}

\begin{attention}[Les 5 erreurs qui coûtent des points au Bac]
\begin{enumerate}
\item \textbf{Confondre le rayon de l'orbite $r$ et l'altitude $h$.} Toutes les formules ($v = \sqrt{GM/r}$, $T = 2\pi\sqrt{r^3/GM}$, $E_m = -GMm/2r$) s'écrivent avec $r = R + h$, distance au \emph{centre} de l'astre. Utiliser $h$ à la place de $r$ est l'erreur la plus fréquente ; conclure « l'altitude est $\SI{42160}{km}$ » pour un géostationnaire l'est tout autant : il faut retrancher $R_T$ pour obtenir $h \approx \SI{35800}{km}$.
\item \textbf{Oublier de préciser le référentiel et le caractère galiléen.} Toute application de la deuxième loi de Newton doit commencer par : « Dans le référentiel géocentrique (ou héliocentrique), supposé galiléen... ». Sans cette phrase, le raisonnement est incomplet et pénalisé.
\item \textbf{Appliquer la troisième loi de Kepler à des astres différents.} Le rapport $T^2/r^3$ n'est constant que pour des satellites du \emph{même} astre attracteur. Comparer directement la Lune (satellite de la Terre) et Vénus (planète du Soleil) est faux : les constantes $\dfrac{4\pi^2}{GM}$ diffèrent.
\item \textbf{Prendre $T = \SI{24}{h} = \SI{86400}{s}$ pour le satellite géostationnaire.} La période de révolution d'un géostationnaire est le \emph{jour sidéral} $T_{\text{sid}} = \SI{86164}{s}$ (durée de rotation de la Terre par rapport aux étoiles). L'écart de près de 4 minutes fausse le calcul de l'altitude d'environ $\SI{100}{km}$.
\item \textbf{Se tromper de signe ou de facteur 2 dans les énergies.} On a $E_p = -\dfrac{GMm}{r}$ (négative, nulle à l'infini), $E_c = +\dfrac{GMm}{2r}$ et $E_m = -\dfrac{GMm}{2r}$. Écrire $E_m = +\dfrac{GMm}{2r}$ ou oublier le signe moins de $E_p$ conduit à une énergie mécanique positive, ce qui décrirait un satellite \emph{non lié} : contradiction physique qu'un contrôle rapide permet de détecter. De même, $v_\ell = \sqrt{2GM/R}$ et non $\sqrt{GM/R}$ (qui est la vitesse de satellisation rasante).
\end{enumerate}
\end{attention}

\courssub{Exercices d'entraînement}

\begin{exercice}[Exercice 1 : Satellite d'observation terrestre]
Un satellite d'observation de masse $m = \SI{1200}{kg}$ est placé sur une orbite circulaire héliosynchrone à l'altitude $h = \SI{800}{km}$. Données : $G = \num{6,67e-11}\ \mathrm{SI}$, $M_T = \num{5,97e24}\ \mathrm{kg}$, $R_T = \SI{6370}{km}$.
\begin{enumerate}
\item Dans quel référentiel étudie-t-on son mouvement ? Ce référentiel est-il galiléen ? Justifier.
\item Faire le bilan des forces et établir l'expression de la vitesse orbitale. Calculer sa valeur numérique.
\item Calculer la période de révolution en minutes.
\item Déterminer son énergie mécanique. Le satellite est-il lié à la Terre ?
\end{enumerate}
\end{exercice}

\begin{exercice}[Exercice 2 : La lune Phobos de Mars]
Phobos, satellite naturel de Mars, décrit une orbite quasi circulaire de rayon $r = \num{9,38e6}\ \mathrm{m}$ avec une période $T = \num{7,66e3}\ \mathrm{s}$. $G = \num{6,67e-11}\ \mathrm{SI}$.
\begin{enumerate}
\item En déduire la masse $M_M$ de Mars.
\item Un second satellite, Déimos, orbite à $r' = \num{2,35e7}\ \mathrm{m}$. Sans utiliser $G$ ni $M_M$, calculer sa période $T'$ en heures.
\item Comparer les vitesses orbitales de Phobos et Déimos.
\end{enumerate}
\end{exercice}

\begin{exercice}[Exercice 3 : Mise en orbite géostationnaire par étage]
Un lanceur place un satellite de masse $m = \SI{2000}{kg}$ sur une orbite de transfert elliptique (GTO) de périgée $r_p = \SI{6600}{km}$ et d'apogée $r_a = \SI{42160}{km}$ (orbite géostationnaire). On néglige les frottements.
\begin{enumerate}
\item Rappeler les trois conditions pour qu'un satellite soit géostationnaire.
\item À l'apogée, le satellite a une vitesse $v_a$. En utilisant la conservation du moment cinétique (aire constante) et de l'énergie mécanique, exprimer $v_a$ en fonction de $G, M_T, r_p, r_a$.
\item Calculer l'apport d'énergie (en $\mathrm{J}$) nécessaire au périgée pour circulariser l'orbite à $r_a$ (allumage du moteur d'apogée).
\end{enumerate}
\end{exercice}

\courssub{Corrigés des exercices}

\begin{corrige}[Exercice 1]
\textbf{1.} Référentiel géocentrique (origine centre Terre, axes fixes vers étoiles). Supposé galiléen car accélération du centre de Terre et perturbations (Lune, Soleil) négligées au premier ordre.
\textbf{2.} Bilan : $\vec{F} = -\frac{GM_T m}{r^2}\vec{u}_r$. 2e loi de Newton : $m\frac{v^2}{r} = \frac{GM_T m}{r^2} \implies v = \sqrt{\frac{GM_T}{r}}$.
$r = 6370+800 = \SI{7170}{km} = \num{7,17e6}\ \mathrm{m}$.
$v = \sqrt{\frac{\num{3,98e14}}{\num{7,17e6}}} = \sqrt{\num{5,55e7}} \approx \SI{7,45e3}{m.s^{-1}}$.
\textbf{3.} $T = \frac{2\pi r}{v} = \frac{2\pi \times \num{7,17e6}}{\num{7,45e3}} \approx \SI{6040}{s} \approx \SI{101}{min}$.
\textbf{4.} $E_m = -\frac{GM_T m}{2r} = -\frac{\num{3,98e14} \times 1200}{2 \times \num{7,17e6}} \approx \num{-3,33e10}\ \mathrm{J}$. $E_m < 0 \implies$ satellite lié.
\end{corrige}

\begin{corrige}[Exercice 2]
\textbf{1.} $M_M = \frac{4\pi^2 r^3}{G T^2} = \frac{4\pi^2 (\num{9,38e6})^3}{\num{6,67e-11} \times (\num{7,66e3})^2}$.
$r^3 \approx \num{8,25e20}$, $T^2 \approx \num{5,87e7}$.
$M_M \approx \frac{4\pi^2 \times \num{8,25e20}}{\num{6,67e-11} \times \num{5,87e7}} \approx \frac{\num{3,26e22}}{\num{3,91e-3}} \approx \num{6,4e23}\ \mathrm{kg}$.
\textbf{2.} Même astre (Mars) $\implies \frac{T'^2}{r'^3} = \frac{T^2}{r^3} \implies T' = T \left(\frac{r'}{r}\right)^{3/2}$.
$\frac{r'}{r} = \frac{2,35e7}{9,38e6} = 2,505$. $(2,505)^{3/2} = 2,505 \times \sqrt{2,505} \approx 2,505 \times 1,583 = 3,966$.
$T' = \num{7,66e3} \times 3,966 \approx \num{3,04e4}\ \mathrm{s} \approx \SI{8,44}{h}$.
\textbf{3.} $v \propto 1/\sqrt{r}$. $r' > r \implies v' < v$. Phobos (plus proche) est plus rapide que Déimos.
$v_{\text{Phobos}} = 2\pi r/T \approx \SI{7,69e3}{m/s}$. $v_{\text{Déimos}} = v_{\text{Phobos}} \sqrt{r/r'} \approx 7690 \times \sqrt{1/2,505} \approx \SI{4,86e3}{m/s}$.
\end{corrige}

\begin{corrige}[Exercice 3]
\textbf{1.} (Voir définition) : Orbite circulaire équatoriale, sens direct, $T = T_{\text{sid}} = \SI{86164}{s}$.
\textbf{2.} Orbite elliptique : conservation moment cinétique (2e loi Kepler) : $r_p v_p = r_a v_a$.
Conservation énergie mécanique : $\frac{1}{2}mv_p^2 - \frac{GMm}{r_p} = \frac{1}{2}mv_a^2 - \frac{GMm}{r_a}$.
Isoler $v_a$ : de la 1ère, $v_p = \frac{r_a}{r_p}v_a$. Substituer dans la 2e :
$\frac{1}{2}\frac{r_a^2}{r_p^2}v_a^2 - \frac{GM}{r_p} = \frac{1}{2}v_a^2 - \frac{GM}{r_a}$.
$v_a^2 \left(\frac{r_a^2}{r_p^2} - 1\right) = 2GM\left(\frac{1}{r_p} - \frac{1}{r_a}\right) = 2GM\frac{r_a-r_p}{r_p r_a}$.
$v_a^2 \frac{(r_a-r_p)(r_a+r_p)}{r_p^2} = 2GM\frac{r_a-r_p}{r_p r_a}$.
$v_a^2 = \frac{2GM r_p}{r_a(r_a+r_p)}$.
Numérique : $r_p = \num{6,6e6}$, $r_a = \num{4,216e7}$.
$v_a = \sqrt{\frac{2 \times \num{3,98e14} \times \num{6,6e6}}{\num{4,216e7} \times \num{4,876e7}}} \approx \sqrt{\frac{\num{5,25e21}}{\num{2,056e15}}} \approx \sqrt{\num{2,55e6}} \approx \SI{1,60e3}{m/s}$.
\textbf{3.} Circularisation à $r_a$ : vitesse circulaire $v_{\text{geo}} = \sqrt{GM/r_a} \approx \SI{3,07e3}{m/s}$.
Énergie initiale (à l'apogée GTO) : $E_i = \frac{1}{2}m v_a^2 - \frac{GMm}{r_a}$.
Énergie finale (circulaire $r_a$) : $E_f = -\frac{GMm}{2r_a}$.
Apport $\Delta E = E_f - E_i = -\frac{GMm}{2r_a} - \left(\frac{1}{2}m v_a^2 - \frac{GMm}{r_a}\right) = \frac{GMm}{2r_a} - \frac{1}{2}m v_a^2$.
En utilisant $v_a^2 = \frac{2GM r_p}{r_a(r_a+r_p)}$ :
$\Delta E = \frac{GMm}{2r_a} \left[ 1 - \frac{2r_p}{r_a+r_p} \right] = \frac{GMm}{2r_a} \frac{r_a - r_p}{r_a + r_p}$.
Numérique : $\Delta E = \frac{\num{3,98e14} \times 2000}{2 \times \num{4,216e7}} \times \frac{\num{4,216e7} - \num{6,6e6}}{\num{4,216e7} + \num{6,6e6}} \approx \num{9,44e9} \times \frac{\num{3,556e7}}{\num{4,876e7}} \approx \num{9,44e9} \times 0,73 \approx \num{6,9e9}\ \mathrm{J}$.
\end{corrige}

\newpage

\coursec{Exercices}

\courssub{Je vérifie mes connaissances}

\begin{exercice}[QCM : référentiels et lois de Kepler]
Pour chaque affirmation, choisir la (ou les) bonne(s) réponse(s). Justifier brièvement.
\begin{enumerate}
\item Le référentiel géocentrique est :
\begin{enumerate}
\item lié à la surface de la Terre ;
\item lié au centre de la Terre, ses axes pointant vers des étoiles lointaines ;
\item lié au centre du Soleil.
\end{enumerate}
\item Selon la première loi de Kepler, la trajectoire d'une planète autour du Soleil est :
\begin{enumerate}
\item un cercle centré sur le Soleil ;
\item une ellipse dont le Soleil est un foyer ;
\item une ellipse dont le Soleil est le centre.
\end{enumerate}
\item Pour un satellite terrestre en orbite circulaire de rayon $r$, la vitesse est :
\begin{enumerate}
\item $v = \sqrt{\dfrac{G\,M_T}{r}}$ ;
\item $v = \sqrt{\dfrac{G\,m}{r}}$ où $m$ est la masse du satellite ;
\item $v = \sqrt{G\,M_T\,r}$.
\end{enumerate}
\item La troisième loi de Kepler s'écrit, pour les satellites d'un même astre de masse $M$ :
\begin{enumerate}
\item $\dfrac{T^2}{r^3} = \dfrac{4\pi^2}{GM}$ ;
\item $\dfrac{T^2}{r^3} = \dfrac{GM}{4\pi^2}$ ;
\item $\dfrac{T}{r^3} = \dfrac{4\pi^2}{GM}$.
\end{enumerate}
\end{enumerate}
\end{exercice}

\begin{exercice}[Vrai ou faux justifié]
Répondre par vrai ou faux en justifiant chaque réponse par une phrase ou une formule.
\begin{enumerate}
\item La vitesse d'un satellite en orbite circulaire dépend de sa masse.
\item Un satellite géostationnaire peut se trouver immobile au-dessus de la ville de Douala.
\item Plus l'altitude d'un satellite terrestre est grande, plus sa période de révolution est grande.
\item La force de gravitation exercée par la Terre sur un satellite en orbite circulaire ne travaille pas.
\item Un satellite géostationnaire tourne dans le même sens que la rotation de la Terre.
\end{enumerate}
\end{exercice}

\begin{exercice}[Définitions et conditions]
\begin{enumerate}
\item Définir les référentiels héliocentrique et géocentrique. Préciser dans lequel on étudie le mouvement de la Lune et dans lequel on étudie le mouvement de Mars.
\item Énoncer les trois lois de Kepler.
\item Donner les trois conditions qu'un satellite doit remplir pour être géostationnaire (plan de la trajectoire, sens de rotation, période).
\end{enumerate}
\end{exercice}

\begin{exercice}[Calculs directs]
On donne : $G = \num{6,67}\times 10^{-11}\ \mathrm{SI}$ ; masse de la Terre $M_T = \num{5,98}\times 10^{24}\ \mathrm{kg}$ ; rayon terrestre $R_T = \num{6,38}\times 10^{3}\ \mathrm{km}$.
\begin{enumerate}
\item Calculer la vitesse d'un satellite en orbite circulaire à l'altitude $h = \num{400}\ \mathrm{km}$ (orbite de la station spatiale internationale).
\item En déduire sa période de révolution $T$, exprimée en minutes.
\item Vérifier par analyse dimensionnelle que l'expression $v = \sqrt{GM_T/r}$ est bien homogène à une vitesse.
\end{enumerate}
\end{exercice}

\courssub{Je m'entraîne}

\begin{exercice}[Deux satellites d'un même astre]
Deux satellites A et B évoluent sur des orbites circulaires autour d'un même astre, avec $r_B = 4\,r_A$.
\begin{enumerate}
\item À l'aide de la troisième loi de Kepler, déterminer sans calculatrice le rapport $\dfrac{T_B}{T_A}$.
\item À l'aide de l'expression de la vitesse, déterminer le rapport $\dfrac{v_B}{v_A}$.
\item Conclure : comment évoluent la période et la vitesse lorsque le rayon de l'orbite augmente ?
\end{enumerate}
\end{exercice}

\begin{exercice}[La Lune et la constante de Kepler]
On donne pour la Lune : $r_L = \num{3,84}\times 10^{8}\ \mathrm{m}$, $T_L = \num{27,3}\ \mathrm{j}$ ; pour un satellite géostationnaire : $r_G = \num{4,22}\times 10^{7}\ \mathrm{m}$, $T_G = \num{23}\ \mathrm{h}\ \num{56}\ \mathrm{min}$.
\begin{enumerate}
\item Convertir les deux périodes en secondes.
\item Calculer le rapport $\dfrac{T^2}{r^3}$ pour la Lune puis pour le satellite géostationnaire. Comparer.
\item En déduire la valeur de la masse de la Terre. On donne $G = \num{6,67}\times 10^{-11}\ \mathrm{SI}$.
\item Expliquer en quoi ce résultat constitue une vérification de la loi de gravitation universelle.
\end{enumerate}
\end{exercice}

\begin{exercice}[Masse du Soleil et masse de Jupiter]
On donne $G = \num{6,67}\times 10^{-11}\ \mathrm{SI}$.
\begin{enumerate}
\item La Terre décrit autour du Soleil une orbite quasi circulaire de rayon $r = \num{1,50}\times 10^{11}\ \mathrm{m}$ en $T = \num{365,25}\ \mathrm{j}$. Établir l'expression de la masse du Soleil $M_S$ en fonction de $G$, $r$ et $T$, puis la calculer.
\item Le satellite Io de Jupiter a une période $T_{Io} = \num{1,77}\ \mathrm{j}$ et un rayon orbital $r_{Io} = \num{4,22}\times 10^{8}\ \mathrm{m}$. Calculer la masse de Jupiter $M_J$.
\item Calculer le rapport $\dfrac{M_J}{M_S}$ et commenter.
\end{enumerate}
\end{exercice}

\begin{exercice}[Satellite météorologique]
Un satellite météorologique de masse $m = \num{850}\ \mathrm{kg}$ décrit une orbite circulaire à l'altitude $h = \num{850}\ \mathrm{km}$. On donne $M_T = \num{5,98}\times 10^{24}\ \mathrm{kg}$, $R_T = \num{6,38}\times 10^{6}\ \mathrm{m}$, $G = \num{6,67}\times 10^{-11}\ \mathrm{SI}$.
\begin{enumerate}
\item Faire le bilan des forces et représenter la force sur un schéma.
\item Montrer que le mouvement du satellite est uniforme.
\item Calculer sa vitesse $v$ et sa période $T$.
\item Calculer la valeur de la force de gravitation qu'il subit et la comparer à son poids à la surface de la Terre ($g_0 = \num{9,8}\ \mathrm{N.kg^{-1}}$).
\end{enumerate}
\end{exercice}

\courssub{Je me prépare au Bac}

\begin{exercice}[Type Bac : satellite de télécommunications au service du Cameroun]
Pour acheminer les signaux de télévision et de téléphonie entre Yaoundé et les régions éloignées, on utilise un satellite de télécommunications $S$ de masse $m = \num{2,5}\times 10^{3}\ \mathrm{kg}$, assimilé à un point matériel, en orbite circulaire de rayon $r$ autour du centre $O$ de la Terre. On se place dans le référentiel géocentrique supposé galiléen.

\textbf{Données :} $G = \num{6,67}\times 10^{-11}\ \mathrm{SI}$ ; $M_T = \num{5,98}\times 10^{24}\ \mathrm{kg}$ ; $R_T = \num{6,38}\times 10^{6}\ \mathrm{m}$ ; période de rotation propre de la Terre : $T_0 = \num{86164}\ \mathrm{s}$.
\begin{enumerate}
\item Définir le référentiel géocentrique et rappeler pourquoi on peut le considérer galiléen pour cette étude.
\item Faire le bilan des forces appliquées au satellite et représenter la force sur un schéma où figurent le repère de Frenet $(S, \vec{u}_t, \vec{u}_n)$.
\item Appliquer la deuxième loi de Newton au satellite. En projetant sur $\vec{u}_t$, montrer que le mouvement est uniforme.
\item En projetant sur $\vec{u}_n$, établir l'expression de la vitesse $v = \sqrt{\dfrac{G\,M_T}{r}}$.
\item Le satellite doit être géostationnaire.
\begin{enumerate}
\item Rappeler les trois conditions à remplir.
\item Établir l'expression de la troisième loi de Kepler $\dfrac{T^2}{r^3} = \dfrac{4\pi^2}{G\,M_T}$.
\item En déduire l'expression littérale du rayon $r$ de l'orbite, puis montrer que l'altitude du satellite est unique. Calculer cette altitude $h$.
\end{enumerate}
\item Calculer la vitesse $v$ du satellite sur son orbite.
\item Calculer la vitesse $v_D$ d'un point de l'équateur terrestre (par exemple un habitant de Kribi) due à la rotation de la Terre. Comparer $v$ et $v_D$ et expliquer pourquoi le satellite paraît immobile vu du sol.
\end{enumerate}
\end{exercice}

\begin{exercice}[Type Bac : peser le Soleil et vérifier Kepler]
Un professeur de Terminale C propose à ses élèves de « peser » le Soleil sans jamais quitter la salle de classe.

\textbf{Données :} $G = \num{6,67}\times 10^{-11}\ \mathrm{SI}$ ; orbite terrestre assimilée à un cercle de rayon $r_T = \num{1,50}\times 10^{11}\ \mathrm{m}$ ; $T_T = \num{365,25}\ \mathrm{j}$ ; $1\ \mathrm{j} = \num{86400}\ \mathrm{s}$.
\begin{enumerate}
\item Dans quel référentiel étudie-t-on le mouvement de la Terre autour du Soleil ? Le définir.
\item Énoncer les trois lois de Kepler appliquées au système solaire.
\item La Terre, de masse $m$, est soumise à la seule attraction gravitationnelle du Soleil de masse $M_S$.
\begin{enumerate}
\item Écrire l'expression vectorielle de cette force.
\item Montrer que le mouvement circulaire de la Terre est uniforme.
\item Établir l'expression de sa vitesse $v$ en fonction de $G$, $M_S$ et $r_T$.
\end{enumerate}
\item En déduire que $M_S = \dfrac{4\pi^2\,r_T^3}{G\,T_T^2}$ et calculer $M_S$.
\item La constante de Kepler du système solaire vaut $K = \num{2,97}\times 10^{-19}\ \mathrm{s^2.m^{-3}}$. Vérifier la cohérence avec le résultat de la question 4.
\item Mars a une période de révolution $T_M = \num{687}\ \mathrm{j}$. En exploitant la troisième loi de Kepler, déterminer le rayon $r_M$ de son orbite (supposée circulaire) en fonction de $r_T$, $T_T$ et $T_M$, puis calculer sa valeur.
\end{enumerate}
\end{exercice}

\begin{exercice}[Type Bac : mise sur orbite géostationnaire — étude énergétique]
Une agence spatiale souhaite transférer un satellite de masse $m = \num{1,2}\times 10^{3}\ \mathrm{kg}$ depuis une orbite basse circulaire de rayon $r_1 = R_T + \num{300}\ \mathrm{km}$ jusqu'à l'orbite géostationnaire de rayon $r_2 = \num{4,22}\times 10^{7}\ \mathrm{m}$.

\textbf{Données :} $G = \num{6,67}\times 10^{-11}\ \mathrm{SI}$ ; $M_T = \num{5,98}\times 10^{24}\ \mathrm{kg}$ ; $R_T = \num{6,38}\times 10^{6}\ \mathrm{m}$. On rappelle que l'énergie potentielle de gravitation du satellite est $E_p = -\dfrac{G\,M_T\,m}{r}$ (nulle à l'infini).
\begin{enumerate}
\item Pour une orbite circulaire de rayon $r$, établir l'expression de l'énergie cinétique $E_c = \dfrac{G\,M_T\,m}{2r}$ du satellite.
\item En déduire que l'énergie mécanique du satellite s'écrit $E_m = -\dfrac{G\,M_T\,m}{2r}$. Commenter son signe.
\item Calculer les énergies mécaniques $E_{m1}$ et $E_{m2}$ du satellite sur les deux orbites.
\item En déduire l'énergie minimale $W$ que les propulseurs doivent fournir au satellite pour réaliser le transfert. Exprimer le résultat en joules puis en kilowattheures ($1\ \mathrm{kWh} = \num{3,6}\times 10^{6}\ \mathrm{J}$).
\item Calculer les vitesses $v_1$ et $v_2$ du satellite sur les deux orbites. Le satellite va-t-il plus vite sur l'orbite haute ? Justifier alors que l'énergie fournie sert surtout à « éloigner » le satellite de la Terre.
\item On appelle vitesse de libération la vitesse minimale $v_l$ qu'il faut communiquer à un corps à la surface de la Terre pour qu'il échappe définitivement à son attraction. En écrivant que son énergie mécanique totale est nulle à l'infini, établir $v_l = \sqrt{\dfrac{2\,G\,M_T}{R_T}}$ et calculer sa valeur. Comparer à $v_1$.
\end{enumerate}
\end{exercice}

\newpage

\coursec{Corrigés des exercices}

\courssub{Je vérifie mes connaissances}

\begin{corrige}[Exercice 1 — QCM : référentiels et lois de Kepler]
\begin{enumerate}
\item \textbf{Réponse b.} Le référentiel géocentrique a pour origine le centre de la Terre et ses axes pointent vers trois étoiles lointaines considérées comme fixes. Il ne tourne pas avec la Terre (contrairement au référentiel terrestre, réponse a) et n'est pas centré sur le Soleil (réponse c : c'est le référentiel héliocentrique).
\item \textbf{Réponse b.} Première loi de Kepler (loi des ellipses) : la trajectoire d'une planète est une ellipse dont le Soleil occupe \emph{l'un des foyers}, et non le centre.
\item \textbf{Réponse a.} L'application de la deuxième loi de Newton en projection normale donne $v = \sqrt{\dfrac{G\,M_T}{r}}$ : la vitesse ne dépend que de la masse $M_T$ de l'astre attracteur et du rayon $r$, jamais de la masse $m$ du satellite. La réponse c n'est pas homogène : $G\,M_T\,r$ a pour unité $\mathrm{m^4.kg^{-1}.s^{-2}}$, dont la racine carrée n'est pas une vitesse.
\item \textbf{Réponse a.} Troisième loi de Kepler : $\dfrac{T^2}{r^3} = \dfrac{4\pi^2}{GM}$. Vérification dimensionnelle : $[G] = \mathrm{m^3.kg^{-1}.s^{-2}}$, donc $\left[\dfrac{4\pi^2}{GM}\right] = \dfrac{\mathrm{m^3.kg^{-1}.s^{-2}}}{\mathrm{kg}} \times \mathrm{kg^0} = \mathrm{s^2.m^{-3}}$... reprenons proprement : $\dfrac{1}{GM}$ s'exprime en $\dfrac{\mathrm{kg}}{\mathrm{m^3.s^{-2}}} = \mathrm{s^2.m^{-3}}$, ce qui correspond bien à l'unité de $\dfrac{T^2}{r^3}$.
\end{enumerate}
\end{corrige}

\begin{corrige}[Exercice 2 — Vrai ou faux justifié]
\begin{enumerate}
\item \textbf{Faux.} $v = \sqrt{\dfrac{G\,M_T}{r}}$ : la masse $m$ du satellite n'intervient pas ; seule compte la masse de l'astre attracteur. Tous les satellites sur la même orbite ont la même vitesse.
\item \textbf{Faux.} Un satellite géostationnaire évolue obligatoirement dans le plan équatorial. Or Douala (latitude $\approx 4^{\circ}$ N) n'est pas sur l'équateur : le satellite ne peut pas rester à sa verticale. Vu de Douala, il paraîtrait fixe mais décalé vers le sud, au-dessus d'un point de l'équateur.
\item \textbf{Vrai.} D'après la troisième loi de Kepler, $T^2 = \dfrac{4\pi^2}{GM_T}\,r^3$ : $T$ est une fonction croissante de $r = R_T + h$, donc de l'altitude $h$.
\item \textbf{Vrai.} Sur une orbite circulaire, la force de gravitation est radiale (centripète) alors que le déplacement élémentaire est tangentiel : $\vec{F} \perp d\vec{r}$, donc la puissance $\vec{F}\cdot\vec{v} = 0$ et le travail est nul. C'est cohérent avec le fait que la vitesse (et donc l'énergie cinétique) reste constante.
\item \textbf{Vrai.} Le satellite géostationnaire tourne dans le sens de rotation propre de la Terre (d'ouest en est), avec la même période $T = T_0$ : c'est la condition pour qu'il paraisse immobile vu du sol.
\end{enumerate}
\end{corrige}

\begin{corrige}[Exercice 3 — Définitions et conditions]
\begin{enumerate}
\item Le \textbf{référentiel héliocentrique} a pour origine le centre du Soleil et ses axes pointent vers trois étoiles lointaines fixes ; il est galiléen avec une excellente approximation. Le \textbf{référentiel géocentrique} a pour origine le centre de la Terre et des axes pointant vers des étoiles lointaines ; il est en translation quasi circulaire autour du Soleil et peut être considéré galiléen pour des durées courtes devant l'année. On étudie le mouvement de la \textbf{Lune} dans le référentiel géocentrique (elle tourne autour de la Terre) et celui de \textbf{Mars} dans le référentiel héliocentrique (elle tourne autour du Soleil).
\item \textbf{Lois de Kepler :}
\begin{itemize}
\item \emph{Loi des orbites} : chaque planète décrit une ellipse dont le Soleil occupe l'un des foyers.
\item \emph{Loi des aires} : le segment reliant le Soleil à la planète balaie des aires égales pendant des durées égales.
\item \emph{Loi des périodes} : le rapport $\dfrac{T^2}{a^3}$ (où $a$ est le demi-grand axe) est le même pour toutes les planètes d'un même astre attracteur : $\dfrac{T^2}{a^3} = \dfrac{4\pi^2}{GM}$.
\end{itemize}
\item Un satellite est \textbf{géostationnaire} si et seulement si :
\begin{itemize}
\item sa trajectoire est un cercle situé dans le \textbf{plan équatorial} ;
\item il tourne dans le \textbf{même sens} que la rotation propre de la Terre ;
\item sa période de révolution est \textbf{égale à la période de rotation propre} de la Terre : $T = T_0 \approx \num{86164}\ \mathrm{s}$ (jour sidéral, soit environ 23 h 56 min 4 s).
\end{itemize}
\end{enumerate}
\end{corrige}

\begin{corrige}[Exercice 4 — Calculs directs]
\begin{enumerate}
\item Le rayon de l'orbite est $r = R_T + h = \num{6,38}\times 10^{6} + \num{0,400}\times 10^{6} = \num{6,78}\times 10^{6}\ \mathrm{m}$. On calcule d'abord $G\,M_T = \num{6,67}\times 10^{-11} \times \num{5,98}\times 10^{24} = \num{3,99}\times 10^{14}\ \mathrm{m^3.s^{-2}}$.
\[ v = \sqrt{\dfrac{G\,M_T}{r}} = \sqrt{\dfrac{\num{3,99}\times 10^{14}}{\num{6,78}\times 10^{6}}} = \sqrt{\num{5,88}\times 10^{7}} \approx \num{7,67}\times 10^{3}\ \mathrm{m.s^{-1}} \]
soit $v \approx \num{7,7}\ \mathrm{km.s^{-1}}$ (environ \num{27600} km/h !).
\item La période est la durée d'un tour complet :
\[ T = \dfrac{2\pi r}{v} = \dfrac{2\pi \times \num{6,78}\times 10^{6}}{\num{7,67}\times 10^{3}} \approx \num{5,55}\times 10^{3}\ \mathrm{s} \approx \num{92,5}\ \mathrm{min}. \]
La station spatiale fait donc environ 16 tours de Terre par jour : ses occupants assistent à 16 levers de Soleil quotidiens !
\item Analyse dimensionnelle : $[G] = \mathrm{m^3.kg^{-1}.s^{-2}}$, donc
\[ \left[\dfrac{G\,M_T}{r}\right] = \dfrac{\mathrm{m^3.kg^{-1}.s^{-2}} \times \mathrm{kg}}{\mathrm{m}} = \mathrm{m^2.s^{-2}} \quad\Longrightarrow\quad \left[\sqrt{\dfrac{G\,M_T}{r}}\right] = \mathrm{m.s^{-1}}. \]
L'expression est bien homogène à une vitesse.
\end{enumerate}
\end{corrige}

\courssub{Je m'entraîne}

\begin{corrige}[Exercice 5 — Deux satellites d'un même astre]
\begin{enumerate}
\item Troisième loi de Kepler pour le même astre attracteur : $\dfrac{T_A^2}{r_A^3} = \dfrac{T_B^2}{r_B^3}$, donc
\[ \dfrac{T_B^2}{T_A^2} = \dfrac{r_B^3}{r_A^3} = 4^3 = 64 \quad\Longrightarrow\quad \dfrac{T_B}{T_A} = \sqrt{64} = 8. \]
\item $v = \sqrt{\dfrac{GM}{r}}$, donc
\[ \dfrac{v_B}{v_A} = \sqrt{\dfrac{r_A}{r_B}} = \sqrt{\dfrac{1}{4}} = \dfrac{1}{2}. \]
\item \textbf{Conclusion :} lorsque le rayon de l'orbite augmente, la période \emph{augmente} (comme $r^{3/2}$) et la vitesse \emph{diminue} (comme $r^{-1/2}$). Les satellites lointains tournent plus lentement et mettent plus de temps à faire un tour.
\end{enumerate}
\end{corrige}

\begin{corrige}[Exercice 6 — La Lune et la constante de Kepler]
\begin{enumerate}
\item $T_L = \num{27,3} \times \num{86400} = \num{2,36}\times 10^{6}\ \mathrm{s}$ ; \quad $T_G = 23\times 3600 + 56\times 60 = \num{86160}\ \mathrm{s}$.
\item Pour la Lune :
\[ \dfrac{T_L^2}{r_L^3} = \dfrac{(\num{2,36}\times 10^{6})^2}{(\num{3,84}\times 10^{8})^3} = \dfrac{\num{5,56}\times 10^{12}}{\num{5,66}\times 10^{25}} \approx \num{9,83}\times 10^{-14}\ \mathrm{s^2.m^{-3}}. \]
Pour le satellite géostationnaire :
\[ \dfrac{T_G^2}{r_G^3} = \dfrac{(\num{86160})^2}{(\num{4,22}\times 10^{7})^3} = \dfrac{\num{7,42}\times 10^{9}}{\num{7,52}\times 10^{22}} \approx \num{9,88}\times 10^{-14}\ \mathrm{s^2.m^{-3}}. \]
Les deux rapports sont \textbf{égaux} (à moins de 1 \% près, écart dû aux arrondis des données) : la constante de Kepler est la même pour tous les satellites de la Terre, qu'il s'agisse d'un satellite artificiel ou de la Lune.
\item De $\dfrac{T^2}{r^3} = \dfrac{4\pi^2}{G\,M_T}$, on tire :
\[ M_T = \dfrac{4\pi^2}{G \times \dfrac{T^2}{r^3}} = \dfrac{39,48}{\num{6,67}\times 10^{-11} \times \num{9,85}\times 10^{-14}} \approx \num{6,0}\times 10^{24}\ \mathrm{kg}. \]
On retrouve la valeur admise $M_T \approx \num{5,98}\times 10^{24}\ \mathrm{kg}$.
\item La loi de gravitation universelle, combinée à la deuxième loi de Newton, prédit que $\dfrac{T^2}{r^3}$ est une constante \emph{indépendante du satellite}. Le fait que la Lune (corps naturel) et un satellite artificiel (corps minuscule devant la Lune) vérifient la même constante montre que la même loi régit la chute des corps, les satellites et les astres : c'est une confirmation spectaculaire du caractère \emph{universel} de la gravitation.
\end{enumerate}
\end{corrige}

\begin{corrige}[Exercice 7 — Masse du Soleil et masse de Jupiter]
\begin{enumerate}
\item Troisième loi de Kepler appliquée à la Terre autour du Soleil : $\dfrac{T^2}{r^3} = \dfrac{4\pi^2}{G\,M_S}$, d'où
\[ M_S = \dfrac{4\pi^2\,r^3}{G\,T^2}. \]
Conversion : $T = \num{365,25} \times \num{86400} = \num{3,156}\times 10^{7}\ \mathrm{s}$, donc $T^2 = \num{9,96}\times 10^{14}\ \mathrm{s^2}$ et $r^3 = (\num{1,50}\times 10^{11})^3 = \num{3,375}\times 10^{33}\ \mathrm{m^3}$.
\[ M_S = \dfrac{39,48 \times \num{3,375}\times 10^{33}}{\num{6,67}\times 10^{-11} \times \num{9,96}\times 10^{14}} = \dfrac{\num{1,332}\times 10^{35}}{\num{6,64}\times 10^{4}} \approx \num{2,0}\times 10^{30}\ \mathrm{kg}. \]
\item Même formule avec Io, satellite de Jupiter : $T_{Io} = \num{1,77} \times \num{86400} = \num{1,529}\times 10^{5}\ \mathrm{s}$, $T_{Io}^2 = \num{2,34}\times 10^{10}\ \mathrm{s^2}$, $r_{Io}^3 = (\num{4,22}\times 10^{8})^3 = \num{7,52}\times 10^{25}\ \mathrm{m^3}$.
\[ M_J = \dfrac{39,48 \times \num{7,52}\times 10^{25}}{\num{6,67}\times 10^{-11} \times \num{2,34}\times 10^{10}} = \dfrac{\num{2,97}\times 10^{27}}{\num{1,56}} \approx \num{1,9}\times 10^{27}\ \mathrm{kg}. \]
\item $\dfrac{M_J}{M_S} = \dfrac{\num{1,9}\times 10^{27}}{\num{2,0}\times 10^{30}} \approx \num{9,5}\times 10^{-4} \approx \dfrac{1}{1050}$. Jupiter, la plus massive des planètes, ne pèse qu'environ un millième de la masse du Soleil : le Soleil concentre à lui seul plus de 99,8 \% de la masse du système solaire, ce qui justifie de le considérer comme immobile au centre du référentiel héliocentrique.
\end{enumerate}
\end{corrige}

\begin{corrige}[Exercice 8 — Satellite météorologique]
\begin{enumerate}
\item Dans le référentiel géocentrique supposé galiléen, le satellite n'est soumis qu'à la \textbf{force de gravitation terrestre} :
\[ \vec{F} = -\dfrac{G\,M_T\,m}{r^2}\,\vec{u}_n \quad \text{avec } r = R_T + h, \]
où $\vec{u}_n$ est le vecteur unitaire dirigé du satellite $S$ vers le centre $O$ de la Terre.
\begin{popfigure}
\begin{center}
\begin{tikzpicture}[scale=1.0]
\fill[poppink!60] (0,0) circle (0.9);
\node at (0,0) {\small Terre};
\draw[popmuted, dashed] (0,0) circle (2.2);
\fill[popblue] (2.2,0) circle (0.09) node[right=2pt] {$S$};
\draw[-{Stealth}, poporange, very thick] (2.2,0) -- (0.95,0) node[midway, above, poporange] {$\vec{F}$};
\draw[-{Stealth}, popdark] (2.2,0) -- (2.2,0.9) node[above] {$\vec{u}_t$};
\draw[-{Stealth}, popdark] (2.2,0) -- (1.3,0) node[midway, below] {$\vec{u}_n$};
\draw[-{Stealth}, popgreen, thick] (2.2,0) -- (3.1,0.7) node[right] {$\vec{v}$};
\end{tikzpicture}
\end{center}
\legende{Bilan des forces : la seule force de gravitation est centripète.}
\end{popfigure}
\item Appliquons la deuxième loi de Newton dans le repère de Frenet : $\vec{F} = m\,\vec{a}$. En projection sur $\vec{u}_t$ : $m\,\dfrac{dv}{dt} = 0$, car la force est purement normale. Donc $\dfrac{dv}{dt} = 0$ : la vitesse est constante, le mouvement circulaire est \textbf{uniforme}.
\item $r = \num{6,38}\times 10^{6} + \num{0,85}\times 10^{6} = \num{7,23}\times 10^{6}\ \mathrm{m}$.
\[ v = \sqrt{\dfrac{\num{3,99}\times 10^{14}}{\num{7,23}\times 10^{6}}} = \sqrt{\num{5,52}\times 10^{7}} \approx \num{7,43}\times 10^{3}\ \mathrm{m.s^{-1}}. \]
\[ T = \dfrac{2\pi r}{v} = \dfrac{2\pi \times \num{7,23}\times 10^{6}}{\num{7,43}\times 10^{3}} \approx \num{6,11}\times 10^{3}\ \mathrm{s} \approx \num{102}\ \mathrm{min} \approx \num{1,7}\ \mathrm{h}. \]
\item Force de gravitation subie sur orbite :
\[ F = \dfrac{G\,M_T\,m}{r^2} = \dfrac{\num{3,99}\times 10^{14} \times 850}{(\num{7,23}\times 10^{6})^2} = \dfrac{\num{3,39}\times 10^{17}}{\num{5,23}\times 10^{13}} \approx \num{6,5}\times 10^{3}\ \mathrm{N}. \]
Poids à la surface : $P_0 = m\,g_0 = 850 \times \num{9,8} \approx \num{8,3}\times 10^{3}\ \mathrm{N}$. La force en orbite vaut environ 78 \% du poids au sol : à \num{850} km d'altitude, la gravité n'a diminué que d'environ 22 \%. L'« impesanteur » des astronautes n'est donc pas due à l'absence de gravité, mais au fait que le satellite et son contenu sont en chute libre permanente autour de la Terre.
\end{enumerate}
\end{corrige}

\courssub{Je me prépare au Bac}

\begin{corrige}[Exercice 9 — Type Bac : satellite de télécommunications au service du Cameroun]
\textbf{Barème indicatif (sur 10) :} 1 pt (q.1) ; 1 pt (q.2) ; 1,5 pt (q.3) ; 1,5 pt (q.4) ; 3 pts (q.5) ; 1 pt (q.6) ; 1 pt (q.7).

\begin{enumerate}
\item Le référentiel géocentrique a pour origine le centre $O$ de la Terre et ses axes pointent vers trois étoiles lointaines fixes. Il est en translation quasi circulaire autour du Soleil ; pour une étude de durée très inférieure à l'année, son accélération d'entraînement est négligeable : on peut le considérer \textbf{galiléen}.
\item Le satellite n'est soumis qu'à la force de gravitation $\vec{F} = -\dfrac{G\,M_T\,m}{r^2}\,\vec{u}_n$, dirigée de $S$ vers $O$.
\begin{popfigure}
\begin{center}
\begin{tikzpicture}[scale=1.0]
\fill[poppink!60] (0,0) circle (0.9);
\node at (0,0) {\small Terre};
\draw[popmuted, dashed] (0,0) circle (2.3);
\fill[popblue] (2.3,0) circle (0.09) node[right=2pt] {$S$};
\draw[-{Stealth}, poporange, very thick] (2.3,0) -- (0.95,0) node[midway, above, poporange] {$\vec{F}$};
\draw[-{Stealth}, popdark] (2.3,0) -- (2.3,0.95) node[above] {$\vec{u}_t$};
\draw[-{Stealth}, popdark] (2.3,0) -- (1.35,0) node[midway, below] {$\vec{u}_n$};
\end{tikzpicture}
\end{center}
\legende{Force de gravitation et repère de Frenet $(S,\vec{u}_t,\vec{u}_n)$.}
\end{popfigure}
\item Deuxième loi de Newton : $\vec{F} = m\,\vec{a}$. Dans le repère de Frenet, $\vec{a} = \dfrac{dv}{dt}\,\vec{u}_t + \dfrac{v^2}{r}\,\vec{u}_n$. La force n'a aucune composante tangentielle : la projection sur $\vec{u}_t$ donne $m\,\dfrac{dv}{dt} = 0$, donc $v = \text{constante}$ : le mouvement circulaire est \textbf{uniforme}.
\item Projection sur $\vec{u}_n$ : $m\,\dfrac{v^2}{r} = \dfrac{G\,M_T\,m}{r^2}$. En simplifiant par $m$ et en multipliant par $r$ : $v^2 = \dfrac{G\,M_T}{r}$, d'où
\[ v = \sqrt{\dfrac{G\,M_T}{r}}. \]
\item 
\begin{enumerate}
\item Conditions géostationnaires : orbite circulaire dans le \textbf{plan équatorial} ; rotation dans le \textbf{sens de rotation propre de la Terre} ; période $T = T_0 = \num{86164}\ \mathrm{s}$.
\item Le mouvement étant uniforme, $v = \dfrac{2\pi r}{T}$. En égalant avec $v^2 = \dfrac{G\,M_T}{r}$ :
\[ \dfrac{4\pi^2 r^2}{T^2} = \dfrac{G\,M_T}{r} \quad\Longrightarrow\quad \dfrac{T^2}{r^3} = \dfrac{4\pi^2}{G\,M_T}. \]
\item On en tire $r^3 = \dfrac{G\,M_T\,T_0^2}{4\pi^2}$, soit
\[ r = \left(\dfrac{G\,M_T\,T_0^2}{4\pi^2}\right)^{1/3}. \]
Toutes les grandeurs du membre de droite sont fixées ($G$, $M_T$, $T_0$) : le rayon est \textbf{unique}, donc l'altitude $h = r - R_T$ est unique.
\[ r^3 = \dfrac{\num{3,99}\times 10^{14} \times (\num{86164})^2}{39,48} = \dfrac{\num{3,99}\times 10^{14} \times \num{7,42}\times 10^{9}}{39,48} = \num{7,50}\times 10^{22}\ \mathrm{m^3} \]
\[ r \approx \num{4,22}\times 10^{7}\ \mathrm{m} \quad\Longrightarrow\quad h = r - R_T \approx \num{4,22}\times 10^{7} - \num{6,38}\times 10^{6} \approx \num{3,58}\times 10^{7}\ \mathrm{m} \approx \num{35800}\ \mathrm{km}. \]
\end{enumerate}
\item $v = \sqrt{\dfrac{G\,M_T}{r}} = \sqrt{\dfrac{\num{3,99}\times 10^{14}}{\num{4,22}\times 10^{7}}} = \sqrt{\num{9,45}\times 10^{6}} \approx \num{3,08}\times 10^{3}\ \mathrm{m.s^{-1}}$.
\item Un point de l'équateur (Kribi est proche de l'équateur) décrit un cercle de rayon $R_T$ en $T_0$ :
\[ v_D = \dfrac{2\pi R_T}{T_0} = \dfrac{2\pi \times \num{6,38}\times 10^{6}}{\num{86164}} \approx \num{4,65}\times 10^{2}\ \mathrm{m.s^{-1}}. \]
On a $v \approx \num{6,6}\,v_D$ : le satellite va bien plus vite en valeur absolue. Mais sa \emph{vitesse angulaire} $\omega = \dfrac{2\pi}{T_0}$ est exactement celle de la Terre : satellite et habitant de Kribi tournent « ensemble », d'où l'immobilité apparente du satellite vu du sol.
\end{enumerate}
\textbf{Points de vigilance :} citer explicitement le référentiel et son caractère galiléen avant d'appliquer la deuxième loi de Newton ; écrire la projection sur $\vec{u}_t$ \emph{avant} celle sur $\vec{u}_n$ ; ne pas oublier de convertir $T_0$ en secondes ; distinguer soigneusement $r$ (rayon, mesuré depuis le centre de la Terre) et $h$ (altitude, mesurée depuis la surface) : $h = r - R_T$.
\end{corrige}

\begin{corrige}[Exercice 10 — Type Bac : peser le Soleil et vérifier Kepler]
\textbf{Barème indicatif (sur 10) :} 1 pt (q.1) ; 1,5 pt (q.2) ; 3 pts (q.3) ; 1,5 pt (q.4) ; 1 pt (q.5) ; 2 pts (q.6).

\begin{enumerate}
\item On étudie la Terre dans le \textbf{référentiel héliocentrique} : origine au centre du Soleil, axes dirigés vers trois étoiles lointaines fixes ; il est galiléen avec une excellente approximation.
\item \emph{Loi des orbites} : chaque planète décrit une ellipse dont le Soleil est un foyer. \emph{Loi des aires} : le rayon Soleil--planète balaie des aires égales pendant des durées égales. \emph{Loi des périodes} : $\dfrac{T^2}{a^3}$ est identique pour toutes les planètes du système solaire.
\item 
\begin{enumerate}
\item $\vec{F}_{S/T} = -\dfrac{G\,M_S\,m}{r_T^2}\,\vec{u}_n$, où $\vec{u}_n$ pointe de la Terre vers le Soleil.
\item Deuxième loi de Newton dans le repère de Frenet, projection sur $\vec{u}_t$ : $m\,\dfrac{dv}{dt} = 0$ car $\vec{F}$ est radiale. Donc $v$ est constante : le mouvement circulaire est uniforme.
\item Projection sur $\vec{u}_n$ : $\dfrac{m\,v^2}{r_T} = \dfrac{G\,M_S\,m}{r_T^2}$, d'où $v = \sqrt{\dfrac{G\,M_S}{r_T}}$.
\end{enumerate}
\item Avec $v = \dfrac{2\pi r_T}{T_T}$ : $\dfrac{4\pi^2 r_T^2}{T_T^2} = \dfrac{G\,M_S}{r_T}$, donc
\[ M_S = \dfrac{4\pi^2\,r_T^3}{G\,T_T^2}. \]
Numériquement : $T_T = \num{365,25}\times\num{86400} = \num{3,156}\times 10^{7}\ \mathrm{s}$ ; $T_T^2 = \num{9,96}\times 10^{14}\ \mathrm{s^2}$ ; $r_T^3 = \num{3,375}\times 10^{33}\ \mathrm{m^3}$.
\[ M_S = \dfrac{39,48 \times \num{3,375}\times 10^{33}}{\num{6,67}\times 10^{-11} \times \num{9,96}\times 10^{14}} \approx \num{2,0}\times 10^{30}\ \mathrm{kg}. \]
Le Soleil a bien été « pesé » depuis la salle de classe !
\item Calculons la constante de Kepler avec notre masse :
\[ K = \dfrac{4\pi^2}{G\,M_S} = \dfrac{39,48}{\num{6,67}\times 10^{-11} \times \num{2,0}\times 10^{30}} \approx \num{2,96}\times 10^{-19}\ \mathrm{s^2.m^{-3}}. \]
C'est en excellent accord avec la valeur de référence $K = \num{2,97}\times 10^{-19}\ \mathrm{s^2.m^{-3}}$ (écart inférieur à 1 \%) : le résultat est cohérent.
\item Troisième loi de Kepler pour la Terre et Mars (même astre attracteur) :
\[ \dfrac{T_M^2}{r_M^3} = \dfrac{T_T^2}{r_T^3} \quad\Longrightarrow\quad r_M = r_T\left(\dfrac{T_M}{T_T}\right)^{2/3}. \]
\[ r_M = \num{1,50}\times 10^{11} \times \left(\dfrac{687}{\num{365,25}}\right)^{2/3} = \num{1,50}\times 10^{11} \times (\num{1,88})^{2/3} \approx \num{1,50}\times 10^{11} \times \num{1,52} \approx \num{2,29}\times 10^{11}\ \mathrm{m}. \]
Mars orbite donc environ 1,5 fois plus loin du Soleil que la Terre (1,52 unité astronomique).
\end{enumerate}
\textbf{Points de vigilance :} la démonstration de la troisième loi de Kepler dans le cas circulaire est exigible ; bien justifier l'uniformité du mouvement \emph{avant} d'écrire $v = \dfrac{2\pi r}{T}$ ; dans la question 6, exploiter le rapport des constantes de Kepler plutôt que de recalculer $M_S$ (méthode plus rapide et plus sûre).
\end{corrige}

\begin{corrige}[Exercice 11 — Type Bac : mise sur orbite géostationnaire — étude énergétique]
\textbf{Barème indicatif (sur 10) :} 1,5 pt (q.1) ; 1,5 pt (q.2) ; 2 pts (q.3) ; 1,5 pt (q.4) ; 2 pts (q.5) ; 1,5 pt (q.6).

\begin{enumerate}
\item Sur une orbite circulaire, $v^2 = \dfrac{G\,M_T}{r}$ (question de cours). Donc
\[ E_c = \dfrac{1}{2}m\,v^2 = \dfrac{1}{2}m\,\dfrac{G\,M_T}{r} = \dfrac{G\,M_T\,m}{2r}. \]
\item $E_m = E_c + E_p = \dfrac{G\,M_T\,m}{2r} - \dfrac{G\,M_T\,m}{r} = -\dfrac{G\,M_T\,m}{2r}$. L'énergie mécanique est \textbf{négative} : le satellite est dans un état lié, il ne peut pas s'échapper seul de l'attraction terrestre (l'état de référence $E_m = 0$ correspond à un satellite au repos à l'infini).
\item Orbite basse : $r_1 = \num{6,38}\times 10^{6} + \num{0,30}\times 10^{6} = \num{6,68}\times 10^{6}\ \mathrm{m}$. Avec $G\,M_T\,m = \num{3,99}\times 10^{14} \times \num{1,2}\times 10^{3} = \num{4,79}\times 10^{17}\ \mathrm{J.m}$ :
\[ E_{m1} = -\dfrac{\num{4,79}\times 10^{17}}{2 \times \num{6,68}\times 10^{6}} \approx -\num{3,58}\times 10^{10}\ \mathrm{J}. \]
\[ E_{m2} = -\dfrac{\num{4,79}\times 10^{17}}{2 \times \num{4,22}\times 10^{7}} \approx -\num{5,67}\times 10^{9}\ \mathrm{J}. \]
\item L'énergie minimale à fournir est l'augmentation d'énergie mécanique :
\[ W = E_{m2} - E_{m1} = -\num{5,67}\times 10^{9} - (-\num{3,58}\times 10^{10}) \approx \num{3,02}\times 10^{10}\ \mathrm{J}. \]
En kilowattheures : $W = \dfrac{\num{3,02}\times 10^{10}}{\num{3,6}\times 10^{6}} \approx \num{8,4}\times 10^{3}\ \mathrm{kWh}$, soit la consommation électrique annuelle d'environ deux ménages camerounais urbains.
\item $v_1 = \sqrt{\dfrac{\num{3,99}\times 10^{14}}{\num{6,68}\times 10^{6}}} = \sqrt{\num{5,97}\times 10^{7}} \approx \num{7,73}\times 10^{3}\ \mathrm{m.s^{-1}}$ ;
$v_2 = \sqrt{\dfrac{\num{3,99}\times 10^{14}}{\num{4,22}\times 10^{7}}} \approx \num{3,08}\times 10^{3}\ \mathrm{m.s^{-1}}$.
\textbf{Non} : le satellite va \emph{moins vite} sur l'orbite haute ($v_2 < v_1$). Pourtant il a fallu fournir de l'énergie : celle-ci a servi essentiellement à augmenter l'énergie potentielle ($E_p$ passe de $-\num{7,2}\times 10^{10}$ J à $-\num{1,1}\times 10^{10}$ J), c'est-à-dire à « éloigner » le satellite de la Terre, pendant que l'énergie cinétique \emph{diminuait}.
\item À l'infini, le corps libéré arrive avec une vitesse nulle : $E_m(\infty) = 0$. Conservation de l'énergie mécanique entre la surface et l'infini :
\[ \dfrac{1}{2}m\,v_l^2 - \dfrac{G\,M_T\,m}{R_T} = 0 \quad\Longrightarrow\quad v_l = \sqrt{\dfrac{2\,G\,M_T}{R_T}}. \]
\[ v_l = \sqrt{\dfrac{2 \times \num{3,99}\times 10^{14}}{\num{6,38}\times 10^{6}}} = \sqrt{\num{1,25}\times 10^{8}} \approx \num{1,12}\times 10^{4}\ \mathrm{m.s^{-1}} \approx \num{11,2}\ \mathrm{km.s^{-1}}. \]
Comparaison : $v_l = \sqrt{2}\,v_{\text{surface}} \approx \num{1,41}\,v_1$ (car $v_1 \approx \sqrt{\dfrac{G\,M_T}{R_T}}$ à basse altitude). Il faut donc gagner encore environ 41 \% de vitesse par rapport à l'orbite basse pour échapper définitivement à l'attraction terrestre.
\end{enumerate}
\textbf{Points de vigilance :} ne pas oublier le signe négatif de $E_p$ et de $E_m$ ; l'énergie à fournir est $W = \Delta E_m = E_{m2} - E_{m1} > 0$ (attention au piège des signes : $E_{m2}$ est moins négative que $E_{m1}$) ; pour la vitesse de libération, écrire explicitement la conservation de l'énergie mécanique et la condition « à peine libéré » : vitesse nulle à l'infini.
\end{corrige}

\newpage

\coursec{Fiche bilan}

\begin{popfigure}
\centering
\begin{tikzpicture}[
    mindmap,
    grow cyclic,
    every node/.style={concept, execute at begin node=\hskip0pt},
    concept color=popblue,
    level 1/.append style={level distance=4.5cm, sibling angle=360/6},
    level 2/.append style={level distance=3cm, sibling angle=360/6},
    text=white,
    font=\small\bfseries
]
\node[concept, scale=1.2] {Mouvement Planètes \& Satellites (P07)}
    child[concept color=poporange] { node[concept] {1. Référentiels \& Kepler}
        child { node[concept, scale=0.7] {Géocentrique / Héliocentrique} }
        child { node[concept, scale=0.7] {3 Lois de Kepler} }
    }
    child[concept color=popgreen] { node[concept] {2. Satellite Circulaire}
        child { node[concept, scale=0.7] {Vitesse $v = \sqrt{GM/R}$} }
        child { node[concept, scale=0.7] {Période $T = 2\pi\sqrt{R^3/GM}$} }
        child { node[concept, scale=0.7] {3ème Loi : $T^2 \propto R^3$} }
    }
    child[concept color=poppurple] { node[concept] {3. Géostationnaire}
        child { node[concept, scale=0.7] {Période = 23 h 56 min 4 s} }
        child { node[concept, scale=0.7] {Plan équatorial} }
        child { node[concept, scale=0.7] {Altitude $\approx \SI{35 786}{km}$} }
    }
    child[concept color=poppink] { node[concept] {4. Énergie (Complément)}
        child { node[concept, scale=0.7] {$E_m = -\dfrac{GMm}{2R}$} }
        child { node[concept, scale=0.7] {Vitesse libération $v_{\text{lib}} = \sqrt{2GM/R}$} }
    }
    child[concept color=popdark] { node[concept] {5. Méthodes \& Applications}
        child { node[concept, scale=0.7] {Masse astre $M = \frac{4\pi^2 R^3}{GT^2}$} }
        child { node[concept, scale=0.7] {Bilan forces : $\vec{F}_g = m\vec{a}$} }
        child { node[concept, scale=0.7] {Analyse dimensionnelle} }
    };
\end{tikzpicture}
\legende{Carte mentale de la leçon P07 : Mouvement des planètes et des satellites}
\end{popfigure}

\begin{bilan}
\courssub{Définitions et référentiels}
\begin{itemize}
    \item \textbf{Référentiel géocentrique} : origine au centre de la Terre, axes fixes par rapport aux étoiles (quasi-galiléen pour l'étude des satellites terrestres).
    \item \textbf{Référentiel héliocentrique} : origine au centre du Soleil, axes fixes par rapport aux étoiles (galiléen pour l'étude des planètes).
    \item \textbf{Satellite} : corps en orbite autour d'un astre principal sous l'effet exclusif de la gravitation (poids = force centripète).
\end{itemize}

\courssub{Lois de Kepler (énoncés à connaître)}
\begin{enumerate}
    \item \textbf{Loi des orbites} : Les planètes décrivent des ellipses dont le Soleil occupe un foyer. (Pour un satellite : orbite circulaire ou elliptique, astre au foyer).
    \item \textbf{Loi des aires} : Le rayon vecteur (astre $\to$ planète) balaye des aires égales en temps égaux $\implies$ moment cinétique constant (mouvement plan).
    \item \textbf{Loi des périodes (3ème loi)} : Le carré de la période de révolution $T$ est proportionnel au cube du demi-grand axe $a$ de l'orbite : $T^2 = \dfrac{4\pi^2}{GM}a^3$. Pour un cercle, $a = R$.
\end{enumerate}

\courssub{Formules clés (mouvement circulaire uniforme, rayon $R$, masse astre $M$, masse satellite $m$)}
\begin{center}
\begin{tabularx}{\linewidth}{|l|X|X|}\hline
\rowcolor{popblueL}
\textbf{Grandeur} & \textbf{Expression} & \textbf{Unité SI} \\ \hline
Vitesse orbitale & $v = \sqrt{\dfrac{GM}{R}}$ & $\SI{}{m.s^{-1}}$ \\ \hline
Période de révolution & $T = 2\pi\sqrt{\dfrac{R^3}{GM}} = \dfrac{2\pi R}{v}$ & $\SI{}{s}$ \\ \hline
Accélération gravitationnelle & $g(R) = \dfrac{GM}{R^2}$ & $\SI{}{m.s^{-2}}$ \\ \hline
3ème Loi de Kepler & $\dfrac{T^2}{R^3} = \dfrac{4\pi^2}{GM} = \text{cste}$ & $\SI{}{s^2.m^{-3}}$ \\ \hline
Masse de l'astre (si satellite connu) & $M = \dfrac{4\pi^2 R^3}{GT^2}$ & $\SI{}{kg}$ \\ \hline
Énergie potentielle & $E_p = -\dfrac{GMm}{R}$ & $\SI{}{J}$ \\ \hline
Énergie cinétique & $E_c = \dfrac{1}{2}mv^2 = \dfrac{GMm}{2R}$ & $\SI{}{J}$ \\ \hline
Énergie mécanique & $E_m = E_c + E_p = -\dfrac{GMm}{2R}$ & $\SI{}{J}$ \\ \hline
Vitesse de libération & $v_{\text{lib}} = \sqrt{\dfrac{2GM}{R}} = \sqrt{2}\,v$ & $\SI{}{m.s^{-1}}$ \\ \hline
\end{tabularx}
\end{center}

\courssub{Satellite géostationnaire (conditions nécessaires et suffisantes)}
\begin{enumerate}
    \item Mouvement circulaire uniforme.
    \item Orbite dans le plan équatorial (inclinaison $0^\circ$).
    \item Sens de rotation Est $\to$ Ouest (identique à la Terre).
    \item Période sidérale $T = \SI{23}{h}\ \SI{56}{min}\ \SI{4}{s} = \SI{86164}{s}$.
    \item Altitude $h \approx \SI{35 786}{km}$ au-dessus du niveau de la mer ($R \approx \SI{42 164}{km}$).
\end{enumerate}

\courssub{Méthode express : Déterminer la masse d'un astre}
\begin{enumerate}
    \item Identifier un satellite (lune, satellite artificiel) en orbite \emph{quasi-circulaire}.
    \item Mesurer ou relever : rayon orbital $R$ (distance centre-à-centre) et période $T$.
    \item Appliquer $M = \dfrac{4\pi^2 R^3}{GT^2}$.
    \item Vérifier l'homogénéité : $[M] = \dfrac{[L]^3}{[M^{-1}L^3T^{-2}][T]^2} = [M]$. \checkmark
\end{enumerate}

\courssub{Méthode express : Altitude satellite géostationnaire}
\begin{enumerate}
    \item Écrire $T = 2\pi\sqrt{R^3/GM_{\text{Terre}}}$.
    \item Isoler $R = \sqrt[3]{\dfrac{GM_{\text{Terre}}T^2}{4\pi^2}}$.
    \item Calculer $h = R - R_{\text{Terre}}$ ($R_{\text{Terre}} \approx \SI{6371}{km}$).
\end{enumerate}
\end{bilan}

\begin{aretenir}[Checklist avant le Bac \textemdash{} P07]
\begin{itemize}
    \item $\square$ Différencier référentiel géocentrique (satellites Terre) et héliocentrique (planètes).
    \item $\square$ Énoncer les 3 lois de Kepler sans erreur (ellipses, aires, $T^2 \propto a^3$).
    \item $\square$ Démontrer $v = \sqrt{GM/R}$ et $T = 2\pi\sqrt{R^3/GM}$ à partir du PFD ($\vec{F}_g = m\vec{a}$).
    \item $\square$ Savoir exploiter $T^2/R^3 = \text{cste}$ pour comparer deux satellites d'un même astre.
    \item $\square$ Calculer la masse $M$ d'un astre connaissant $R$ et $T$ d'un satellite ($M = 4\pi^2 R^3/GT^2$).
    \item $\square$ Lister les 4 conditions du satellite géostationnaire et donner son altitude ($\approx \SI{36 000}{km}$).
    \item $\square$ Écrire les expressions de $E_p$, $E_c$, $E_m$ pour un satellite circulaire ($E_m < 0$, $|E_m| = E_c = -E_p/2$).
    \item $\square$ Définir la vitesse de libération $v_{\text{lib}} = \sqrt{2GM/R}$ et la comparer à $v_{\text{circ}}$ ($v_{\text{lib}} = \sqrt{2}\,v_{\text{circ}}$).
    \item $\square$ Vérifier systématiquement l'homogénéité des formules (analyse dimensionnelle).
    \item $\square$ Attention : $R$ = distance centre-à-centre (pas l'altitude $h$) ; $T$ en secondes (pas en heures).
    \item $\square$ Contexte Cameroun : Satellites de télécoms (MTN/Orange/CAMTEL) géostationnaires ; barrage de Lom Pangar / Nachtigal (pesanteur locale $g$).
\end{itemize}
\end{aretenir}

\findecours

\end{document}
